% Special Discrete Probability Distributions — Further Mathematics Upper Sixth — Cours signé OnBuch+
% Official MINESEC syllabus. Compiler avec : tectonic FMA18-special-discrete-probability-distributions.tex
\newcommand{\DOCMATIERE}{Further Mathematics}
\newcommand{\DOCNIVEAU}{Upper Sixth}
\newcommand{\DOCMODULE}{Upper Sixth — Further mathematics}
\newcommand{\DOCLECON}{18}
\newcommand{\DOCTITRE}{Special Discrete Probability Distributions}
% OnBuch+ — préambule « cours signés » ENRICHI (compilé avec tectonic / XeLaTeX)
% Système visuel « Pop » d'OnBuch+ : fond crème (jamais blanc), bordures encre
% épaisses, ombres « dures » décalées, titres Archivo Black en capitales.
% Version MATHÉMATIQUES : graphes (pgfplots/tikz), boîtes pédagogiques
% (définition, propriété, méthode, attention, le savais-tu, exercice résolu,
% exercice, TP, bilan), page de garde et signature OnBuch+ ancrée en bas.
%
% Macros attendues dans main.tex AVANT \input{preamble} :
%   \DOCMATIERE  \DOCNIVEAU  \DOCMODULE  \DOCLECON (numéro)  \DOCTITRE
\documentclass[11pt]{article}

\usepackage{fontspec}
\usepackage[english]{babel}
\usepackage[a4paper,margin=2cm,top=3.4cm,bottom=2.8cm,headheight=1.4cm,footskip=1.3cm]{geometry}
\usepackage{amsmath,amssymb,amsfonts}
\usepackage{mathtools} % \xmapsto, \coloneqq... souvent utilisés par les modèles
\usepackage{cancel} % simplifications barrées (bilans de forces)
% \abs{z} (module, valeur absolue) et \norm{u} (norme), utilisés par les modèles
\providecommand{\abs}{}\let\abs\relax\DeclarePairedDelimiter{\abs}{\lvert}{\rvert}
\providecommand{\norm}{}\let\norm\relax\DeclarePairedDelimiter{\norm}{\lVert}{\rVert}
\usepackage[table]{xcolor} % [table] : fournit aussi \rowcolors (alternance de lignes)
\usepackage{pagecolor}
\usepackage{tikz}
\usetikzlibrary{arrows.meta,positioning,calc,shapes.geometric,shapes.misc,shapes.arrows,decorations.pathmorphing,decorations.pathreplacing,decorations.markings,patterns,shadows,mindmap,trees,fit,backgrounds,matrix,intersections,angles,quotes}
\usepackage{pgfplots}
\usepackage[version=4]{mhchem} % équations nucléaires \ce{^{235}_{92}U}
\usepackage[european,siunitx]{circuitikz} % schémas électriques
\pgfplotsset{compat=1.18}
\usepgfplotslibrary{fillbetween,groupplots} % aires entre courbes (calcul intégral)
\usepackage{enumitem}
\usepackage{fancyhdr}
% [most] charge toutes les bibliothèques tcolorbox (listings, minted, raster,
% documentation...) et gonfle inutilement la mémoire TeX sur un document long
% (~300 boîtes) : on ne charge que ce qui est réellement utilisé.
\usepackage[skins,breakable]{tcolorbox}
\usepackage{graphicx}
\usepackage{colortbl}
\usepackage{booktabs}
\usepackage{tabularx}
\usepackage{diagbox} % case d'en-tête diagonale des tableaux à double entrée
% Colonnes extensibles centrée / gauche / droite (C, L, R), souvent utilisées par les modèles
\newcolumntype{C}{>{\centering\arraybackslash}X}
\newcolumntype{L}{>{\raggedright\arraybackslash}X}
\newcolumntype{R}{>{\raggedleft\arraybackslash}X}
\usepackage{array}
\usepackage{multicol}
\usepackage{siunitx}
% parse-numbers=false : accepte aussi \SI{2,5 \times 10^{-3}}{...}
\sisetup{locale=UK,output-decimal-marker={.},group-minimum-digits=4,parse-numbers=false}
\DeclareSIUnit\ohm{\ensuremath{\symup{\Omega}}} % U+2126 absent de la police
\DeclareSIUnit\division{div} % réglages d'oscilloscope (ms/div, V/div)
\DeclareSIUnit\tr{tr} % tours (vitesse de rotation, ex. \SI{1500}{\tr\per\minute})
\usepackage{qrcode}
% \degree : commande fréquemment utilisée par les modèles pour un angle ou une
% température en dehors de \SI{}{} (non fournie par les paquets chargés).
\providecommand{\degree}{^{\circ}}
% \qed (fin de démonstration), utilisé par les modèles sans amsthm
\providecommand{\qed}{\hfill$\square$}
% \barre : double barre des valeurs interdites dans un tableau de variations (tkz-tab)
\providecommand{\barre}{\ensuremath{\|}}
% environnement proof (amsthm non chargé) utilisé par les modèles
\ifdefined\proof\else\newenvironment{proof}[1][Proof]{\par\noindent\textit{#1.}\ }{\qed\par}\fi
\usepackage{lastpage}
\usepackage{hyperref}
\hypersetup{hidelinks}

% ============================================================
% Palette « Pop » OnBuch+ — mêmes hex que la classe P (plus_ui.dart)
% ============================================================
\definecolor{poporange}{HTML}{F59321}
\definecolor{popdark}{HTML}{E07A0C}
\definecolor{popcream}{HTML}{FFF6E8}
\definecolor{popink}{HTML}{1C1714}
\definecolor{poppurple}{HTML}{7A5AE0}
\definecolor{popgreen}{HTML}{2E9E6A}
\definecolor{poppink}{HTML}{E0457B}
\definecolor{popblue}{HTML}{2F7DE1}
\definecolor{popgold}{HTML}{C98A12}
\definecolor{popmuted}{HTML}{8A7F76}
\definecolor{popline}{HTML}{EADFD2}
\definecolor{popgray}{HTML}{8A7F76}
\colorlet{popgrey}{popgray}
% Alias tolérants (le modèle nomme parfois les couleurs différemment)
\colorlet{popwhite}{white}
\colorlet{popblack}{popink}
\colorlet{popred}{poppink}
\colorlet{popyellow}{popgold}
% Teintes claires pour fonds de tableaux / zones
\colorlet{popblueL}{popblue!10!white}
\colorlet{poporangeL}{poporange!14!white}
\colorlet{popgreenL}{popgreen!12!white}
\colorlet{poppurpleL}{poppurple!12!white}
\colorlet{poppinkL}{poppink!12!white}
\colorlet{popgoldL}{popgold!14!white}

\pagecolor{popcream}

% ============================================================
% Typographie
% ============================================================
\defaultfontfeatures{Ligatures=TeX}
\setmainfont{PlusJakartaSans-Medium.ttf}[Path=../../../fonts/,BoldFont=PlusJakartaSans-Bold.ttf]
% Maths en sans-serif (Fira Math) avec chiffres et lettres droites de Plus
% Jakarta, pour que les formules se fondent dans le texte.
\usepackage[math-style=ISO,bold-style=ISO]{unicode-math}
\setmathfont{FiraMath-Regular.otf}[Path=../../../fonts/]
\setmathfont{PlusJakartaSans-Medium.ttf}[Path=../../../fonts/,range={up/num,up/{Latin,latin},\mathup}]
% (icomma retiré : décimales avec point en anglais)
% Caractères mathématiques Unicode absents des polices de texte (Plus Jakarta,
% Archivo Black) : rendus par leur équivalent mathématique, en texte comme en
% formule — sinon ils s'affichent en carré vide (ex. « Arithmétique dans ℤ »).
\usepackage{newunicodechar}
\newunicodechar{ℝ}{\ensuremath{\mathbb{R}}}
\newunicodechar{ℤ}{\ensuremath{\mathbb{Z}}}
\newunicodechar{ℂ}{\ensuremath{\mathbb{C}}}
\newunicodechar{ℕ}{\ensuremath{\mathbb{N}}}
\newunicodechar{ℚ}{\ensuremath{\mathbb{Q}}}
\newunicodechar{𝔻}{\ensuremath{\mathbb{D}}}
\newunicodechar{α}{\ensuremath{\alpha}}
\newunicodechar{β}{\ensuremath{\beta}}
\newunicodechar{γ}{\ensuremath{\gamma}}
\newunicodechar{δ}{\ensuremath{\delta}}
\newunicodechar{ε}{\ensuremath{\varepsilon}}
\newunicodechar{θ}{\ensuremath{\theta}}
\newunicodechar{λ}{\ensuremath{\lambda}}
\newunicodechar{μ}{\ensuremath{\mu}}
\newunicodechar{σ}{\ensuremath{\sigma}}
\newunicodechar{τ}{\ensuremath{\tau}}
\newunicodechar{φ}{\ensuremath{\varphi}}
\newunicodechar{ψ}{\ensuremath{\psi}}
\newunicodechar{ω}{\ensuremath{\omega}}
\newunicodechar{Σ}{\ensuremath{\Sigma}}
% majuscules produites par \MakeUppercase dans les titres (α-aminés -> Α)
\newunicodechar{Α}{\ensuremath{\alpha}}
\newunicodechar{Β}{\ensuremath{\beta}}
\newunicodechar{Γ}{\ensuremath{\Gamma}}
\newunicodechar{Δ}{\ensuremath{\Delta}}
\newunicodechar{Ε}{\ensuremath{\varepsilon}}
\newunicodechar{Θ}{\ensuremath{\Theta}}
\newunicodechar{Λ}{\ensuremath{\Lambda}}
\newunicodechar{Μ}{\ensuremath{\mu}}
\newunicodechar{Τ}{\ensuremath{\tau}}
\newunicodechar{Φ}{\ensuremath{\Phi}}
\newunicodechar{Ψ}{\ensuremath{\Psi}}
\newunicodechar{Ω}{\ensuremath{\Omega}}
\newunicodechar{↦}{\ensuremath{\mapsto}}
\newunicodechar{∈}{\ensuremath{\in}}
\newunicodechar{∉}{\ensuremath{\notin}}
\newunicodechar{∧}{\ensuremath{\wedge}}
\newunicodechar{⇔}{\ensuremath{\Leftrightarrow}}
\newunicodechar{⟺}{\ensuremath{\Longleftrightarrow}}
\newunicodechar{⇒}{\ensuremath{\Rightarrow}}
\newunicodechar{⟹}{\ensuremath{\Longrightarrow}}
\newunicodechar{∘}{\ensuremath{\circ}}
\newunicodechar{∪}{\ensuremath{\cup}}
\newunicodechar{∩}{\ensuremath{\cap}}
\newunicodechar{≡}{\ensuremath{\equiv}}
\newunicodechar{⊂}{\ensuremath{\subset}}
\newunicodechar{⊥}{\ensuremath{\perp}}
\newunicodechar{∃}{\ensuremath{\exists}}
\newunicodechar{∀}{\ensuremath{\forall}}
\newunicodechar{∛}{\ensuremath{\sqrt[3]{\,}}}
\newunicodechar{∜}{\ensuremath{\sqrt[4]{\,}}}
\newunicodechar{⃗}{\ensuremath{{}^{\to}}}
\newunicodechar{ⁿ}{\ifmmode^{n}\else\textsuperscript{\ensuremath{n}}\fi}
\newunicodechar{ˣ}{\ifmmode^{x}\else\textsuperscript{\ensuremath{x}}\fi}
\newunicodechar{ᵇ}{\ifmmode^{b}\else\textsuperscript{\ensuremath{b}}\fi}
\newunicodechar{ʳ}{\ifmmode^{r}\else\textsuperscript{\ensuremath{r}}\fi}
\newunicodechar{ᵘ}{\ifmmode^{u}\else\textsuperscript{\ensuremath{u}}\fi}
\newunicodechar{ʸ}{\ifmmode^{y}\else\textsuperscript{\ensuremath{y}}\fi}
\newunicodechar{ᵃ}{\ifmmode^{a}\else\textsuperscript{\ensuremath{a}}\fi}
\newunicodechar{ᵉ}{\ifmmode^{e}\else\textsuperscript{\ensuremath{e}}\fi}
\newunicodechar{ᵏ}{\ifmmode^{k}\else\textsuperscript{\ensuremath{k}}\fi}
\newunicodechar{ᵖ}{\ifmmode^{p}\else\textsuperscript{\ensuremath{p}}\fi}
\newunicodechar{ᴺ}{\ifmmode^{N}\else\textsuperscript{\ensuremath{N}}\fi}
\newunicodechar{ᶜ}{\ifmmode^{c}\else\textsuperscript{\ensuremath{c}}\fi}
\newunicodechar{ʰ}{\ifmmode^{h}\else\textsuperscript{\ensuremath{h}}\fi}
\newunicodechar{ᵠ}{\ifmmode^{q}\else\textsuperscript{\ensuremath{q}}\fi}
\newunicodechar{ₙ}{\ifmmode_{n}\else\textsubscript{\ensuremath{n}}\fi}
\newunicodechar{ₐ}{\ifmmode_{a}\else\textsubscript{\ensuremath{a}}\fi}
\newunicodechar{ᵢ}{\ifmmode_{i}\else\textsubscript{\ensuremath{i}}\fi}
\newunicodechar{ᵣ}{\ifmmode_{r}\else\textsubscript{\ensuremath{r}}\fi}
\newunicodechar{ₖ}{\ifmmode_{k}\else\textsubscript{\ensuremath{k}}\fi}
\newunicodechar{ᵦ}{\ifmmode_{\beta}\else\textsubscript{\ensuremath{\beta}}\fi}
\newunicodechar{ₓ}{\ifmmode_{x}\else\textsubscript{\ensuremath{x}}\fi}
\newfontfamily\popdisplay{ArchivoBlack-Regular.ttf}[Path=../../../fonts/]
\newfontfamily\poplogo{SpaceGrotesk-Bold.ttf}[Path=../../../fonts/]
\newfontfamily\popsemi{PlusJakartaSans-SemiBold.ttf}[Path=../../../fonts/]
\newfontfamily\popxbold{PlusJakartaSans-ExtraBold.ttf}[Path=../../../fonts/]
\newfontfamily\popxboldls{PlusJakartaSans-ExtraBold.ttf}[Path=../../../fonts/,LetterSpace=3.0]

\setlist[enumerate]{leftmargin=1.7em,itemsep=5pt,topsep=4pt}
\setlist[itemize]{leftmargin=1.7em,itemsep=4pt,topsep=4pt,label={\color{poporange}\textbullet}}
\setlength{\parindent}{0pt}
\setlength{\parskip}{5pt}
\renewcommand{\arraystretch}{1.25}

% pgfplots : style Pop par défaut
\pgfplotsset{
  every axis/.append style={axis line style={popink,line width=1pt},tick style={popink},
    grid style={popline,line width=0.5pt},label style={font=\small},tick label style={font=\footnotesize}},
  popcurve/.style={poporange,line width=1.8pt,no markers,smooth},
}
% Même style disponible dans un \draw TikZ simple (hors pgfplots)
\tikzset{popcurve/.style={poporange,line width=1.8pt}}
\tikzset{popgrid/.style={popline,very thin}}
\colorlet{popcurve}{poporange} % le nom du style est parfois utilisé comme couleur

% ============================================================
% Pill / squiggle
% ============================================================
\newcommand{\poppill}[3]{%
  \tikz[baseline=-0.55ex]\node[draw=popink,line width=1.2pt,rounded corners=2.6pt,%
    fill=#2,inner xsep=6.5pt,inner ysep=3pt]{%
    \popxboldls\fontsize{7}{7}\selectfont\color{#3}\MakeUppercase{#1}};%
}
\newcommand{\popsquiggle}[1]{%
  \tikz[baseline=-0.4ex]{
    \draw[#1,line width=2.6pt,line cap=round]
      plot [smooth] coordinates {(0,0.09) (0.34,0.01) (0.68,0.21) (1.02,0.01) (1.36,0.10)};
  }%
}
% Mot-clé surligné (marqueur orange sous le texte)
\newcommand{\cle}[1]{{\popxbold\color{popdark}#1}}

% ============================================================
% En-tête / pied
% ============================================================
\pagestyle{fancy}
\fancyhf{}
\renewcommand{\headrulewidth}{0pt}
\renewcommand{\footrulewidth}{0pt}
\lhead{\raisebox{-0.15ex}{\poplogo\fontsize{13}{13}\selectfont\color{popink}OnBuch\color{poporange}+}}
\rhead{\poppill{\DOCMATIERE\ \textbullet\ \DOCNIVEAU\ \textbullet\ Chapter \DOCLECON}{white}{popink}}
\lfoot{\popxbold\fontsize{7.6}{7.6}\selectfont\color{popmuted}\MakeUppercase{Signed course OnBuch+}\ \textbullet\ {\popsemi\DOCTITRE}}
\rfoot{\tikz[baseline=-0.6ex]\node[rounded corners=8.5pt,fill=popink,minimum height=17pt,minimum width=17pt,inner xsep=5pt,inner sep=0]{\popxbold\fontsize{8}{8}\selectfont\color{popcream}\thepage\,/\,\pageref*{LastPage}};}

% ============================================================
% Titres
% ============================================================
\newcommand{\chaptitre}[2]{%
  \begin{tcolorbox}[enhanced,colback=poporange,colframe=popink,boxrule=2.6pt,arc=11pt,
      drop shadow=popink,left=15pt,right=15pt,top=12pt,bottom=11pt,before skip=3pt,after skip=18pt]
    \poppill{#2}{popink}{popcream}\par\vspace{9pt}
    {\raggedright\hyphenpenalty=10000\exhyphenpenalty=10000\popdisplay\fontsize{22}{24}\selectfont\color{popcream}\MakeUppercase{#1}\par}
  \end{tcolorbox}
}
% Section numérotée : pastille orange avec le numéro + titre Archivo
\newcounter{popsec}
\newcommand{\coursec}[1]{%
  \stepcounter{popsec}\setcounter{popsub}{0}%
  \par\vspace{16pt}\noindent
  \begin{minipage}{\linewidth}
  \tikz[baseline=-0.7ex]\node[draw=popink,line width=1.6pt,fill=poporange,rounded corners=4pt,minimum size=22pt,inner sep=0]{\popdisplay\fontsize{12}{12}\selectfont\color{popcream}\thepopsec};%
  \hspace{8pt}{\popdisplay\fontsize{15}{17}\selectfont\color{popink}\MakeUppercase{#1}}\par
  \vspace{4pt}\noindent{\color{poporange}\rule{36pt}{3.6pt}}
  \end{minipage}\par\nopagebreak\vspace{8pt}
}
\newcounter{popsub}
\newcommand{\courssub}[1]{%
  \stepcounter{popsub}%
  \par\vspace{9pt}\noindent{\popxbold\fontsize{11.5}{13}\selectfont\color{popink}{\color{poporange}\thepopsec.\thepopsub}\hspace{6pt}#1}\par\nopagebreak\vspace{4pt}
}

% ============================================================
% Boîtes pédagogiques — même recette : fond blanc, bordure épaisse colorée,
% ombre dure encre, badge plein. Titre optionnel [#1].
% ============================================================
\tcbset{popbase/.style={breakable,enhanced,colback=white,boxrule=2.3pt,arc=9pt,
  shadow={3pt}{-3pt}{0pt}{popink},left=14pt,right=14pt,top=11pt,bottom=11pt,before skip=11pt,after skip=15pt,
  fontupper=\popsemi\fontsize{10.6}{14.5}\selectfont\color{popink},
  fontlower=\popsemi\fontsize{10.6}{14.5}\selectfont\color{popink}}}
% Coche dessinée (le glyphe ✓ n'existe pas dans les polices du document)
\newcommand{\popcheck}[1]{\tikz[baseline=-0.5ex]\draw[#1,line width=1.6pt,line cap=round,line join=round](-0.13,0.0)--(-0.04,-0.09)--(0.13,0.1);}
\providecommand{\coche}{\popcheck{popgreen}}
\providecommand{\resultat}[1]{\par\smallskip\noindent\fcolorbox{popgreen}{popgreenL}{\parbox{\dimexpr\linewidth-2\fboxsep-2\fboxrule\relax}{\textbf{Result.} #1}}\par\smallskip}
\newcommand{\popboxhead}[3]{\poppill{#1}{#2}{white}\ifx\relax#3\relax\else\hspace{6pt}{\popxbold\fontsize{10.6}{12}\selectfont\color{popink}#3}\fi\par\vspace{7pt}}

\newtcolorbox{aretenir}[1][]{popbase,colframe=popgreen,before upper={\popboxhead{Key points}{popgreen}{#1}}}
\newtcolorbox{exemplebox}[1][]{popbase,colframe=poporange,before upper={\popboxhead{Exemple}{poporange}{#1}}}
\newtcolorbox{definition}[1][]{popbase,colframe=popblue,before upper={\popboxhead{Definition}{popblue}{#1}}}
\newtcolorbox{propriete}[1][]{popbase,colframe=poppurple,before upper={\popboxhead{Property}{poppurple}{#1}}}
\newtcolorbox{methode}[1][]{popbase,colframe=popgold,before upper={\popboxhead{Method}{popgold}{#1}}}
\newtcolorbox{attention}[1][]{popbase,colframe=poppink,before upper={\popboxhead{Warning}{poppink}{#1}}}
\newtcolorbox{experience}[1][]{popbase,colframe=popblue,colback=popblueL,before upper={\popboxhead{Experiment}{popblue}{#1}}}
% « Le savais-tu ? » : carte violette pleine (contexte, culture, Cameroun)
\newtcolorbox{savaistu}[1][]{popbase,colback=poppurple,colframe=popink,
  fontupper=\popsemi\fontsize{10.4}{14.2}\selectfont\color{white},
  before upper={\poppill{Did you know?}{white}{poppurple}\ifx\relax#1\relax\else\hspace{6pt}{\popxbold\color{white}#1}\fi\par\vspace{7pt}}}
% Exercice résolu : énoncé en haut, \tcblower puis la solution sur fond crème
\newcounter{popexo}\newcounter{popexr}
\newtcolorbox{exoresolu}[1][]{popbase,colframe=poporange,colbacklower=popcream,
  segmentation style={popink,line width=1.2pt,dashed},
  before upper={\stepcounter{popexr}\popboxhead{Worked example \thepopexr}{poporange}{#1}},
  before lower={\poppill{Solution}{popink}{popcream}\par\vspace{6pt}}}
% Exercice (non résolu en place) — la correction est dans la partie Corrigés
\newtcolorbox{exercice}[1][]{popbase,colframe=popink,
  before upper={\stepcounter{popexo}\popboxhead{Exercise \thepopexo}{popink}{#1}}}
% Corrigé
\newtcolorbox{corrige}[1][]{popbase,colframe=popgreen,colback=popgreenL,
  before upper={\popboxhead{Solution}{popgreen}{#1}}}
% Objectifs / prérequis / bilan
\newtcolorbox{objectifs}{popbase,colframe=popink,colback=poporangeL,
  before upper={\popboxhead{Chapter objectives}{popink}{}}}
\newtcolorbox{prerequis}{popbase,colframe=popink,colback=popblueL,
  before upper={\popboxhead{Prerequisites}{popblue}{}}}
\newtcolorbox{bilan}{popbase,colframe=popink,colback=white,boxrule=2.8pt,
  before upper={\popboxhead{Fiche bilan}{poporange}{L'essentiel à connaître pour l'examen}}}
% Figure encadrée avec légende
\newtcolorbox{popfigure}[1][]{enhanced,breakable=false,colback=white,colframe=popline,boxrule=1.4pt,arc=8pt,
  left=10pt,right=10pt,top=10pt,bottom=8pt,before skip=10pt,after skip=12pt,halign=center,
  lower separated=false,halign lower=center,
  fontlower=\popsemi\fontsize{9}{11.5}\selectfont\color{popmuted},
  #1}
\newcommand{\legende}[1]{\par\vspace{6pt}{\popsemi\fontsize{9}{11.5}\selectfont\color{popmuted}#1}}

% ============================================================
% Signature OnBuch+
% ============================================================
\newcommand{\poplogoinline}[1]{{\poplogo\fontsize{#1}{#1}\selectfont\color{popink}OnBuch\color{poporange}+}}
\newcommand{\signatureonbuch}{%
  \begin{tcolorbox}[enhanced,colback=popink,colframe=popink,boxrule=2.2pt,arc=10pt,
      drop shadow=poporange,left=14pt,right=14pt,top=10pt,bottom=10pt,before skip=0pt,after skip=0pt]
    \tikz[baseline=-0.7ex]\node[circle,fill=popgreen,draw=popcream,line width=1.3pt,minimum size=22pt,inner sep=0]{\popcheck{white}};%
    \hspace{9pt}\begin{minipage}[c]{0.62\linewidth}
      {\popxbold\fontsize{12}{13}\selectfont\color{popcream}Signed\ \textbullet\ {\poplogo OnBuch\color{poporange}+}}\par\vspace{2pt}
      {\raggedright\popsemi\fontsize{8}{10.5}\selectfont\color{popline}\MakeUppercase{OnBuch+ teaching team}\\\MakeUppercase{Follows the official MINESEC syllabus}\par}
    \end{minipage}\hfill
    {\poplogo\fontsize{22}{22}\selectfont\color{popcream}OnBuch\color{poporange}+}
  \end{tcolorbox}%
}

% ============================================================
% Page de garde
% #1 titre · #2 matière · #3 niveau · #4 module · #5 n° leçon · #6 durée ·
% #7 description / objectifs (texte ou itemize)
% ============================================================
\newcommand{\pagegarde}[7]{%
  \thispagestyle{empty}
  \newgeometry{a4paper,margin=1.7cm,top=1.6cm,bottom=1.4cm}%
  \begin{tikzpicture}[remember picture,overlay]
    \fill[poporange!18] ([xshift=-3.2cm,yshift=-3.6cm]current page.north east) circle (4.6cm);
    \fill[poppurple!14] ([xshift=2.4cm,yshift=7.6cm]current page.south west) circle (3.8cm);
  \end{tikzpicture}%
  {\poplogo\fontsize{30}{30}\selectfont\color{popink}OnBuch\color{poporange}+}\hfill
  \raisebox{0.6ex}{\poppill{Signed course \textbullet\ Premium}{popink}{popcream}}\par
  \vspace{1pt}\popsquiggle{poppurple}\par
  \vspace{8pt}
  \begin{tcolorbox}[enhanced,colback=poporange,colframe=popink,boxrule=2.8pt,arc=13pt,
      drop shadow=popink,left=18pt,right=18pt,top=12pt,bottom=13pt,after skip=10pt]
    \poppill{#2 \textbullet\ #3}{popink}{popcream}\hspace{5pt}\poppill{#4}{white}{popink}\par\vspace{8pt}
    {\popdisplay\fontsize{14}{14}\selectfont\color{popink}CHAPTER #5}\par\vspace{4pt}
    {\raggedright\hyphenpenalty=10000\exhyphenpenalty=10000\popdisplay\fontsize{25}{27}\selectfont\color{popcream}\MakeUppercase{#1}\par}
    \vspace{8pt}
    {\popxbold\fontsize{9.5}{9.5}\selectfont\color{popink}\MakeUppercase{Suggested time: #6}}
  \end{tcolorbox}
  \begin{tcolorbox}[enhanced,colback=white,colframe=popink,boxrule=2.2pt,arc=10pt,
      drop shadow=popink,left=16pt,right=16pt,top=10pt,bottom=10pt,after skip=8pt]
    \poppill{In this chapter}{poppurple}{white}\par\vspace{5pt}
    \popsemi\fontsize{9.3}{12.6}\selectfont\color{popink}#7
  \end{tcolorbox}
  \noindent\begin{minipage}[t]{0.487\linewidth}
    \begin{tcolorbox}[enhanced,colback=popgreen,colframe=popink,boxrule=2.2pt,arc=10pt,
        drop shadow=popink,left=11pt,right=11pt,top=10pt,bottom=10pt,height=3.1cm,valign=center]
      \tcbox[colback=white,colframe=popink,boxrule=1.3pt,arc=5pt,boxsep=0pt,nobeforeafter,
        left=3pt,right=3pt,top=3pt,bottom=3pt,valign=center]{\qrcode[height=1.9cm]{https://play.google.com/store/apps/details?id=com.luvvix.onbuch&hl=en}}%
      \hspace{8pt}\begin{minipage}[c]{0.5\linewidth}
        {\popdisplay\fontsize{11}{12.5}\selectfont\color{white}\MakeUppercase{Download OnBuch}}\par\vspace{4pt}
        {\raggedright\popsemi\fontsize{8.4}{11}\selectfont\color{white}Free \textbullet\ Google Play\\AI tutor, past papers, results\par}
      \end{minipage}
    \end{tcolorbox}
  \end{minipage}\hfill
  \begin{minipage}[t]{0.487\linewidth}
    \begin{tcolorbox}[enhanced,colback=poppurple,colframe=popink,boxrule=2.2pt,arc=10pt,
        drop shadow=popink,left=11pt,right=11pt,top=10pt,bottom=10pt,height=3.1cm,valign=center]
      \tikz[baseline=-0.7ex]\node[circle,fill=white,draw=popink,line width=1.3pt,minimum size=1.6cm,inner sep=0]{\popdisplay\fontsize{24}{24}\selectfont\color{poppurple}+};%
      \hspace{8pt}\begin{minipage}[c]{0.6\linewidth}
        {\popdisplay\fontsize{11}{12.5}\selectfont\color{white}\MakeUppercase{Go OnBuch+}}\par\vspace{4pt}
        {\raggedright\popsemi\fontsize{8.4}{11}\selectfont\color{white}All signed courses, unlimited AI tutor, PDF sheets — OnBuch+ tab in the app\par}
      \end{minipage}
    \end{tcolorbox}
  \end{minipage}
  \vspace{8pt}
  \popcontenu
  \vfill
  \signatureonbuch
  \restoregeometry
  \newpage
}

% Rangée « Ce cours contient » (page de garde)
\newcommand{\poptile}[3]{% #1 couleur #2 gros texte #3 légende
  \begin{tcolorbox}[enhanced,colback=#1,colframe=popink,boxrule=2pt,arc=9pt,shadow={2.5pt}{-2.5pt}{0pt}{popink},
      left=6pt,right=6pt,top=9pt,bottom=9pt,halign=center,valign=center,height=1.75cm,width=\linewidth,nobeforeafter]
    {\popdisplay\fontsize{12.5}{13}\selectfont\color{white}\MakeUppercase{#2}}\par\vspace{3pt}
    {\centering\popsemi\fontsize{7.8}{9.5}\selectfont\color{white}#3\par}
  \end{tcolorbox}}
\newcommand{\popcontenu}{%
  \par\noindent{\popxboldls\fontsize{7.5}{7.5}\selectfont\color{popmuted}THIS SIGNED COURSE CONTAINS}\par\vspace{7pt}
  \noindent\begin{minipage}[t]{0.235\linewidth}\poptile{popblue}{Course}{complete, illustrated, step by step}\end{minipage}\hfill
  \begin{minipage}[t]{0.235\linewidth}\poptile{poporange}{Methods}{and worked examples}\end{minipage}\hfill
  \begin{minipage}[t]{0.235\linewidth}\poptile{poppink}{Exercises}{exam style + solutions}\end{minipage}\hfill
  \begin{minipage}[t]{0.235\linewidth}\poptile{popgreen}{Summary sheet}{the essentials to remember}\end{minipage}\par}

% Sommaire
\newtcolorbox{sommaire}{enhanced,colback=white,colframe=popink,boxrule=2.2pt,arc=10pt,
  drop shadow=popink,left=18pt,right=18pt,top=13pt,bottom=13pt,before skip=6pt,after skip=20pt,
  before upper={\poppill{Contents}{poppurple}{white}\par\vspace{9pt}\popsemi\fontsize{10.6}{14.5}\selectfont\color{popink}}}

% Signature de fin de cours — poussée tout en bas de la dernière page
\newcommand{\findecours}{%
  \par\vfill\nopagebreak
  \begin{center}\popsquiggle{poporange}\end{center}\vspace{4pt}
  \signatureonbuch
}
\AtBeginDocument{\def\llbracket{\mathopen{[\mkern-3.5mu[}}\def\rrbracket{\mathclose{]\mkern-3.5mu]}}} % intervalles d'entiers (glyphe ⟦ absent de la police)
% Glyphes absents de Fira Math : redéfinis à partir de glyphes disponibles
\AtBeginDocument{%
  \def\setminus{\mathbin{\backslash}}\let\smallsetminus\setminus
  \def\vdots{\mathord{\vbox{\baselineskip3.2pt\lineskiplimit0pt\kern4pt\hbox{.}\hbox{.}\hbox{.}}}}%
  \def\ddots{\mathinner{\mkern1mu\raise6.5pt\hbox{.}\mkern2mu\raise3.6pt\hbox{.}\mkern2mu\raise0.7pt\hbox{.}\mkern1mu}}%
  \def\triangle{\mathord{\scalebox{0.9}{$\symup{\Delta}$}}}%
  \def\nabla{\mathord{\raisebox{\depth}{\scalebox{1}[-1]{$\symup{\Delta}$}}}}%
  \def\checkmark{\popcheck{popdark}}%
  \def\star{\mathbin{\ast}}\def\bigstar{\mathord{\ast}}%
  \def\diamond{\mathbin{\scalebox{0.6}{$\Diamond$}}}%
  \def\bigcup{\mathop{\mathchoice{\raisebox{-0.3em}{\scalebox{1.7}{$\cup$}}}{\scalebox{1.25}{$\cup$}}{\cup}{\cup}}\displaylimits}%
  \def\bigcap{\mathop{\mathchoice{\raisebox{-0.3em}{\scalebox{1.7}{$\cap$}}}{\scalebox{1.25}{$\cap$}}{\cap}{\cap}}\displaylimits}%
  \def\lesseqqgtr{\mathrel{\lessgtr}}\def\bigoplus{\mathop{\oplus}\displaylimits}%
}
\providecommand{\ssi}{if and only if} % abréviation fréquente
\AtBeginDocument{\providecommand{\wideparen}{\overparen}} % arc de cercle

% Fonctions usuelles que les modèles utilisent sans que amsmath les fournisse
\providecommand{\arsinh}{\operatorname{arsinh}}\providecommand{\arcosh}{\operatorname{arcosh}}
\providecommand{\artanh}{\operatorname{artanh}}\providecommand{\arsech}{\operatorname{arsech}}
\providecommand{\arcsch}{\operatorname{arcsch}}\providecommand{\arcoth}{\operatorname{arcoth}}
\providecommand{\arcsinh}{\operatorname{arsinh}}\providecommand{\arccosh}{\operatorname{arcosh}}
\providecommand{\arctanh}{\operatorname{artanh}}\providecommand{\sech}{\operatorname{sech}}
\providecommand{\csch}{\operatorname{csch}}\providecommand{\arcsec}{\operatorname{arcsec}}
\providecommand{\arccsc}{\operatorname{arccsc}}\providecommand{\arccot}{\operatorname{arccot}}
\providecommand{\sgn}{\operatorname{sgn}}\providecommand{\lcm}{\operatorname{lcm}}
\providecommand{\Var}{\operatorname{Var}}\providecommand{\Cov}{\operatorname{Cov}}

%% FIGFIX-BEGIN v5 — correctifs globaux des figures (docs/onbuch-generation/tools/figfix)
\usepackage{adjustbox}\usepackage{etoolbox}\tcbuselibrary{raster}
% toute figure TikZ/pgfplots plus large que la ligne est réduite à la largeur disponible
\BeforeBeginEnvironment{tikzpicture}{\begin{adjustbox}{max width=\linewidth}}
\AfterEndEnvironment{tikzpicture}{\end{adjustbox}}
% légendes pgfplots lisibles par-dessus les courbes
\pgfplotsset{every axis/.append style={legend style={fill=white,fill opacity=0.92,text opacity=1,draw=black!15,inner sep=3pt}}}
% étiquettes d'arêtes (syntaxe edge["..."]) avec fond blanc
\tikzset{every edge quotes/.append style={fill=white,inner sep=1.2pt}}
%% FIGFIX-END
\begin{document}

\pagegarde{Special Discrete Probability Distributions}{Further Mathematics}{Upper Sixth}{Upper Sixth — Further mathematics}{18}{11 periods}{This chapter studies the four fundamental discrete probability distributions of the GCE Further Mathematics syllabus: binomial, uniform, geometric and Poisson. For each one, we derive the mean and variance, learn to use cumulative tables, and apply the models to real-life counting problems, including the Poisson approximation of the binomial.
\vspace{4pt}
\begin{multicols}{2}\small
\begin{itemize}
\item By the end of this chapter I can recognise a binomial situation and state its parameters n and p.
\item By the end of this chapter I can compute binomial probabilities and derive/use E(X)=np and Var(X)=npq.
\item By the end of this chapter I can read cumulative binomial probability tables to find P(X≤r), P(X≥r) and P(a≤X≤b).
\item By the end of this chapter I can use the recurrence formula P(X=r+1)/P(X=r) to locate the mode of a binomial distribution.
\item By the end of this chapter I can define the discrete uniform distribution and derive its expectation (n+1)/2 and variance (n²−1)/12.
\item By the end of this chapter I can model 'waiting time to first success' with the geometric distribution and use E(X)=1/p and Var(X)=q/p².
\end{itemize}\end{multicols}}

\begin{sommaire}
\begin{multicols}{2}
\begin{enumerate}
\item The Binomial Distribution: Probabilities, Mean and Variance
\item Cumulative Binomial Tables and the Mode
\item The Discrete Uniform Distribution
\item The Geometric Distribution
\item The Poisson Distribution and Its Tables
\item Poisson Approximation to the Binomial and Link with Frequency Distributions
\item Integration activity
\item Methods and worked examples
\item Exercises
\item Solutions to the exercises
\item Summary sheet
\end{enumerate}
\end{multicols}
\end{sommaire}

\begin{objectifs}
\begin{itemize}
\item By the end of this chapter I can recognise a binomial situation and state its parameters n and p.
\item By the end of this chapter I can compute binomial probabilities and derive/use E(X)=np and Var(X)=npq.
\item By the end of this chapter I can read cumulative binomial probability tables to find P(X≤r), P(X≥r) and P(a≤X≤b).
\item By the end of this chapter I can use the recurrence formula P(X=r+1)/P(X=r) to locate the mode of a binomial distribution.
\item By the end of this chapter I can define the discrete uniform distribution and derive its expectation (n+1)/2 and variance (n²−1)/12.
\item By the end of this chapter I can model 'waiting time to first success' with the geometric distribution and use E(X)=1/p and Var(X)=q/p².
\item By the end of this chapter I can model rare random events with the Poisson distribution, use E(X)=Var(X)=λ, and adjust λ to the unit interval.
\item By the end of this chapter I can approximate a binomial distribution by a Poisson distribution when n is large and p is small (np moderate).
\item By the end of this chapter I can use cumulative Poisson tables and the additive property of independent Poisson variables.
\item By the end of this chapter I can relate a probability distribution to an observed frequency distribution and compare theoretical and experimental means.
\end{itemize}
\end{objectifs}

\begin{prerequis}
\begin{itemize}
\item Probability of events, independence and the binomial coefficients nCr (Lower Sixth Pure).
\item Random variables, probability distributions, expectation E(X) and variance Var(X) (Lower Sixth Statistics).
\item Summation of series: sum of first n integers, sum of squares, geometric series (Lower Sixth Pure).
\item The exponential function e\^{}x and its use in limits (Lower Sixth Pure).
\item Factorials and manipulation of algebraic fractions.
\end{itemize}
\end{prerequis}

\newpage

\coursec{The Binomial Distribution: Probabilities, Mean and Variance}

At the Mfoundi market in Yaoundé, a trader who sells 200 avocados a day knows from experience that about 3 out of every 10 are overripe and unsellable. How many bad avocados should she expect each morning, and how likely is she to lose more than 70 in one day? Meanwhile, the MTN call centre in Douala receives on average 4 calls per minute: what is the probability of receiving no call at all in the next minute? And a voter queue at a polling station in Bamenda raises another question: how many people must be checked before the first registered voter is found? These everyday situations are modelled by the \cle{special discrete distributions} of this chapter: \cle{binomial}, \cle{uniform}, \cle{geometric} and \cle{Poisson}. In this first section we build the most fundamental of them all, the binomial distribution, from the ground up: we define it, compute its probabilities, and derive its mean and variance rigorously.

\courssub{Bernoulli Trials: the Building Blocks}

Imagine tossing a coin once and recording only whether it lands Heads. Or inspecting one avocado and recording only whether it is overripe. Each of these experiments has exactly \emph{two} possible outcomes, which we conventionally call \cle{success} and \cle{failure}. Such an experiment is called a \cle{Bernoulli trial}.

\begin{definition}[Bernoulli trial]
A \textbf{Bernoulli trial} is a random experiment with exactly two outcomes: \emph{success} (S), with probability $p$, and \emph{failure} (F), with probability $q = 1 - p$. A sequence of Bernoulli trials is a sequence of \textbf{independent} repetitions of this experiment, with $p$ staying \textbf{constant} from trial to trial.
\end{definition}

The word ``success'' is a label, not a value judgement: if we are counting overripe avocados, then ``success'' means ``the avocado is overripe''. What matters mathematically is the structure, not the story.

\begin{methode}[Recognising a binomial situation]
To decide whether a random variable $X$ follows a binomial distribution, check the \textbf{four conditions}:
\begin{enumerate}
\item \textbf{Fixed number of trials:} the experiment is repeated a fixed number $n$ of times.
\item \textbf{Two outcomes:} each trial results in success or failure only.
\item \textbf{Constant probability:} the probability of success $p$ is the same for every trial.
\item \textbf{Independence:} the outcome of one trial does not affect the others.
\end{enumerate}
If all four hold and $X$ counts the number of successes, then $X$ is binomial.
\end{methode}

\begin{attention}[Sampling without replacement]
Drawing items from a small finite population \emph{without replacement} breaks the independence and constant-$p$ conditions. For example, picking 5 students from a class of 30 and counting the number of girls is \textbf{not} exactly binomial (it is hypergeometric). However, when the population is very large compared with the sample (e.g.\ 5 bolts from a production line of thousands), the binomial model is an excellent approximation and examiners accept it.
\end{attention}

\courssub{Definition and Probability Formula of $X \sim B(n,p)$}

Suppose we perform $n = 3$ independent Bernoulli trials with $P(S) = p$, and we want the probability of exactly $2$ successes. The tree diagram below shows all $2^3 = 8$ possible outcomes.

\begin{popfigure}
\begin{center}
\begin{tikzpicture}[grow=right, level distance=2.6cm,
level 1/.style={sibling distance=3.4cm},
level 2/.style={sibling distance=1.8cm},
level 3/.style={sibling distance=0.95cm},
every node/.style={font=\small}]
\node {Start}
child { node {F ($q$)}
  child { node {F ($q$)}
    child { node {F} edge from parent node[above] {$q$} }
    child { node {S} edge from parent node[above] {$p$} }
    edge from parent node[above] {$q$} }
  child { node {S ($p$)}
    child { node {F} edge from parent node[above] {$q$} }
    child { node {S} edge from parent node[above] {$p$} }
    edge from parent node[above] {$p$} }
  edge from parent node[above] {$q$} }
child { node {S ($p$)}
  child { node {F ($q$)}
    child { node {F} edge from parent node[above] {$q$} }
    child { node {S} edge from parent node[above] {$p$} }
    edge from parent node[above] {$q$} }
  child { node {S ($p$)}
    child { node {F} edge from parent node[above] {$q$} }
    child { node {S} edge from parent node[above] {$p$} }
    edge from parent node[above] {$p$} }
  edge from parent node[above] {$p$} };
\end{tikzpicture}
\end{center}
\legende{Tree diagram of 3 Bernoulli trials. Exactly 2 successes occur along the paths SSF, SFS and FSS: there are $\binom{3}{2}=3$ such paths, each of probability $p^2 q$.}
\end{popfigure}

Each path containing exactly 2 successes and 1 failure has probability $p^2 q$ (multiplication rule, by independence), and the number of such paths is the number of ways of choosing which 2 of the 3 trials are successes, namely $\binom{3}{2} = 3$. Hence
\[
P(\text{exactly 2 successes}) = \binom{3}{2} p^2 q.
\]
This reasoning generalises immediately.

\begin{definition}[Binomial distribution]
A discrete random variable $X$ follows a \textbf{binomial distribution} with parameters $n$ and $p$, written $X \sim B(n,p)$, if $X$ counts the number of successes in $n$ independent Bernoulli trials with $P(\text{success}) = p$. Its probability mass function is
\[
P(X = r) = \binom{n}{r} p^r q^{\,n-r}, \qquad r = 0, 1, 2, \dots, n, \quad q = 1 - p.
\]
\end{definition}

The name ``binomial'' comes from the binomial theorem: summing all the probabilities gives
\[
\sum_{r=0}^{n} \binom{n}{r} p^r q^{n-r} = (p + q)^n = 1^n = 1,
\]
which confirms that the formula defines a genuine probability distribution.

\begin{attention}[Do not forget the coefficient]
A very common mistake is to write $P(X=r) = p^r q^{n-r}$, forgetting $\binom{n}{r}$. The factor $p^r q^{n-r}$ is the probability of \emph{one particular arrangement} of $r$ successes and $n-r$ failures; you must multiply by the number of arrangements, $\binom{n}{r}$.
\end{attention}

\begin{exemplebox}[Building the full distribution table]
A fair coin is tossed 4 times and $X$ counts the number of Heads. Then $X \sim B(4, 0.5)$ and $P(X=r) = \binom{4}{r}(0.5)^4 = \binom{4}{r}/16$:
\begin{center}
\begin{tabular}{|c|ccccc|}
\hline
\rowcolor{popblueL} $r$ & $0$ & $1$ & $2$ & $3$ & $4$ \\
\hline
$P(X=r)$ & $\frac{1}{16}$ & $\frac{4}{16}$ & $\frac{6}{16}$ & $\frac{4}{16}$ & $\frac{1}{16}$ \\
\hline
\end{tabular}
\end{center}
Check: $\frac{1+4+6+4+1}{16} = 1$. \checkmark
\end{exemplebox}

\begin{popfigure}
\begin{center}
\begin{tikzpicture}
\begin{axis}[
ybar, bar width=10pt,
width=12cm, height=6.5cm,
xlabel={Number of successes $r$},
ylabel={$P(X=r)$},
xmin=-0.7, xmax=10.7,
ymin=0, ymax=0.28,
xtick={0,1,2,3,4,5,6,7,8,9,10},
axis lines=left,
ymajorgrids=true, grid style={popline},
every axis plot/.append style={fill=popblue, draw=popdark},
]
\addplot coordinates {
(0,0.0282) (1,0.1211) (2,0.2335) (3,0.2668) (4,0.2001)
(5,0.1029) (6,0.0368) (7,0.0090) (8,0.0014) (9,0.0001) (10,0.0000)
};
\end{axis}
\end{tikzpicture}
\end{center}
\legende{Probability mass function of $X \sim B(10, 0.3)$. The distribution is skewed to the right because $p < 0.5$; the most likely value is $r = 3 = np$.}
\end{popfigure}

The bar chart of $B(10, 0.3)$ shows a typical shape: the probabilities rise to a peak near $np$ and then fall away. When $p = 0.5$ the distribution is symmetric; when $p < 0.5$ it has a tail to the right; when $p > 0.5$ a tail to the left.

\courssub{Derivation of the Mean: $E(X) = np$}

Intuition first: if 3 avocados out of 10 are bad, then out of 200 avocados we ``expect'' $200 \times 0.3 = 60$ bad ones. Expectation should be $n$ times $p$. Let us prove it.

\begin{propriete}[Mean of the binomial distribution — proof examinable]
If $X \sim B(n,p)$, then
\[
E(X) = np.
\]
\end{propriete}

\textbf{Proof.} By definition, $E(X) = \displaystyle\sum_{r=0}^{n} r\, P(X=r) = \sum_{r=0}^{n} r \binom{n}{r} p^r q^{n-r}$. The term $r = 0$ vanishes. For $r \geq 1$ we use the key identity
\[
r \binom{n}{r} = r \cdot \frac{n!}{r!(n-r)!} = \frac{n!}{(r-1)!(n-r)!} = n \cdot \frac{(n-1)!}{(r-1)!\,(n-r)!} = n \binom{n-1}{r-1}.
\]
Therefore
\begin{align*}
E(X) &= \sum_{r=1}^{n} n \binom{n-1}{r-1} p^r q^{n-r}
= np \sum_{r=1}^{n} \binom{n-1}{r-1} p^{r-1} q^{(n-1)-(r-1)}.
\end{align*}
Putting $s = r - 1$, the sum becomes $\displaystyle\sum_{s=0}^{n-1} \binom{n-1}{s} p^s q^{(n-1)-s} = (p+q)^{n-1} = 1$. Hence $E(X) = np$. $\blacksquare$

\courssub{Derivation of the Variance: $\operatorname{Var}(X) = npq$}

For the variance we use $\operatorname{Var}(X) = E(X^2) - [E(X)]^2$, but $E(X^2)$ is awkward to compute directly. The trick is to compute the \emph{factorial moment} $E\big(X(X-1)\big)$ first, because $r(r-1)\binom{n}{r} = n(n-1)\binom{n-2}{r-2}$.

\begin{propriete}[Variance of the binomial distribution — proof examinable]
If $X \sim B(n,p)$, then
\[
\operatorname{Var}(X) = npq, \qquad \text{and} \qquad \sigma = \sqrt{npq}.
\]
\end{propriete}

\textbf{Proof.} First,
\[
E\big(X(X-1)\big) = \sum_{r=0}^{n} r(r-1)\binom{n}{r} p^r q^{n-r}.
\]
The terms $r = 0$ and $r = 1$ vanish. For $r \geq 2$,
\[
r(r-1)\binom{n}{r} = \frac{n!}{(r-2)!(n-r)!} = n(n-1)\binom{n-2}{r-2}.
\]
Hence, factoring out $n(n-1)p^2$ and setting $s = r-2$,
\begin{align*}
E\big(X(X-1)\big) &= n(n-1)p^2 \sum_{r=2}^{n} \binom{n-2}{r-2} p^{r-2} q^{n-r} \\
&= n(n-1)p^2 \sum_{s=0}^{n-2} \binom{n-2}{s} p^s q^{(n-2)-s} = n(n-1)p^2 (p+q)^{n-2} = n(n-1)p^2.
\end{align*}
Now $E(X^2) = E\big(X(X-1)\big) + E(X) = n(n-1)p^2 + np$, so
\begin{align*}
\operatorname{Var}(X) &= n(n-1)p^2 + np - (np)^2 = n^2p^2 - np^2 + np - n^2p^2 \\
&= np(1 - p) = npq. \qquad \blacksquare
\end{align*}

\begin{aretenir}[Key results]
If $X \sim B(n, p)$ with $q = 1 - p$:
\[
P(X = r) = \binom{n}{r} p^r q^{n-r}, \qquad E(X) = np, \qquad \operatorname{Var}(X) = npq.
\]
The mean tells you the \emph{expected number of successes}; the variance measures the \emph{spread} around that expectation. The variance is greatest when $p = 0.5$ (maximum uncertainty) and tends to $0$ as $p \to 0$ or $p \to 1$.
\end{aretenir}

\begin{exoresolu}[The avocado trader]
The Mfoundi trader sells $n = 200$ avocados daily, each independently overripe with probability $p = 0.3$. Let $X$ be the number of unsellable avocados.
\begin{enumerate}
\item State the distribution of $X$.
\item Find the expected number of bad avocados and the standard deviation.
\item Compute $P(X = 60)$.
\end{enumerate}
\tcblower
\textbf{Solution.}
\begin{enumerate}
\item Fixed $n = 200$ trials, two outcomes (overripe or not), constant $p = 0.3$, independence assumed: $X \sim B(200, 0.3)$.
\item $E(X) = np = 200 \times 0.3 = 60$ bad avocados per day. $\operatorname{Var}(X) = npq = 200 \times 0.3 \times 0.7 = 42$, so $\sigma = \sqrt{42} \approx 6.48$. She should expect about 60 bad avocados, typically fluctuating by roughly 6 or 7 from day to day.
\item $P(X = 60) = \dbinom{200}{60}(0.3)^{60}(0.7)^{140} \approx 0.0615$. Even the \emph{most likely} single value has only about a $6\%$ chance — with 201 possible values, the probability is spread thinly.
\end{enumerate}
\end{exoresolu}

\courssub{Finding $n$ and $p$ from the Mean and Variance}

In many GCE questions, $n$ and $p$ are unknown but $E(X)$ and $\operatorname{Var}(X)$ are given. This yields two simultaneous equations:
\[
np = \mu, \qquad npq = \sigma^2.
\]

\begin{methode}[Recovering $n$ and $p$]
\begin{enumerate}
\item Divide the variance by the mean: $\dfrac{npq}{np} = q = \dfrac{\sigma^2}{\mu}$.
\item Deduce $p = 1 - q$.
\item Deduce $n = \dfrac{\mu}{p}$ (check that $n$ is a positive integer).
\end{enumerate}
\end{methode}

\begin{exoresolu}[Unknown $n$ and $p$]
A binomial random variable $X$ has mean $6$ and variance $4.2$. Find $n$, $p$ and $P(X = 2)$.
\tcblower
\textbf{Solution.} We have $np = 6$ and $npq = 4.2$.
\[
q = \frac{npq}{np} = \frac{4.2}{6} = 0.7 \quad\Longrightarrow\quad p = 1 - 0.7 = 0.3, \qquad n = \frac{6}{0.3} = 20.
\]
So $X \sim B(20, 0.3)$ and
\[
P(X = 2) = \binom{20}{2}(0.3)^2(0.7)^{18} = 190 \times 0.09 \times (0.7)^{18} \approx 0.0278.
\]
\end{exoresolu}

\begin{attention}[Sanity checks]
Always verify that your answers make sense: $0 < p < 1$, $n$ a positive integer, and $\operatorname{Var}(X) < E(X)$ for a binomial (since $q < 1$). If a question gives a variance \emph{larger} than the mean, the variable cannot be binomial — this observation will matter when we meet the Poisson distribution, where mean and variance are equal.
\end{attention}

\begin{exercice}[Exercises on the Binomial Distribution]
\begin{enumerate}
\item A fair die is rolled 6 times. Let $X$ be the number of sixes obtained. State the distribution of $X$ and compute $P(X = 2)$, $E(X)$ and $\operatorname{Var}(X)$.
\item It is known that $70\%$ of candidates in a certain school pass the GCE Advanced Level. A random sample of 8 candidates is taken. Find the probability that exactly 5 pass, and the expected number of passes.
\item A binomial variable has mean $12$ and variance $3$. Find $n$, $p$ and $P(X = 12)$.
\end{enumerate}
\end{exercice}

\begin{corrige}[Exercise 1]
\begin{enumerate}
\item $X \sim B(6, \tfrac16)$; $P(X=2) = \binom{6}{2}\left(\tfrac16\right)^2\left(\tfrac56\right)^4 = 15 \times \tfrac{1}{36} \times \tfrac{625}{1296} \approx 0.2009$; $E(X) = 6 \times \tfrac16 = 1$; $\operatorname{Var}(X) = 6 \times \tfrac16 \times \tfrac56 = \tfrac56 \approx 0.833$.
\item $X \sim B(8, 0.7)$; $P(X=5) = \binom{8}{5}(0.7)^5(0.3)^3 = 56 \times 0.16807 \times 0.027 \approx 0.2541$; $E(X) = 5.6$.
\item $q = \tfrac{3}{12} = 0.25$, so $p = 0.75$ and $n = \tfrac{12}{0.75} = 16$. $P(X=12) = \binom{16}{12}(0.75)^{12}(0.25)^4 = 1820 \times (0.75)^{12} \times \tfrac{1}{256} \approx 0.2252$.
\end{enumerate}
\end{corrige}

\begin{savaistu}[Jacob Bernoulli]
The Bernoulli trial is named after the Swiss mathematician Jacob Bernoulli (1655--1705), whose book \emph{Ars Conjectandi} laid the foundations of probability theory. The same family produced no fewer than eight distinguished mathematicians. The binomial distribution you have just studied is the engine behind opinion polls, quality control in factories, and medical trials — including the analysis of examination pass rates in Cameroonian schools.
\end{savaistu}

\coursec{Cumulative Binomial Tables and the Mode}

In the previous section we derived the probability mass function of the binomial distribution $B(n,p)$ and computed its mean $\mu=np$ and variance $\sigma^2=npq$.  While the formula $P(X=r)=\binom{n}{r}p^r q^{n-r}$ is exact, evaluating it by hand for several values of $r$ becomes tedious, especially when $n$ is large.  In examinations and in practice, statisticians use \textbf{cumulative binomial tables} which give $P(X\le r)$ directly.  This section teaches you how to read those tables efficiently, how to handle probabilities that are not directly tabulated, and how to locate the mode of the distribution using a simple recurrence relation.

\courssub{Structure of Cumulative Binomial Tables}

Standard GCE cumulative binomial tables are organised as follows:
\begin{itemize}
  \item Rows correspond to the number of trials $n$ (typically $n=1,2,\dots,50$).
  \item Columns correspond to the success probability $p$ (usually $p=0.05,0.10,\dots,0.50$).
  \item Inside the table, for a given $n$ and $p$, you find a list of values $P(X\le r)$ for $r=0,1,\dots,n$.
\end{itemize}
Only probabilities for $p\le 0.5$ are printed.  If $p>0.5$ we use the symmetry of the binomial distribution (see below).

\begin{popfigure}
\centering
\begin{tikzpicture}[scale=0.75, every node/.style={font=\small}]
  % Table header
  \node[draw, fill=popblueL, minimum width=1.1cm, minimum height=0.65cm] (n) at (0,0) {$n$};
  \node[draw, fill=popblueL, minimum width=1.1cm, minimum height=0.65cm] (r) at (1.1,0) {$r$};
  \node[draw, fill=popblueL, minimum width=2.2cm, minimum height=0.65cm] (p05) at (2.2,0) {$p=0.05$};
  \node[draw, fill=popblueL, minimum width=2.2cm, minimum height=0.65cm] (p10) at (4.4,0) {$p=0.10$};
  \node[draw, fill=popblueL, minimum width=2.2cm, minimum height=0.65cm] (p20) at (6.6,0) {$p=0.20$};
  \node[draw, fill=popblueL, minimum width=2.2cm, minimum height=0.65cm] (p30) at (8.8,0) {$p=0.30$};
  \node[draw, fill=popblueL, minimum width=2.2cm, minimum height=0.65cm] (p40) at (11.0,0) {$p=0.40$};
  \node[draw, fill=popblueL, minimum width=2.2cm, minimum height=0.65cm] (p50) at (13.2,0) {$p=0.50$};

  % Data rows for n=8
  \foreach \rval/\pvalA/\pvalB/\pvalC/\pvalD/\pvalE/\pvalF in {
    0/0.6634/0.4305/0.1678/0.0576/0.0168/0.0039,
    1/0.9428/0.8131/0.5033/0.2553/0.1064/0.0352,
    2/0.9942/0.9619/0.7969/0.5518/0.3154/0.1445,
    3/0.9994/0.9950/0.9437/0.8059/0.5941/0.3633,
    4/0.9999/0.9995/0.9896/0.9420/0.8263/0.6367,
    5/1.0000/0.9999/0.9988/0.9887/0.9502/0.8555,
    6/1.0000/1.0000/0.9999/0.9988/0.9930/0.9648,
    7/1.0000/1.0000/1.0000/0.9999/0.9997/0.9961,
    8/1.0000/1.0000/1.0000/1.0000/1.0000/1.0000
  }{
    \pgfmathsetmacro{\y}{-0.65 - \rval*0.55}
    \node[draw, minimum width=1.1cm, minimum height=0.55cm] at (0,\y) {8};
    \node[draw, minimum width=1.1cm, minimum height=0.55cm] at (1.1,\y) {\rval};
    \node[draw, minimum width=2.2cm, minimum height=0.55cm] at (2.2,\y) {\pvalA};
    \node[draw, minimum width=2.2cm, minimum height=0.55cm] at (4.4,\y) {\pvalB};
    \node[draw, minimum width=2.2cm, minimum height=0.55cm] at (6.6,\y) {\pvalC};
    \node[draw, minimum width=2.2cm, minimum height=0.55cm] at (8.8,\y) {\pvalD};
    \node[draw, minimum width=2.2cm, minimum height=0.55cm] at (11.0,\y) {\pvalE};
    \node[draw, minimum width=2.2cm, minimum height=0.55cm] at (13.2,\y) {\pvalF};
  }

  % Highlight the cell for n=8, p=0.30, r=3
  \draw[poporange, very thick] (8.8,-2.3) rectangle (11.0,-2.85);
  \draw[poporange, very thick, ->] (11.0,-2.575) -- (12.5,-2.575) node[right, text width=3cm, align=left] {P(X\le 3) = 0.8059};
\end{tikzpicture}
\legende{Extract of a cumulative binomial table for $n=8$. The highlighted cell gives $P(X\le 3)$ when $p=0.30$.}
\end{popfigure}

\courssub{Reading Probabilities from the Table}

The table gives \textbf{cumulative probabilities} $P(X\le r)$.  From these we can obtain any probability we need.

\begin{methode}[Using Cumulative Binomial Tables]
\textbf{Step 1:} Identify $n$ and $p$.  If $p>0.5$, see the symmetry trick below.\\
\textbf{Step 2:} Locate the block for your $n$, then the column for your $p$.\\
\textbf{Step 3:} Read $P(X\le r)$ directly from the row corresponding to $r$.\\
\textbf{Step 4:} Convert your required probability to the form $P(X\le r)$ using the rules:
\begin{itemize}
  \item $P(X=r) = P(X\le r) - P(X\le r-1)$ (with $P(X\le -1)=0$).
  \item $P(X\ge r) = 1 - P(X\le r-1)$.
  \item $P(X>r) = 1 - P(X\le r)$.
  \item $P(a\le X\le b) = P(X\le b) - P(X\le a-1)$.
\end{itemize}
\end{methode}

\begin{attention}[Strict vs Non-Strict Inequalities — The \#1 GCE Trap]
Many candidates write $P(X<r)$ as $P(X\le r)$ instead of $P(X\le r-1)$.  Remember:
\[
P(X<r)=P(X\le r-1),\qquad P(X>r)=1-P(X\le r).
\]
Always convert strict inequalities to non-strict ones \emph{before} looking at the table.
\end{attention}

\courssub{Worked Examples: Reading the Table}

\begin{exoresolu}[Using the table for $X\sim B(8,0.3)$]
Using the table extract above, find:
\begin{enumerate}
  \item $P(X=3)$
  \item $P(X\ge 4)$
  \item $P(2\le X\le 5)$
\end{enumerate}
\tcblower
\textbf{Solution.}  From the table, for $n=8$, $p=0.30$:
\[
\begin{array}{c|c}
r & P(X\le r) \\ \hline
0 & 0.0576 \\
1 & 0.2553 \\
2 & 0.5518 \\
3 & 0.8059 \\
4 & 0.9420 \\
5 & 0.9887 \\
\end{array}
\]
\begin{enumerate}
  \item $P(X=3)=P(X\le3)-P(X\le2)=0.8059-0.5518=0.2541$.
  \item $P(X\ge4)=1-P(X\le3)=1-0.8059=0.1941$.
  \item $P(2\le X\le5)=P(X\le5)-P(X\le1)=0.9887-0.2553=0.7334$.
\end{enumerate}
\end{exoresolu}

\courssub{The Symmetry Trick for $p>0.5$}

Tables only go up to $p=0.5$.  If $X\sim B(n,p)$ with $p>0.5$, define $Y=n-X$, the number of \emph{failures}.  Then $Y\sim B(n,q)$ where $q=1-p<0.5$.  Any probability about $X$ can be rewritten in terms of $Y$ and read from the table.

\begin{propriete}[Symmetry of the Binomial Distribution]
If $X\sim B(n,p)$ and $Y=n-X$, then $Y\sim B(n,1-p)$.  Consequently:
\[
P(X\le r)=P(Y\ge n-r)=1-P(Y\le n-r-1).
\]
\end{propriete}

\begin{exemplebox}[Example: $p>0.5$]
Let $X\sim B(10,0.7)$.  Find $P(X\le 8)$ using tables.
\tcblower
Here $p=0.7>0.5$.  Let $Y=10-X\sim B(10,0.3)$.
\[
P(X\le 8)=P(10-Y\le 8)=P(Y\ge 2)=1-P(Y\le 1).
\]
Look up $n=10$, $p=0.30$, $r=1$ in the table to get $P(Y\le1)$, then subtract from 1.
\end{exemplebox}

\courssub{The Recurrence Formula and the Mode}

Instead of computing each $P(X=r)$ from the factorial formula, we can generate the whole distribution recursively.

\begin{propriete}[Recurrence Relation for Binomial Probabilities]
For $X\sim B(n,p)$ with $q=1-p$,
\[
\frac{P(X=r+1)}{P(X=r)} = \frac{(n-r)p}{(r+1)q},\qquad r=0,1,\dots,n-1.
\]
\textbf{Proof:}
\[
\frac{P(X=r+1)}{P(X=r)} = \frac{\binom{n}{r+1}p^{r+1}q^{n-r-1}}{\binom{n}{r}p^r q^{n-r}}
= \frac{\frac{n!}{(r+1)!(n-r-1)!}}{\frac{n!}{r!(n-r)!}}\cdot\frac{p}{q}
= \frac{r!(n-r)!}{(r+1)!(n-r-1)!}\cdot\frac{p}{q}
= \frac{n-r}{r+1}\cdot\frac{p}{q}.
\]
\end{propriete}

\begin{aretenir}[Generating the Distribution Recursively]
Start with $P(X=0)=q^n$.  Then for $r=0,1,\dots,n-1$:
\[
P(X=r+1)=P(X=r)\times\frac{(n-r)p}{(r+1)q}.
\]
This avoids large factorials and is ideal for programming or manual calculation when tables are unavailable.
\end{aretenir}

\courssub{Finding the Mode from the Recurrence Ratio}

The \textbf{mode} is the value of $r$ for which $P(X=r)$ is maximal.  The ratio $\frac{P(X=r+1)}{P(X=r)}$ tells us whether the probabilities are increasing or decreasing.

\begin{itemize}
  \item If the ratio $>1$, then $P(X=r+1)>P(X=r)$ — probabilities are increasing.
  \item If the ratio $<1$, then $P(X=r+1)<P(X=r)$ — probabilities are decreasing.
  \item If the ratio $=1$, then $P(X=r+1)=P(X=r)$ — two adjacent values are equally maximal.
\end{itemize}

The ratio exceeds 1 when
\[
\frac{(n-r)p}{(r+1)q} > 1 \;\Longleftrightarrow\; (n-r)p > (r+1)q \;\Longleftrightarrow\; r < (n+1)p - 1.
\]
Thus probabilities increase while $r < (n+1)p-1$ and decrease after $r > (n+1)p-1$.  The peak occurs at the integer part of $(n+1)p$.

\begin{propriete}[Mode of the Binomial Distribution]
Let $X\sim B(n,p)$ and let $m=(n+1)p$.
\begin{itemize}
  \item If $m$ is \textbf{not an integer}, the unique mode is $\lfloor m\rfloor$ (the integer part of $m$).
  \item If $m$ is \textbf{an integer}, there are \textbf{two modes}: $m$ and $m-1$ (both have the same maximal probability).
\end{itemize}
\end{propriete}

\begin{exemplebox}[Mode of $B(9,0.4)$]
Here $n=9$, $p=0.4$, so $m=(9+1)\times0.4=4$.  Since $m=4$ is an integer, the distribution has \textbf{two modes}: $r=4$ and $r=3$.
\tcblower
We can verify using the recurrence ratio:
\[
\frac{P(X=4)}{P(X=3)} = \frac{(9-3)\times0.4}{4\times0.6} = \frac{2.4}{2.4}=1.
\]
Hence $P(X=4)=P(X=3)$ and both are maximal.
\end{exemplebox}

\begin{popfigure}
\centering
\begin{tikzpicture}
\begin{axis}[
  ybar,
  bar width=6pt,
  width=12cm,
  height=7cm,
  xlabel={$r$},
  ylabel={$P(X=r)$},
  xtick={0,1,2,3,4,5,6,7,8,9},
  ymin=0,
  ymax=0.28,
  enlarge x limits=0.1,
  enlarge y limits={upper=0.1},
  nodes near coords,
  nodes near coords align={vertical},
  every node near coord/.style={font=\tiny, rotate=90, anchor=west, yshift=2pt},
  point meta=explicit symbolic
]
\addplot+[fill=popblue, draw=popdark] coordinates {
  (0,0.0101) [0.010]
  (1,0.0605) [0.060]
  (2,0.1612) [0.161]
  (3,0.2508) [0.251]
  (4,0.2508) [0.251]
  (5,0.1672) [0.167]
  (6,0.0743) [0.074]
  (7,0.0212) [0.021]
  (8,0.0035) [0.004]
  (9,0.0003) [0.000]
};
% Highlight modal bars
\addplot+[fill=poporange, draw=popdark] coordinates {
  (3,0.2508) [0.251]
  (4,0.2508) [0.251]
};
\end{axis}
\end{tikzpicture}
\legende{Probability mass function of $B(9,0.4)$. The two modal bars at $r=3$ and $r=4$ are highlighted in orange.}
\end{popfigure}

\courssub{Worked Example: Full Distribution and Mode}

\begin{exoresolu}[Constructing $B(6,0.3)$ recursively and finding its mode]
\begin{enumerate}
  \item Generate all probabilities $P(X=r)$ for $r=0,\dots,6$ using the recurrence formula.
  \item State the mode.
\end{enumerate}
\tcblower
\textbf{Solution.}
\begin{enumerate}
  \item $q=0.7$.  $P(X=0)=0.7^6=0.117649$.
  \begin{align*}
    P(X=1) &= 0.117649\times\frac{6\times0.3}{1\times0.7}=0.302526\\
    P(X=2) &= 0.302526\times\frac{5\times0.3}{2\times0.7}=0.324135\\
    P(X=3) &= 0.324135\times\frac{4\times0.3}{3\times0.7}=0.185220\\
    P(X=4) &= 0.185220\times\frac{3\times0.3}{4\times0.7}=0.059535\\
    P(X=5) &= 0.059535\times\frac{2\times0.3}{5\times0.7}=0.010206\\
    P(X=6) &= 0.010206\times\frac{1\times0.3}{6\times0.7}=0.000729
  \end{align*}
  (Sum $\approx 1.0000$, a good check.)
  \item $m=(6+1)\times0.3=2.1$.  Not an integer $\Rightarrow$ unique mode at $\lfloor2.1\rfloor=2$.
  Indeed $P(X=2)=0.3241$ is the largest probability.
\end{enumerate}
\end{exoresolu}

\courssub{Linking Probability Distributions to Frequency Distributions}

In Lesson 133 you will see that a theoretical probability distribution (like $B(n,p)$) can be fitted to an observed frequency distribution.  The expected frequency for value $r$ is $N\times P(X=r)$ where $N$ is the total number of observations.  The recurrence formula is extremely useful here: once you have $P(X=0)$, you can compute all expected frequencies recursively without recalculating binomial coefficients.

\begin{exemplebox}[Fitting a Binomial Distribution]
A die is rolled 120 times and the number of sixes in each set of 5 rolls is recorded.  If the die is fair, the number of sixes per set follows $B(5,1/6)$.  Use the recurrence formula to find the expected frequencies for 0,1,2,3,4,5 sixes.
\tcblower
$p=1/6$, $q=5/6$, $N=120$.
$P(X=0)=(5/6)^5\approx0.4019$; expected frequency $=120\times0.4019\approx48.2$.
Recurrence ratio $=\frac{(5-r)\times1/6}{(r+1)\times5/6}=\frac{5-r}{5(r+1)}$.
Compute successively:
$r=0\to1$: ratio $=5/5=1$ $\Rightarrow$ $P(X=1)=0.4019$ (freq 48.2)
$r=1\to2$: ratio $=4/10=0.4$ $\Rightarrow$ $P(X=2)=0.1608$ (freq 19.3)
$r=2\to3$: ratio $=3/15=0.2$ $\Rightarrow$ $P(X=3)=0.0322$ (freq 3.9)
$r=3\to4$: ratio $=2/20=0.1$ $\Rightarrow$ $P(X=4)=0.0032$ (freq 0.4)
$r=4\to5$: ratio $=1/25=0.04$ $\Rightarrow$ $P(X=5)=0.0001$ (freq 0.0)
\end{exemplebox}

\courssub{Summary of Key Points}

\begin{aretenir}[Cumulative Tables and Mode — Key Points]
\begin{itemize}
  \item Tables give $P(X\le r)$ for $p\le0.5$.  Convert all probabilities to this form.
  \item $P(X=r)=P(X\le r)-P(X\le r-1)$; $P(X\ge r)=1-P(X\le r-1)$; $P(a\le X\le b)=P(X\le b)-P(X\le a-1)$.
  \item For $p>0.5$, use $Y=n-X\sim B(n,1-p)$ and read from the $q$-column.
  \item Recurrence: $\dfrac{P(X=r+1)}{P(X=r)}=\dfrac{(n-r)p}{(r+1)q}$.  Start from $P(X=0)=q^n$.
  \item Mode = $\lfloor(n+1)p\rfloor$; if $(n+1)p$ is an integer, modes are $(n+1)p$ and $(n+1)p-1$.
\end{itemize}
\end{aretenir}

\begin{exercice}[Practice Exercise]
$X\sim B(12,0.35)$.
\begin{enumerate}
  \item Using tables, find $P(X\le 4)$, $P(X=5)$, $P(X\ge 7)$, $P(3\le X\le 8)$.
  \item If tables only go up to $p=0.5$, explain how you would find $P(X\le 8)$ for $Y\sim B(12,0.65)$.
  \item Find the mode of $X$ using the formula.  Verify with the recurrence ratio.
  \item Generate $P(X=0), P(X=1), P(X=2)$ recursively.
\end{enumerate}
\end{exercice}

\begin{corrige}[Exercise 1]
\begin{enumerate}
  \item From tables ($n=12$, $p=0.35$):
    $P(X\le4)=0.5833$ (approx),
    $P(X=5)=P(X\le5)-P(X\le4)=0.7745-0.5833=0.1912$,
    $P(X\ge7)=1-P(X\le6)=1-0.8912=0.1088$,
    $P(3\le X\le8)=P(X\le8)-P(X\le2)=0.9737-0.2528=0.7209$.
  \item For $Y\sim B(12,0.65)$, let $Z=12-Y\sim B(12,0.35)$.
    $P(Y\le8)=P(12-Z\le8)=P(Z\ge4)=1-P(Z\le3)$.  Read $P(Z\le3)$ from $p=0.35$ column.
  \item $m=(12+1)\times0.35=4.55$.  Not integer $\Rightarrow$ mode = $\lfloor4.55\rfloor=4$.
    Check ratio at $r=4$: $\frac{(12-4)\times0.35}{5\times0.65}=\frac{2.8}{3.25}<1$; at $r=3$: $\frac{9\times0.35}{4\times0.65}=\frac{3.15}{2.6}>1$.  So peak at $r=4$.
  \item $P(X=0)=0.65^{12}\approx0.00569$.
    $P(X=1)=0.00569\times\frac{12\times0.35}{1\times0.65}\approx0.0369$.
    $P(X=2)=0.0369\times\frac{11\times0.35}{2\times0.65}\approx0.1088$.
\end{enumerate}
\end{corrige}

\courssub{Examination Tips}

\begin{savaistu}[GCE Examiner's Advice]
\begin{itemize}
  \item \textbf{Always show the conversion step:} Write ``$P(X>3)=1-P(X\le3)$'' before giving the numerical answer.  This earns the method mark even if you misread the table.
  \item \textbf{For $p>0.5$:} Explicitly define $Y=n-X$ and state $Y\sim B(n,1-p)$.  Do not just say ``use symmetry'' — show the algebra.
  \item \textbf{Mode questions:} If asked to ``find the mode'', state the formula $m=(n+1)p$, compute it, and give the integer part.  If $m$ is integer, write ``two modes: $m$ and $m-1$''.
  \item \textbf{Recurrence formula:} If asked to ``use the recurrence formula'', you must show the ratio $\frac{P(X=r+1)}{P(X=r)}$ and at least one step of the calculation.
\end{itemize}
\end{savaistu}

This completes the section on cumulative binomial tables and the mode.  In the next section we will study the \textbf{discrete uniform distribution}, where every outcome is equally likely.

\coursec{The Discrete Uniform Distribution}

Having mastered the binomial distribution — its probabilities, cumulative tables, and mode — we now turn to a much simpler but equally important model: the \textbf{discrete uniform distribution}. While the binomial distribution counts successes in a fixed number of independent trials, the uniform distribution describes situations where \emph{every outcome is equally likely}. Think of rolling a fair die, drawing a raffle ticket from a well-mixed drum, or generating a random digit on a calculator. In each case, symmetry dictates that no outcome is favoured over another. This section builds the distribution from first principles, derives its mean and variance rigorously, and shows how to apply it to practical problems.

\courssub{Definition and Intuition}

\begin{definition}[Discrete Uniform Distribution on $\{1,2,\dots,n\}$]
A discrete random variable $X$ follows a \textbf{discrete uniform distribution} on the set $\{1,2,\dots,n\}$ (where $n\in\mathbb{N}$, $n\ge 1$) if its probability mass function is
\[
P(X=r)=\frac{1}{n}\qquad\text{for }r=1,2,\dots,n,
\]
and $P(X=r)=0$ otherwise. We write $X\sim U(1,n)$.
\end{definition}

The definition captures the idea of \emph{equally likely outcomes}. If an experiment has $n$ distinct outcomes and symmetry (or careful design) guarantees that none is more probable than another, then each outcome receives probability $1/n$. Note that the probabilities sum to $1$, as required:
\[
\sum_{r=1}^{n}P(X=r)=\sum_{r=1}^{n}\frac{1}{n}=n\cdot\frac{1}{n}=1.
\]

\begin{exemplebox}[Everyday Examples]
\begin{itemize}
\item \textbf{Fair die:} A standard six-sided die has faces $\{1,2,3,4,5,6\}$. If the die is fair, $X\sim U(1,6)$ and $P(X=r)=1/6$ for $r=1,\dots,6$.
\item \textbf{Raffle ticket:} A raffle at a school fate in Bamenda sells tickets numbered $1$ to $1000$. One ticket is drawn at random. The winning number $X\sim U(1,1000)$.
\item \textbf{Random digit:} A calculator's random-integer function returns an integer from $0$ to $9$ with equal probability. Then $X$ is uniform on $\{0,1,\dots,9\}$ — a shift of the standard definition (see the extension below).
\end{itemize}
\end{exemplebox}

\begin{popfigure}
\begin{tikzpicture}
\begin{axis}[
    width=0.9\linewidth,
    height=7cm,
    ybar,
    bar width=12pt,
    xlabel={Outcome $r$},
    ylabel={Probability $P(X=r)$},
    xtick={1,2,3,4,5,6},
    ytick={0,0.05,0.1,0.15},
    ymin=0,
    ymax=0.2,
    axis lines=left,
    enlarge x limits=0.15,
    enlarge y limits={upper=0.15},
    tick label style={font=\small},
    label style={font=\small},
]
\addplot[fill=popblue, draw=popdark] coordinates {
    (1,0.1666667)
    (2,0.1666667)
    (3,0.1666667)
    (4,0.1666667)
    (5,0.1666667)
    (6,0.1666667)
};
\draw[dashed, thick, poporange] (axis cs:3.5,0) -- (axis cs:3.5,0.18);
\node[poporange, anchor=south, font=\small] at (axis cs:3.5,0.182) {$E(X)=3.5$};
\end{axis}
\end{tikzpicture}
\legende{Probability mass function of $X\sim U(1,6)$ (fair die). All bars have equal height $1/6$; the mean $E(X)=3.5$ is marked by the orange dashed line.}
\end{popfigure}

\courssub{Expectation and Variance}

We derive the mean and variance using two classic summation formulas. These derivations are \textbf{examinable} — you must be able to reproduce them step by step in the GCE Advanced Level examination.

\begin{methode}[Key Summation Formulas to Memorise]
For any positive integer $n$:
\begin{enumerate}
\item $\displaystyle\sum_{r=1}^{n} r = \frac{n(n+1)}{2}$ \quad (sum of the first $n$ integers)
\item $\displaystyle\sum_{r=1}^{n} r^2 = \frac{n(n+1)(2n+1)}{6}$ \quad (sum of the first $n$ squares)
\end{enumerate}
These two formulas drive the derivations of $E(X)$ and $E(X^2)$ for the uniform distribution.
\end{methode}

\begin{propriete}[Expectation of $U(1,n)$ — proof examinable]
If $X\sim U(1,n)$, then
\[
E(X)=\frac{n+1}{2}.
\]
\end{propriete}

\textbf{Proof.} By the definition of expectation for a discrete random variable,
\[
E(X)=\sum_{r=1}^{n} r\cdot P(X=r)=\sum_{r=1}^{n} r\cdot\frac{1}{n}=\frac{1}{n}\sum_{r=1}^{n} r.
\]
Using the first summation formula,
\[
E(X)=\frac{1}{n}\cdot\frac{n(n+1)}{2}=\frac{n+1}{2}. \qquad\blacksquare
\]

\begin{propriete}[Second Moment and Variance of $U(1,n)$ — proof examinable]
If $X\sim U(1,n)$, then
\[
E(X^2)=\frac{(n+1)(2n+1)}{6},\qquad
\operatorname{Var}(X)=\frac{n^2-1}{12}.
\]
\end{propriete}

\textbf{Proof.} First compute $E(X^2)$:
\[
E(X^2)=\sum_{r=1}^{n} r^2\cdot P(X=r)=\frac{1}{n}\sum_{r=1}^{n} r^2
=\frac{1}{n}\cdot\frac{n(n+1)(2n+1)}{6}=\frac{(n+1)(2n+1)}{6}.
\]
Now use the computational formula $\operatorname{Var}(X)=E(X^2)-[E(X)]^2$. Put both terms over the common denominator $12$:
\[
E(X^2)=\frac{(n+1)(2n+1)}{6}=\frac{2(n+1)(2n+1)}{12}=\frac{2(2n^2+3n+1)}{12}=\frac{4n^2+6n+2}{12},
\]
\[
[E(X)]^2=\left(\frac{n+1}{2}\right)^2=\frac{(n+1)^2}{4}=\frac{3(n+1)^2}{12}=\frac{3n^2+6n+3}{12}.
\]
Subtracting,
\[
\operatorname{Var}(X)=\frac{(4n^2+6n+2)-(3n^2+6n+3)}{12}=\frac{n^2-1}{12}. \qquad\blacksquare
\]

\begin{attention}[Algebraic Trap in the Variance Derivation]
When subtracting $\frac{(n+1)(2n+1)}{6}-\frac{(n+1)^2}{4}$, do \textbf{not} expand carelessly. The safest route is the common denominator $12$ as shown above. A frequent error is to write $\frac{2n+1}{6}-\frac{n+1}{4}=\frac{4(2n+1)-3(n+1)}{12}$ and then mis-expand the numerator. Expand fully: $4(2n+1)=8n+4$ and $3(n+1)=3n+3$, giving $\frac{5n+1}{12}$ — which is \emph{wrong} because the factor $(n+1)$ was pulled out incorrectly at the start. Work with the full numerators $4n^2+6n+2$ and $3n^2+6n+3$ instead; the $6n$ terms cancel and the result $\frac{n^2-1}{12}$ follows cleanly.
\end{attention}

\begin{aretenir}[Summary for $X\sim U(1,n)$]
\[
P(X=r)=\frac{1}{n}\quad(r=1,\dots,n),\qquad E(X)=\frac{n+1}{2},\qquad \operatorname{Var}(X)=\frac{n^2-1}{12}.
\]
\end{aretenir}

\courssub{Extension: Uniform Distribution on $\{a,a+1,\dots,b\}$}

In many practical situations the equally likely outcomes do not start at $1$. If $X$ takes the integer values $a,a+1,\dots,b$ with equal probability, where $a,b\in\mathbb{Z}$ and $a\le b$, then the number of outcomes is $n=b-a+1$. We can write $X=(a-1)+Y$ where $Y\sim U(1,n)$. Using linearity of expectation, $E(X)=(a-1)+E(Y)$, and the fact that adding a constant does not change variance, $\operatorname{Var}(X)=\operatorname{Var}(Y)$. This gives the following general result.

\begin{propriete}[Uniform Distribution on $\{a,a+1,\dots,b\}$]
If $X$ is uniformly distributed on $\{a,a+1,\dots,b\}$ (with $n=b-a+1$ outcomes), then
\[
P(X=r)=\frac{1}{n}\quad(r=a,\dots,b),\qquad
E(X)=\frac{a+b}{2},\qquad
\operatorname{Var}(X)=\frac{n^2-1}{12}=\frac{(b-a+1)^2-1}{12}.
\]
\end{propriete}

\textbf{Justification of the mean.} $E(X)=(a-1)+\dfrac{n+1}{2}=\dfrac{2a-2+n+1}{2}=\dfrac{2a+(b-a+1)-1}{2}=\dfrac{a+b}{2}$. $\blacksquare$

\begin{exemplebox}[Fair Die Revisited]
For a fair die, $a=1$, $b=6$, $n=6$.
\[
E(X)=\frac{1+6}{2}=3.5,\qquad
\operatorname{Var}(X)=\frac{6^2-1}{12}=\frac{35}{12}\approx 2.9167.
\]
The standard deviation is $\sigma=\sqrt{35/12}\approx 1.708$.
\end{exemplebox}

\courssub{Worked Examples}

\begin{exoresolu}[Example 1: Expectation and Variance of a Fair Die Score]
A fair six-sided die is rolled once. Let $X$ be the score obtained.
\begin{enumerate}
\item Write down the distribution of $X$.
\item Calculate $E(X)$ and $\operatorname{Var}(X)$.
\item If the die is rolled twice and the scores are added, find the mean and variance of the total score $T$.
\end{enumerate}
\tcblower
\textbf{Solution.}
\begin{enumerate}
\item $X\sim U(1,6)$, so $P(X=r)=1/6$ for $r=1,\dots,6$.
\item Using the formulas with $n=6$:
\[
E(X)=\frac{6+1}{2}=3.5,\qquad
\operatorname{Var}(X)=\frac{6^2-1}{12}=\frac{35}{12}.
\]
\item Let $X_1,X_2$ be the scores of the two independent rolls. Then $T=X_1+X_2$. By linearity of expectation,
\[
E(T)=E(X_1)+E(X_2)=3.5+3.5=7.
\]
Since the rolls are independent, variances add:
\[
\operatorname{Var}(T)=\operatorname{Var}(X_1)+\operatorname{Var}(X_2)=\frac{35}{12}+\frac{35}{12}=\frac{35}{6}\approx 5.833.
\]
\end{enumerate}
\end{exoresolu}

\begin{exoresolu}[Example 2: Checking Fairness of a Game]
A game at a school fair uses a spinner with sectors numbered $2,3,4,5,6,7$. A player pays $200$ FCFA to spin once and wins $50$ FCFA times the number shown. Determine whether the game is fair (i.e., expected net gain is zero) from the player's perspective.
\tcblower
\textbf{Solution.}
Let $X$ be the number shown. $X$ is uniform on $\{2,3,4,5,6,7\}$, so $a=2$, $b=7$, $n=6$.
\[
E(X)=\frac{a+b}{2}=\frac{2+7}{2}=4.5.
\]
The winnings are $W=50X$ (in FCFA). The cost to play is $C=200$ FCFA.
Net gain $G=W-C=50X-200$.
\[
E(G)=50E(X)-200=50\times4.5-200=225-200=25\text{ FCFA}.
\]
Since $E(G)=25>0$, the game is \textbf{not fair} — it favours the player. (A fair game would have $E(G)=0$, which would require a cost of $50\times4.5=225$ FCFA per spin.)
\end{exoresolu}

\courssub{Common Mistakes and Examination Tips}

\begin{attention}[Discrete vs. Continuous Uniform]
Do \textbf{not} confuse the discrete uniform distribution (this section) with the continuous uniform distribution (met later in the course).
\begin{itemize}
\item \textbf{Discrete:} $X$ takes integer values $1,\dots,n$ with $P(X=r)=1/n$. Mean $=(n+1)/2$, Variance $=(n^2-1)/12$.
\item \textbf{Continuous:} $X$ takes any real value in $[a,b]$ with density $f(x)=1/(b-a)$. Mean $=(a+b)/2$, Variance $=(b-a)^2/12$.
\end{itemize}
The formulas look similar but the domains and the nature of probability (point masses vs. density) are completely different. In the GCE exam, the context (integer outcomes vs. continuous measurements) will tell you which one to use.
\end{attention}

\begin{attention}[Summation Limits and the Factor $1/n$]
When deriving $E(X)$ and $E(X^2)$, always write the summation from $r=1$ to $n$ explicitly, and do not forget the factor $1/n$ coming from $P(X=r)=1/n$. A common error is to quote $E(X)=\sum r = n(n+1)/2$ without dividing by $n$. Remember: $E(X)=\frac{1}{n}\sum_{r=1}^{n} r$.
\end{attention}

\begin{attention}[Variance Depends Only on the Number of Outcomes]
The variance $\frac{n^2-1}{12}$ is for the standard uniform on $\{1,\dots,n\}$. If the distribution is on $\{a,\dots,b\}$, the variance is \textbf{still} $\frac{n^2-1}{12}$ where $n=b-a+1$ — it does \textbf{not} depend on $a$ or $b$ individually, only on the number of outcomes. Shifting the distribution does not change its spread. For example, $U(1,6)$ and $U(2,7)$ both have variance $35/12$.
\end{attention}

\courssub{Link with Frequency Distributions (Preview)}

In the lesson on the relationship between probability distributions and frequency distributions, we will see how a theoretical probability distribution like $U(1,n)$ can be compared with an observed frequency distribution. For instance, if we roll a die $600$ times and record the frequency of each face, we expect each face to appear about $600\times\frac{1}{6}=100$ times. These \emph{expected frequencies} $E_r=N\cdot P(X=r)$ can then be compared with the observed frequencies to judge whether the die is truly fair — a comparison formalised by the $\chi^2$ goodness-of-fit test met in the statistics course.

\begin{exercice}[Practice Exercises]
\begin{enumerate}
\item A random number generator produces integers uniformly from $10$ to $20$ inclusive. Find the mean and variance of the generated number.
\item A bag contains tickets numbered $1$ to $50$. One ticket is drawn at random. Let $X$ be the number on the ticket. Find $P(X>40)$, $E(X)$ and $\operatorname{Var}(X)$.
\item A game uses a spinner with sectors $1,2,3,4,5$. It costs $100$ FCFA to play; you win $30$ FCFA times the number shown. Is the game fair? If not, who does it favour and by how much on average?
\item Prove that for $X\sim U(1,n)$, $E(X^2)=\dfrac{(n+1)(2n+1)}{6}$, and deduce that $\operatorname{Var}(X)=\dfrac{n^2-1}{12}$.
\item A fair die is rolled $360$ times. Write down the expected frequency of each face. What is the expected total of all $360$ scores?
\end{enumerate}
\end{exercice}

\begin{corrige}[Exercise 1]
$X$ is uniform on $\{10,11,\dots,20\}$. Here $a=10$, $b=20$, $n=20-10+1=11$.
$E(X)=\dfrac{10+20}{2}=15$.
$\operatorname{Var}(X)=\dfrac{11^2-1}{12}=\dfrac{120}{12}=10$.
\end{corrige}

\begin{corrige}[Exercise 2]
$X\sim U(1,50)$.
$P(X>40)=P(X=41,\dots,50)=10\times\dfrac{1}{50}=0.2$.
$E(X)=\dfrac{50+1}{2}=25.5$.
$\operatorname{Var}(X)=\dfrac{50^2-1}{12}=\dfrac{2499}{12}=208.25$.
\end{corrige}

\begin{corrige}[Exercise 3]
$X\sim U(1,5)$, so $E(X)=\dfrac{5+1}{2}=3$.
Winnings $W=30X$, cost $C=100$ FCFA.
Net gain $G=30X-100$.
$E(G)=30E(X)-100=30\times3-100=-10$ FCFA.
The game is not fair; it favours the organiser by $10$ FCFA per play on average.
\end{corrige}

\begin{corrige}[Exercise 4]
$E(X^2)=\displaystyle\sum_{r=1}^n r^2 P(X=r)=\frac{1}{n}\sum_{r=1}^n r^2=\frac{1}{n}\cdot\frac{n(n+1)(2n+1)}{6}=\frac{(n+1)(2n+1)}{6}$.
Then
\[
\operatorname{Var}(X)=E(X^2)-[E(X)]^2=\frac{4n^2+6n+2}{12}-\frac{3n^2+6n+3}{12}=\frac{n^2-1}{12}.
\]
\end{corrige}

\begin{corrige}[Exercise 5]
Expected frequency of each face $=360\times\dfrac{1}{6}=60$.
Expected score per roll $=3.5$, so the expected total of all $360$ scores is $360\times3.5=1260$.
\end{corrige}

\coursec{The Geometric Distribution}

In the previous section we studied the discrete uniform distribution, where every outcome in a finite set is equally likely. We now turn to a distribution that models a very different kind of random experiment: we repeat independent trials, each with the same probability of success, until the \emph{first success} occurs. The random variable is the \emph{number of trials needed} to obtain that first success. This is the \cle{geometric distribution}, a cornerstone for modelling waiting times in discrete settings.

\courssub{Definition and Recognition of Geometric Situations}

\begin{experience}[A Motivating Situation]
A trader in Douala sells phone credit. Experience shows that each passer-by she approaches buys credit with probability $0.2$, independently of the others. She wonders: \emph{how many people must she approach before she makes her first sale?} It could be the very first person (probability $0.2$), or the second (a failure then a success: $0.8 \times 0.2$), or the third ($0.8 \times 0.8 \times 0.2$), and so on. The number of approaches needed is a random variable whose probabilities form a geometric sequence: this is exactly the situation modelled by the geometric distribution.
\end{experience}

\begin{definition}[Geometric Distribution]
Let an experiment consist of independent Bernoulli trials, each with probability of success $p$ ($0<p<1$) and probability of failure $q=1-p$. Let $X$ be the random variable representing the \textbf{number of trials up to and including the first success}. Then $X$ follows a \cle{geometric distribution} with parameter $p$, denoted $X \sim \mathrm{Geo}(p)$. Its probability mass function is
\[
P(X = r) = q^{\,r-1}\,p \qquad (r = 1,2,3,\dots)
\]
\end{definition}

The formula is immediate: the event $\{X=r\}$ means that the first $r-1$ trials are failures and the $r$-th trial is a success. By independence,
\[
P(X=r) = \underbrace{q \times q \times \cdots \times q}_{r-1 \text{ failures}} \times p = q^{\,r-1}p.
\]

\begin{aretenir}[Recognising a Geometric Setting]
A random variable $X$ has a geometric distribution if the situation satisfies:
\begin{itemize}
\item Trials are independent.
\item Each trial has only two outcomes: \emph{success} (probability $p$) and \emph{failure} (probability $q=1-p$).
\item The probability $p$ remains constant from trial to trial.
\item We count the number of trials \textbf{until the first success occurs} (including that success).
\end{itemize}
Typical examples: the number of times you roll a die until you get a six; the number of job interviews until the first offer; the number of items inspected until the first defective one is found.
\end{aretenir}

\begin{attention}[Two Common Definitions — Check the Convention!]
Some textbooks (and some software packages) define the geometric variable as the \emph{number of failures before the first success}, taking values $0,1,2,\dots$ with $P(Y=k)=q^k p$. The GCE Advanced Level convention \textbf{uses the definition above}: $X$ is the \emph{number of trials including the first success}, taking values $1,2,3,\dots$ with $P(X=r)=q^{r-1}p$. Using the wrong convention shifts the mean by $1$ and changes the variance formula, so always state clearly which variable you are counting.
\end{attention}

\begin{exemplebox}[Identifying a Geometric Variable]
A factory produces light bulbs; $5\%$ are defective. Bulbs are tested one by one until the first defective bulb is found. Let $X$ be the number of bulbs tested.
\begin{itemize}
\item Trials: testing a bulb. Independent? Yes (large batch).
\item Success: finding a defective bulb, $p=0.05$, $q=0.95$.
\item $X$ counts trials up to and including the first defective, so $X\sim\mathrm{Geo}(0.05)$.
\item $P(X=3) = (0.95)^2 \times 0.05 = 0.045125$.
\end{itemize}
\end{exemplebox}

\courssub{Cumulative Probabilities and the Tail Formula}

Often we need probabilities such as ``more than $r$ trials are needed'' or ``at most $r$ trials''. Summing the probability mass function term by term is possible but tedious. A simple closed form exists for the upper tail.

\begin{propriete}[Tail Probability $P(X>r)$]
For $X\sim\mathrm{Geo}(p)$,
\[
P(X > r) = q^{\,r} \qquad (r = 0,1,2,\dots)
\]
\emph{Proof:} The event $\{X>r\}$ means that the first $r$ trials are all failures. By independence, its probability is $q \times q \times \cdots \times q = q^r$. \qed
\end{propriete}

From this we immediately obtain:
\[
P(X \le r) = 1 - q^{\,r}, \qquad
P(X \ge r) = q^{\,r-1}, \qquad
P(X < r) = 1 - q^{\,r-1}.
\]

\begin{attention}[Strict or Non-Strict Inequalities]
A very common examination error is to confuse $P(X>r)=q^r$ with $P(X\ge r)=q^{r-1}$. Remember: $P(X>r)$ requires $r$ failures, while $P(X\ge r)$ requires only $r-1$ failures (the $r$-th trial is free). Always translate the sentence into an event on the trials before applying the formula.
\end{attention}

\begin{exemplebox}[Using the Tail Formula]
With $p=0.05$ ($q=0.95$) as above, find the probability that more than $10$ bulbs are tested before the first defective appears.
\[
P(X > 10) = (0.95)^{10} \approx 0.5987.
\]
Hence $P(X \le 10) = 1 - 0.5987 = 0.4013$.
\end{exemplebox}

\courssub{Expectation and Variance}

The mean and variance of the geometric distribution have elegant closed forms. Their derivations use standard power-series manipulations and are examinable at this level.

\begin{propriete}[Mean and Variance of $\mathrm{Geo}(p)$]
If $X \sim \mathrm{Geo}(p)$ with $q=1-p$, then
\[
E(X) = \frac{1}{p}, \qquad \mathrm{Var}(X) = \frac{q}{p^2}.
\]
\emph{Proof (examinable).} We derive both results from the geometric series and its derivatives.

\textbf{Step 1: the mean.}
\[
E(X) = \sum_{r=1}^{\infty} r\,q^{\,r-1}p = p \sum_{r=1}^{\infty} r\,q^{\,r-1}.
\]
Recall the geometric series, valid for $|q|<1$:
\[
\sum_{r=0}^{\infty} q^r = \frac{1}{1-q} = \frac{1}{p}.
\]
Differentiate both sides with respect to $q$:
\[
\sum_{r=1}^{\infty} r\,q^{\,r-1} = \frac{1}{(1-q)^2} = \frac{1}{p^2}.
\]
Hence
\[
E(X) = p \cdot \frac{1}{p^2} = \frac{1}{p}.
\]

\textbf{Step 2: the variance.} We first compute $E[X(X-1)]$, which avoids subtracting large nearly equal numbers:
\[
E[X(X-1)] = \sum_{r=2}^{\infty} r(r-1)\,q^{\,r-1}p = pq \sum_{r=2}^{\infty} r(r-1)\,q^{\,r-2}.
\]
Differentiate the identity $\sum_{r=1}^{\infty} r\,q^{\,r-1} = \dfrac{1}{(1-q)^2}$ once more with respect to $q$:
\[
\sum_{r=2}^{\infty} r(r-1)\,q^{\,r-2} = \frac{2}{(1-q)^3} = \frac{2}{p^3}.
\]
Therefore
\[
E[X(X-1)] = pq \cdot \frac{2}{p^3} = \frac{2q}{p^2}.
\]
Now $E(X^2) = E[X(X-1)] + E(X) = \dfrac{2q}{p^2} + \dfrac{1}{p} = \dfrac{2q+p}{p^2} = \dfrac{1+q}{p^2}$, since $2q+p = q+(q+p)=1+q$. Finally,
\[
\mathrm{Var}(X) = E(X^2) - [E(X)]^2 = \frac{1+q}{p^2} - \frac{1}{p^2} = \frac{q}{p^2}.
\]
\qed
\end{propriete}

\begin{exemplebox}[Mean and Variance in Context]
A basketball player has a free-throw success rate of $p=0.7$. Let $X$ be the number of shots until the first basket.
\begin{itemize}
\item $E(X) = 1/0.7 \approx 1.43$ shots.
\item $\mathrm{Var}(X) = 0.3/(0.7)^2 = 0.3/0.49 \approx 0.612$.
\item Standard deviation $\sigma \approx 0.78$ shots.
\end{itemize}
On average the player needs about $1.4$ shots to score; the small standard deviation reflects that success usually comes quickly when $p$ is large.
\end{exemplebox}

\courssub{The Mode of the Geometric Distribution}

\begin{propriete}[Mode of $\mathrm{Geo}(p)$]
For $X\sim\mathrm{Geo}(p)$ with $0<p<1$, the mode is always $1$.
\emph{Proof:} $P(X=r) = q^{\,r-1}p$. Since $0<q<1$, the sequence is strictly decreasing:
\[
\frac{P(X=r+1)}{P(X=r)} = q < 1 \quad\text{for all } r \ge 1,
\]
so $P(X=1)=p$ is the largest probability. \qed
\end{propriete}

This matches intuition: the single most likely scenario is that the very first trial is a success, whatever the value of $p$.

\courssub{The Memoryless Property}

The geometric distribution is the \emph{only} discrete distribution with the \cle{memoryless property}: past failures do not affect the distribution of the remaining waiting time.

\begin{propriete}[Memoryless Property]
For $X \sim \mathrm{Geo}(p)$ and any non-negative integers $m,n$,
\[
P(X > m+n \mid X > m) = P(X > n).
\]
\emph{Proof:} Since $\{X>m+n\} \subset \{X>m\}$,
\[
P(X > m+n \mid X > m) = \frac{P(X > m+n)}{P(X > m)} = \frac{q^{\,m+n}}{q^{\,m}} = q^{\,n} = P(X > n).
\]
\qed
\end{propriete}

\begin{savaistu}[Interpretation]
If you have already observed $m$ failures, the conditional distribution of the \emph{additional} number of trials needed for the first success is again $\mathrm{Geo}(p)$. The process ``restarts'' with no memory of the past. This is why the geometric distribution models systems that do not age, such as repeated independent attempts or certain queueing situations.
\end{savaistu}

\begin{popfigure}
\begin{tikzpicture}[scale=0.95, every node/.style={font=\small}]
\draw[->] (-0.5,0) -- (9.5,0) node[right] {trials};
\foreach \x in {1,...,8} {
  \draw (\x,0.08) -- (\x,-0.08) node[below] {\x};
}
\foreach \x in {1,2,3} {
  \fill[popblue] (\x,0.35) circle (2.2pt);
  \node[popblue, above] at (\x,0.35) {F};
}
\foreach \x in {4,5,6} {
  \fill[poporange] (\x,0.35) circle (2.2pt);
  \node[poporange, above] at (\x,0.35) {?};
}
\draw[popblue, thick] (1,-0.55) -- (3,-0.55);
\node[popblue, below] at (2,-0.55) {$m=3$ failures observed};
\draw[poporange, thick] (4,-0.55) -- (6,-0.55);
\node[poporange, below] at (5,-0.55) {$n=3$ further trials};
\node[align=center, text width=6.5cm] at (5,1.6) {Given $X>3$, the chance of needing more than $3$ further trials is $P(X>3)=q^3$, exactly as if we started again.};
\end{tikzpicture}
\legende{The memoryless property: after $m$ observed failures, the remaining waiting time has the same geometric distribution as at the start.}
\end{popfigure}

\courssub{Visualising the Geometric Distribution}

The probability mass function decreases geometrically: each bar is $q$ times the previous one. The following bar chart shows $X \sim \mathrm{Geo}(0.25)$ for $r=1,\dots,12$.

\begin{popfigure}
\begin{tikzpicture}
\begin{axis}[
    width=0.9\linewidth,
    height=6cm,
    ybar,
    bar width=6pt,
    xlabel={$r$ (trial of first success)},
    ylabel={$P(X=r)$},
    xtick={1,2,3,4,5,6,7,8,9,10,11,12},
    ymin=0,
    ymax=0.28,
    tick label style={font=\small},
    label style={font=\small},
    axis lines=left,
    enlarge x limits=0.05,
]
\addplot[fill=popblue, draw=popdark] coordinates {
    (1,0.25)
    (2,0.1875)
    (3,0.140625)
    (4,0.10546875)
    (5,0.0791015625)
    (6,0.0593261719)
    (7,0.0444946289)
    (8,0.0333709717)
    (9,0.0250282288)
    (10,0.0187711716)
    (11,0.0140783787)
    (12,0.0105587840)
};
\end{axis}
\end{tikzpicture}
\legende{Bar chart of $X\sim\mathrm{Geo}(0.25)$ for $r=1$ to $12$. The probabilities decay by a factor of $q=0.75$ at each step, and the mode is $1$.}
\end{popfigure}

\courssub{Applied Problems and Worked Examples}

We now solve typical examination questions that combine recognition, tail probabilities, expectation, variance, and the memoryless property.

\begin{exoresolu}[Finding $p$ from a Given Probability]
The random variable $X \sim \mathrm{Geo}(p)$. It is given that $P(X > 4) = 0.2401$. Find the value of $p$, and hence calculate $E(X)$ and $\mathrm{Var}(X)$.
\tcblower
\textbf{Solution:}
\textbf{Step 1:} Use the tail formula $P(X > r) = q^r$ with $r=4$:
\[
q^4 = 0.2401.
\]
\textbf{Step 2:} Take the fourth root. Since $0.7^2 = 0.49$ and $0.49^2 = 0.2401$, we get
\[
q = 0.7, \qquad p = 1 - q = 0.3.
\]
\textbf{Step 3:} Compute the mean and variance:
\[
E(X) = \frac{1}{p} = \frac{1}{0.3} = \frac{10}{3} \approx 3.33,
\qquad
\mathrm{Var}(X) = \frac{q}{p^2} = \frac{0.7}{0.09} = \frac{70}{9} \approx 7.78.
\]
\end{exoresolu}

\begin{exoresolu}[Memoryless Property in Action]
A machine produces components; $2\%$ are defective. Components are tested one by one until the first defective is found. Given that the first $20$ components tested are all non-defective, find the probability that more than $15$ further components need to be tested.
\tcblower
\textbf{Solution:}
Let $X \sim \mathrm{Geo}(0.02)$, so $q=0.98$. The condition ``the first 20 are non-defective'' means $X > 20$. We require
\[
P(X > 20+15 \mid X > 20) = P(X > 35 \mid X > 20).
\]
By the memoryless property this equals
\[
P(X > 15) = q^{15} = (0.98)^{15} \approx 0.7386.
\]
So the probability is approximately $0.739$: the twenty good components already tested give no information about the future.
\end{exoresolu}

\begin{exoresolu}[Expected Cost of Testing]
In a production line, each item tested costs $500$ CFA francs. The probability that an item is defective is $0.04$. Testing stops at the first defective item. Find the expected total cost and the standard deviation of the cost.
\tcblower
\textbf{Solution:}
Let $X \sim \mathrm{Geo}(0.04)$ be the number of items tested; the cost is $C = 500X$ CFA.
\[
E(X) = \frac{1}{0.04} = 25, \qquad \mathrm{Var}(X) = \frac{0.96}{(0.04)^2} = \frac{0.96}{0.0016} = 600.
\]
Using the linearity of expectation and $\mathrm{Var}(aX)=a^2\mathrm{Var}(X)$:
\[
E(C) = 500 \times 25 = 12\,500 \text{ CFA},
\]
\[
\mathrm{Var}(C) = 500^2 \times 600 = 150\,000\,000 \text{ CFA}^2,
\qquad
\sigma_C = \sqrt{150\,000\,000} \approx 12\,247 \text{ CFA}.
\]
\end{exoresolu}

\begin{exoresolu}[A Two-Stage Problem]
A fair coin is tossed repeatedly. Find the probability that the first head appears on an even-numbered toss.
\tcblower
\textbf{Solution:}
Let $X \sim \mathrm{Geo}(0.5)$ be the toss of the first head, with $p=q=\tfrac12$. We need
\[
P(X \text{ even}) = P(X=2) + P(X=4) + P(X=6) + \cdots = qp + q^3p + q^5p + \cdots
\]
This is a geometric series with first term $qp = \tfrac14$ and common ratio $q^2 = \tfrac14$:
\[
P(X \text{ even}) = \frac{qp}{1-q^2} = \frac{\tfrac14}{1-\tfrac14} = \frac{1}{3}.
\]
\emph{Check:} by symmetry of the argument, $P(X \text{ odd}) = \dfrac{p}{1-q^2} = \dfrac{2}{3}$, and indeed $\tfrac13+\tfrac23=1$.
\end{exoresolu}

\begin{aretenir}[Key Formulas for $\mathrm{Geo}(p)$]
\begin{itemize}
\item PMF: $P(X=r)=q^{\,r-1}p$, $r=1,2,3,\dots$
\item Tail: $P(X>r)=q^{\,r}$, \quad $P(X\le r)=1-q^{\,r}$, \quad $P(X\ge r)=q^{\,r-1}$.
\item Mean: $E(X)=\dfrac{1}{p}$.
\item Variance: $\mathrm{Var}(X)=\dfrac{q}{p^2}$.
\item Mode: $1$ (always).
\item Memoryless: $P(X>m+n\mid X>m)=P(X>n)=q^{\,n}$.
\end{itemize}
\end{aretenir}

\begin{methode}[Solving Geometric Distribution Problems]
\begin{enumerate}
\item \textbf{Identify} the success event and its probability $p$. Confirm independence and constant $p$.
\item \textbf{Define} $X$ = number of trials up to and including the first success. State $X\sim\mathrm{Geo}(p)$.
\item \textbf{Use} the tail formula $P(X>r)=q^r$ for ``more than $r$ trials''; avoid summing long series, and watch strict versus non-strict inequalities.
\item \textbf{Apply} $E(X)=1/p$ and $\mathrm{Var}(X)=q/p^2$ when expectations or variances are needed.
\item \textbf{Check} whether the memoryless property applies when the problem conditions on past failures.
\item \textbf{Interpret} the result in the context of the problem (expected cost, number of attempts, etc.).
\end{enumerate}
\end{methode}

\begin{exercice}[Practice Problems]
\begin{enumerate}
\item A fair die is rolled until a six appears. Find (a) the probability that exactly 4 rolls are needed, (b) the probability that more than 5 rolls are needed, (c) the expected number of rolls and the variance.
\item In a certain region, the probability that a randomly selected household has a solar panel is $0.12$. Households are visited until one with a solar panel is found. Find the probability that at most 3 households are visited. Given that the first 5 households visited have no solar panel, find the probability that more than 4 further households must be visited.
\item A random variable $X\sim\mathrm{Geo}(p)$ has variance $6$. Find $p$ and $E(X)$.
\item Each trial costs $2000$ CFA francs. If $X\sim\mathrm{Geo}(0.2)$ is the number of trials until the first success, find the expected cost and the standard deviation of the total cost.
\item Show that for $X\sim\mathrm{Geo}(p)$, $P(X \text{ is even}) = \dfrac{q}{1+q}$, and deduce the value of this probability when $p = 0.4$.
\end{enumerate}
\end{exercice}

\begin{corrige}[Exercise 1]
$p=1/6$, $q=5/6$.
(a) $P(X=4)=q^3p = (5/6)^3(1/6)=125/1296\approx0.0965$.
(b) $P(X>5)=q^5=(5/6)^5=3125/7776\approx0.4019$.
(c) $E(X)=1/p=6$; $\mathrm{Var}(X)=q/p^2=(5/6)/(1/36)=30$.
\end{corrige}

\begin{corrige}[Exercise 2]
$p=0.12$, $q=0.88$.
$P(X\le3)=1-q^3=1-(0.88)^3=1-0.681472=0.318528\approx0.319$.
Given $X>5$, the memoryless property gives
\[
P(X>5+4\mid X>5)=P(X>4)=q^4=(0.88)^4\approx0.5997.
\]
\end{corrige}

\begin{corrige}[Exercise 3]
$\mathrm{Var}(X)=q/p^2=6$ with $q=1-p$ gives $1-p=6p^2$, i.e. $6p^2+p-1=0$.
\[
p = \frac{-1+\sqrt{1+24}}{12} = \frac{-1+5}{12} = \frac{1}{3}
\]
(taking the positive root, since $0<p<1$). Hence $E(X)=1/p=3$.
\end{corrige}

\begin{corrige}[Exercise 4]
$X\sim\mathrm{Geo}(0.2)$: $E(X)=1/0.2=5$, $\mathrm{Var}(X)=0.8/(0.2)^2=0.8/0.04=20$.
Cost $C=2000X$: $E(C)=2000\times5=10\,000$ CFA; $\mathrm{Var}(C)=2000^2\times20=80\,000\,000$ CFA$^2$; $\sigma_C=\sqrt{80\,000\,000}\approx8944$ CFA.
\end{corrige}

\begin{corrige}[Exercise 5]
$P(X \text{ even}) = qp + q^3p + q^5p + \cdots = \dfrac{qp}{1-q^2} = \dfrac{qp}{(1-q)(1+q)} = \dfrac{q}{1+q}$, since $p=1-q$.
For $p=0.4$, $q=0.6$: $P(X \text{ even}) = \dfrac{0.6}{1.6} = \dfrac{3}{8} = 0.375$.
\end{corrige}

This completes our study of the geometric distribution. In the next section we will explore the Poisson distribution, which models the number of events occurring in a fixed interval of time or space, and its intimate connection with the binomial distribution.

\coursec{The Poisson Distribution and Its Tables}

In the previous section we studied the \cle{geometric distribution}, which counts the number of \emph{trials} needed to obtain the first success. In many real situations, however, there is no fixed sequence of trials at all: phone calls arrive at a call centre in Douala at unpredictable instants, vehicles pass a checkpoint on the Yaound\'e--Bafoussam road at random times, and typing errors appear scattered through the pages of a manuscript. What we can count is the \emph{number of events occurring in a fixed interval of time or space}. The distribution that models exactly this situation is the \cle{Poisson distribution}, named after the French mathematician Sim\'eon Denis Poisson. It is one of the most frequently examined distributions at the GCE Advanced Level.

\courssub{Definition and Conditions for a Poisson Process}

\begin{experience}[Counting random events]
Suppose a small clinic in Bamenda receives, on average, $3$ emergency patients per hour. If we observe the arrivals during one hour, we might record $0, 1, 2, 3, \dots$ patients. The arrivals are not scheduled: they occur \emph{at random}, \emph{independently} of one another, but the \emph{average rate} of $3$ per hour stays roughly constant throughout the day. Any count of events satisfying these three features is a candidate for a Poisson model.
\end{experience}

\begin{definition}[Poisson distribution]
A discrete random variable $X$ follows a \cle{Poisson distribution} with parameter $\lambda > 0$, written $X \sim \mathrm{Po}(\lambda)$, if
\[ P(X = r) = e^{-\lambda}\,\frac{\lambda^{r}}{r!}, \qquad r = 0, 1, 2, 3, \dots \]
$X$ models the \textbf{number of events} occurring in a fixed interval (of time, length, area or volume) when the events occur:
\begin{itemize}
\item \textbf{randomly} (no predictable pattern);
\item \textbf{independently} (one event does not influence the next);
\item at a \textbf{constant average rate} $\lambda$ per unit interval;
\item \textbf{singly} (two events cannot occur at exactly the same instant).
\end{itemize}
The parameter $\lambda$ is the \emph{mean number of events per unit interval}.
\end{definition}

\begin{attention}[When the Poisson model fails]
Do not use a Poisson model if events arrive in clusters (e.g. students leaving school in groups at closing time: not independent), or if the rate changes over the interval (traffic at 7\,a.m.\ versus midnight: rate not constant). Always check the four conditions before writing $X \sim \mathrm{Po}(\lambda)$.
\end{attention}

\begin{propriete}[The Poisson probabilities form a genuine distribution]
The probabilities $P(X=r) = e^{-\lambda}\lambda^{r}/r!$ are non-negative and sum to $1$. The proof is examinable.
\end{propriete}

\textbf{Proof.} Using the exponential series $e^{x} = \displaystyle\sum_{r=0}^{\infty} \frac{x^{r}}{r!}$, valid for every real $x$:
\[ \sum_{r=0}^{\infty} P(X=r) = \sum_{r=0}^{\infty} e^{-\lambda}\frac{\lambda^{r}}{r!} = e^{-\lambda}\sum_{r=0}^{\infty}\frac{\lambda^{r}}{r!} = e^{-\lambda}\,e^{\lambda} = e^{0} = 1. \]
Since also $e^{-\lambda} > 0$ and $\lambda^{r}/r! \geq 0$, every probability is non-negative. $\blacksquare$

\courssub{Mean and Variance of the Poisson Distribution}

The most remarkable feature of the Poisson distribution is that its mean and variance are \emph{equal}.

\begin{propriete}[Mean and variance]
If $X \sim \mathrm{Po}(\lambda)$, then
\[ E(X) = \lambda \qquad \text{and} \qquad \mathrm{Var}(X) = \lambda. \]
The equality $E(X) = \mathrm{Var}(X)$ is the \emph{signature} of the Poisson distribution. The proof is examinable.
\end{propriete}

\textbf{Proof of $E(X) = \lambda$.}
\begin{align*}
E(X) &= \sum_{r=0}^{\infty} r\, e^{-\lambda}\frac{\lambda^{r}}{r!} = e^{-\lambda}\sum_{r=1}^{\infty} \frac{\lambda^{r}}{(r-1)!} \quad (\text{the } r=0 \text{ term vanishes})\\
&= e^{-\lambda}\,\lambda \sum_{r=1}^{\infty} \frac{\lambda^{r-1}}{(r-1)!} = \lambda\, e^{-\lambda} \sum_{k=0}^{\infty} \frac{\lambda^{k}}{k!} = \lambda\, e^{-\lambda} e^{\lambda} = \lambda.
\end{align*}

\textbf{Proof of $\mathrm{Var}(X) = \lambda$ via $E\big(X(X-1)\big)$.} We use the identity $\mathrm{Var}(X) = E\big(X(X-1)\big) + E(X) - \big(E(X)\big)^{2}$, which follows from $X^{2} = X(X-1) + X$. First:
\begin{align*}
E\big(X(X-1)\big) &= \sum_{r=0}^{\infty} r(r-1)\, e^{-\lambda}\frac{\lambda^{r}}{r!} = e^{-\lambda}\sum_{r=2}^{\infty} \frac{\lambda^{r}}{(r-2)!}\\
&= \lambda^{2} e^{-\lambda} \sum_{r=2}^{\infty} \frac{\lambda^{r-2}}{(r-2)!} = \lambda^{2} e^{-\lambda} e^{\lambda} = \lambda^{2}.
\end{align*}
Therefore $\mathrm{Var}(X) = \lambda^{2} + \lambda - \lambda^{2} = \lambda$. $\blacksquare$

\begin{methode}[Quick check of a Poisson model from data]
\begin{enumerate}
\item Compute the sample mean $\bar{x}$ and the sample variance $s^{2}$ of the counts.
\item If $\bar{x} \approx s^{2}$, a Poisson model is plausible; take $\lambda \approx \bar{x}$.
\item If $s^{2}$ is very different from $\bar{x}$, reject the Poisson model.
\end{enumerate}
\end{methode}

\courssub{The Recurrence Formula and the Mode}

Computing $P(X=r)$ directly involves large powers and factorials. The following recurrence makes successive probabilities easy to obtain and also locates the \cle{mode}.

\begin{propriete}[Recurrence formula for the Poisson distribution]
If $X \sim \mathrm{Po}(\lambda)$, then
\[ \frac{P(X = r+1)}{P(X = r)} = \frac{\lambda}{r+1}, \qquad \text{i.e.} \qquad P(X = r+1) = \frac{\lambda}{r+1}\,P(X=r). \]
Consequently the probabilities \emph{increase} while $r + 1 < \lambda$ and \emph{decrease} once $r + 1 > \lambda$: the mode is the integer part of $\lambda$ (and if $\lambda$ is itself an integer, both $\lambda - 1$ and $\lambda$ are modes, with equal probabilities).
\end{propriete}

\textbf{Proof.}
\[ \frac{P(X=r+1)}{P(X=r)} = \frac{e^{-\lambda}\lambda^{r+1}/(r+1)!}{e^{-\lambda}\lambda^{r}/r!} = \frac{\lambda}{r+1}. \]
The ratio exceeds $1$ exactly when $r + 1 < \lambda$, which gives the statement on the mode. $\blacksquare$

\begin{exemplebox}[Using the recurrence to find the mode]
Let $X \sim \mathrm{Po}(4)$. Starting from $P(X=0) = e^{-4} \approx 0.0183$:
\begin{align*}
P(X=1) &= \tfrac{4}{1}P(X=0) \approx 0.0733, &
P(X=2) &= \tfrac{4}{2}P(X=1) \approx 0.1465,\\
P(X=3) &= \tfrac{4}{3}P(X=2) \approx 0.1954, &
P(X=4) &= \tfrac{4}{4}P(X=3) \approx 0.1954,\\
P(X=5) &= \tfrac{4}{5}P(X=4) \approx 0.1563.
\end{align*}
The probabilities stop increasing at $r = 4$; since $\lambda = 4$ is an integer, $P(X=3) = P(X=4)$ and the distribution has the two modes $3$ and $4$.
\end{exemplebox}

\courssub{The Shape of the Distribution}

The bar charts below show $\mathrm{Po}(1)$, $\mathrm{Po}(4)$ and $\mathrm{Po}(10)$. For small $\lambda$ the distribution is strongly skewed to the right; as $\lambda$ grows it becomes more symmetric and bell-shaped.

\begin{popfigure}
\begin{center}
\begin{tikzpicture}
\begin{axis}[ybar, bar width=6pt, width=0.32\linewidth, height=5cm, title={$\mathrm{Po}(1)$}, xlabel={$r$}, ylabel={$P(X=r)$}, xmin=-0.5, xmax=6.5, ymin=0, ymax=0.42, xtick={0,1,2,3,4,5,6}, tick label style={font=\scriptsize}, title style={font=\small}, label style={font=\scriptsize}]
\addplot[fill=popblue, draw=popdark] coordinates {(0,0.3679) (1,0.3679) (2,0.1839) (3,0.0613) (4,0.0153) (5,0.0031) (6,0.0005)};
\end{axis}
\end{tikzpicture}\hfill
\begin{tikzpicture}
\begin{axis}[ybar, bar width=5pt, width=0.32\linewidth, height=5cm, title={$\mathrm{Po}(4)$}, xlabel={$r$}, xmin=-0.5, xmax=11.5, ymin=0, ymax=0.22, xtick={0,2,4,6,8,10}, tick label style={font=\scriptsize}, title style={font=\small}, label style={font=\scriptsize}]
\addplot[fill=poporange, draw=popdark] coordinates {(0,0.0183) (1,0.0733) (2,0.1465) (3,0.1954) (4,0.1954) (5,0.1563) (6,0.1042) (7,0.0595) (8,0.0298) (9,0.0132) (10,0.0053) (11,0.0019)};
\end{axis}
\end{tikzpicture}\hfill
\begin{tikzpicture}
\begin{axis}[ybar, bar width=4pt, width=0.32\linewidth, height=5cm, title={$\mathrm{Po}(10)$}, xlabel={$r$}, xmin=-0.5, xmax=20.5, ymin=0, ymax=0.14, xtick={0,5,10,15,20}, tick label style={font=\scriptsize}, title style={font=\small}, label style={font=\scriptsize}]
\addplot[fill=popgreen, draw=popdark] coordinates {(0,0.0000) (1,0.0005) (2,0.0023) (3,0.0076) (4,0.0189) (5,0.0378) (6,0.0631) (7,0.0901) (8,0.1126) (9,0.1251) (10,0.1251) (11,0.1137) (12,0.0948) (13,0.0729) (14,0.0521) (15,0.0347) (16,0.0217) (17,0.0128) (18,0.0071) (19,0.0037) (20,0.0019)};
\end{axis}
\end{tikzpicture}
\end{center}
\legende{Bar charts of $\mathrm{Po}(1)$, $\mathrm{Po}(4)$ and $\mathrm{Po}(10)$: the distribution becomes more symmetric as $\lambda$ increases.}
\end{popfigure}

\courssub{The Unit Interval: Rescaling $\lambda$}

The parameter $\lambda$ is always attached to a \emph{particular} interval. If the question changes the interval, you must rescale $\lambda$ in proportion, because the average rate is constant.

\begin{methode}[Rescaling the parameter]
\begin{enumerate}
\item Identify the rate and its interval (e.g. $4$ calls per minute).
\item Identify the interval asked in the question (e.g. $30$ seconds $= \frac{1}{2}$ minute).
\item Rescale proportionally: $\lambda_{\text{new}} = 4 \times \frac{1}{2} = 2$ calls per $30$ seconds.
\item Use $X \sim \mathrm{Po}(\lambda_{\text{new}})$ for all calculations.
\end{enumerate}
\end{methode}

\begin{popfigure}
\begin{center}
\begin{tikzpicture}[scale=1]
\draw[->, thick, popdark] (0,0) -- (10.5,0) node[below left, font=\scriptsize] {time};
\foreach \x/\lab in {0/0, 2.5/30\,s, 5/60\,s, 7.5/90\,s, 10/120\,s}{
  \draw[thick, popdark] (\x,0.08) -- (\x,-0.08) node[below, font=\scriptsize] {\lab};
}
\foreach \x in {0.4, 1.3, 2.1, 3.2, 4.6, 5.5, 6.4, 7.2, 8.6, 9.4}{
  \draw[poporange, very thick] (\x,0.25) -- (\x,0.75);
  \fill[poporange] (\x,0.5) circle (1.6pt);
}
\draw[decorate, decoration={brace, amplitude=5pt}, popblue, thick] (0,1.0) -- (2.5,1.0) node[midway, above=6pt, font=\scriptsize, popblue] {$\lambda = 2$};
\draw[decorate, decoration={brace, amplitude=5pt}, popgreen, thick] (0,1.7) -- (5,1.7) node[midway, above=6pt, font=\scriptsize, popgreen] {$\lambda = 4$ per minute};
\draw[decorate, decoration={brace, amplitude=5pt}, poppurple, thick] (0,-0.9) -- (10,-0.9) node[midway, below=6pt, font=\scriptsize, poppurple] {$\lambda = 8$ per $2$ minutes};
\end{tikzpicture}
\end{center}
\legende{Random events (orange ticks) on a timeline: the rate is constant, so $\lambda$ scales with the length of the interval.}
\end{popfigure}

\begin{exoresolu}[Rescaling and ``at least one'']
Calls arrive at an office in Yaound\'e at random, independently, at an average rate of $4$ per minute. Find the probability that (a) exactly $2$ calls arrive in a given $30$-second period; (b) at least one call arrives in a given $30$-second period.
\tcblower
\textbf{Solution.} Rescale: $\lambda = 4 \times \frac{30}{60} = 2$ per $30$ seconds, so $X \sim \mathrm{Po}(2)$.

(a) $P(X=2) = e^{-2}\dfrac{2^{2}}{2!} = 2e^{-2} \approx 0.2707$.

(b) $P(X \geq 1) = 1 - P(X=0) = 1 - e^{-2} \approx 1 - 0.1353 = 0.8647$.
\end{exoresolu}

\begin{attention}[Exam tip: the fastest route for ``at least one'']
For $X \sim \mathrm{Po}(\lambda)$,
\[ P(X \geq 1) = 1 - P(X=0) = 1 - e^{-\lambda}. \]
Never sum $P(X=1) + P(X=2) + \dots$: the series is infinite. The complement is one line of work.
\end{attention}

\courssub{Using Cumulative Poisson Tables}

Cumulative Poisson tables give $P(X \leq r)$ for selected values of $\lambda$ and $r$. Every other probability is obtained from these:

\begin{aretenir}[Reading cumulative Poisson tables]
For $X \sim \mathrm{Po}(\lambda)$, with $F(r) = P(X \leq r)$ read from the table:
\begin{itemize}
\item $P(X = r) = P(X \leq r) - P(X \leq r-1)$;
\item $P(X \geq r) = 1 - P(X \leq r-1)$;
\item $P(X > r) = 1 - P(X \leq r)$;
\item $P(a \leq X \leq b) = P(X \leq b) - P(X \leq a-1)$.
\end{itemize}
If $\lambda$ is not tabulated, compute $e^{-\lambda}\lambda^{r}/r!$ directly with a calculator.
\end{aretenir}

\begin{exemplebox}[Worked table reading]
Let $X \sim \mathrm{Po}(3.5)$. From cumulative tables: $P(X \leq 1) = 0.1359$, $P(X \leq 2) = 0.3208$, $P(X \leq 3) = 0.5366$, $P(X \leq 5) = 0.8576$.
\begin{itemize}
\item $P(X = 3) = P(X \leq 3) - P(X \leq 2) = 0.5366 - 0.3208 = 0.2158$;
\item $P(X \geq 4) = 1 - P(X \leq 3) = 1 - 0.5366 = 0.4634$;
\item $P(2 \leq X \leq 5) = P(X \leq 5) - P(X \leq 1) = 0.8576 - 0.1359 = 0.7217$.
\end{itemize}
\end{exemplebox}

\begin{attention}[Off-by-one errors]
The most common table mistake is confusing $P(X \geq r) = 1 - P(X \leq r-1)$ with $1 - P(X \leq r)$. Remember: the Poisson variable is \emph{discrete}, so $P(X \geq r)$ \emph{includes} $r$, and its complement is $P(X \leq r-1)$.
\end{attention}

\courssub{Sum of Two Independent Poisson Variables}

\begin{propriete}[Additivity of the Poisson distribution]
If $X \sim \mathrm{Po}(\lambda)$ and $Y \sim \mathrm{Po}(\mu)$ are \textbf{independent}, then
\[ X + Y \sim \mathrm{Po}(\lambda + \mu). \]
The proof is instructive and may be required.
\end{propriete}

\textbf{Proof.} Let $Z = X + Y$. For $n \geq 0$, the event $\{Z = n\}$ is the disjoint union of the events $\{X = r\} \cap \{Y = n-r\}$ for $r = 0, 1, \dots, n$. Using independence and then the binomial theorem:
\begin{align*}
P(Z = n) &= \sum_{r=0}^{n} P(X=r)\,P(Y = n-r) = \sum_{r=0}^{n} e^{-\lambda}\frac{\lambda^{r}}{r!}\, e^{-\mu}\frac{\mu^{n-r}}{(n-r)!}\\
&= e^{-(\lambda+\mu)} \frac{1}{n!} \sum_{r=0}^{n} \frac{n!}{r!(n-r)!}\, \lambda^{r} \mu^{n-r} = e^{-(\lambda+\mu)}\, \frac{(\lambda+\mu)^{n}}{n!},
\end{align*}
since $\displaystyle\sum_{r=0}^{n}\binom{n}{r}\lambda^{r}\mu^{n-r} = (\lambda+\mu)^{n}$. This is exactly the $\mathrm{Po}(\lambda+\mu)$ formula. $\blacksquare$

Intuitively this is natural: if accidents occur at junction $A$ at an average rate of $2$ per week and independently at junction $B$ at a rate of $1.5$ per week, then the \emph{combined} stream of accidents still occurs randomly, independently and at the constant combined rate $3.5$ per week.

\begin{attention}[The difference is NOT Poisson]
Only the \emph{sum} of independent Poisson variables is Poisson. The difference $X - Y$ takes negative values, so it \textbf{cannot} be Poisson (a Poisson variable only takes values $0, 1, 2, \dots$). For $X - Y$ you must work from first principles, e.g. $P(X > Y) = \sum_{r} P(Y = r)\,P(X \geq r+1)$.
\end{attention}

\begin{exoresolu}[Accidents at two junctions]
Accidents occur independently at two junctions of Douala: at junction $A$ at an average rate of $2$ per week, and at junction $B$ at an average rate of $1.5$ per week. Find the probability that, in a given week, (a) there are exactly $4$ accidents in total at the two junctions; (b) there is at least one accident in total.
\tcblower
\textbf{Solution.} Let $X \sim \mathrm{Po}(2)$ and $Y \sim \mathrm{Po}(1.5)$, independent. Then $T = X + Y \sim \mathrm{Po}(3.5)$.

(a) $P(T = 4) = e^{-3.5}\dfrac{3.5^{4}}{4!} = e^{-3.5}\dfrac{150.0625}{24} \approx 0.0302 \times 6.2526 \approx 0.1888$.

(b) $P(T \geq 1) = 1 - e^{-3.5} \approx 1 - 0.0302 = 0.9698$.
\end{exoresolu}

\begin{exercice}[Consolidation]
\begin{enumerate}
\item Bacteria occur randomly in a liquid at an average rate of $5$ per millilitre. A sample of $0.4$ ml is examined. Find the probability that it contains (a) exactly $3$ bacteria; (b) at least one bacterium.
\item Given $X \sim \mathrm{Po}(6)$ and cumulative values $P(X \leq 4) = 0.2851$, $P(X \leq 8) = 0.8472$, find $P(5 \leq X \leq 8)$ and $P(X \geq 9)$.
\item $X \sim \mathrm{Po}(2.4)$ and $Y \sim \mathrm{Po}(3.1)$ are independent. State the distribution of $X + Y$ and compute $P(X + Y = 0)$ and the common value of $E(X+Y)$ and $\mathrm{Var}(X+Y)$.
\item For $X \sim \mathrm{Po}(2.5)$, use the recurrence formula to find the mode of the distribution.
\end{enumerate}
\end{exercice}

\begin{corrige}[Exercise 1]
\begin{enumerate}
\item Rescale: $\lambda = 5 \times 0.4 = 2$, so $X \sim \mathrm{Po}(2)$. (a) $P(X=3) = e^{-2}\frac{2^{3}}{3!} = \frac{4}{3}e^{-2} \approx 0.1804$. (b) $P(X \geq 1) = 1 - e^{-2} \approx 0.8647$.
\item $P(5 \leq X \leq 8) = P(X \leq 8) - P(X \leq 4) = 0.8472 - 0.2851 = 0.5621$; $P(X \geq 9) = 1 - P(X \leq 8) = 0.1528$.
\item $X + Y \sim \mathrm{Po}(5.5)$; $P(X+Y=0) = e^{-5.5} \approx 0.0041$; $E(X+Y) = \mathrm{Var}(X+Y) = 5.5$.
\item The probabilities increase while $r + 1 < 2.5$, i.e. for $r = 0, 1$, and decrease from $r + 1 = 3$ onwards. Hence the mode is $r = 2$: indeed $P(X=2) = e^{-2.5}\frac{2.5^{2}}{2!} \approx 0.2565$, larger than $P(X=1) \approx 0.2052$ and $P(X=3) \approx 0.2138$.
\end{enumerate}
\end{corrige}

\begin{savaistu}[Poisson and the Prussian army]
Poisson introduced this law in 1837 in a work on the probability of judgments in criminal and civil matters. It became famous through Bortkiewicz's 1898 study of the number of Prussian soldiers killed each year by horse kicks: the yearly counts matched a Poisson distribution remarkably well. Today the same law models calls to mobile networks, arrivals at bank counters in Douala, and radioactive decays.
\end{savaistu}

\begin{aretenir}[Key points of the section]
\begin{itemize}
\item $X \sim \mathrm{Po}(\lambda)$: $P(X=r) = e^{-\lambda}\lambda^{r}/r!$, for events random, independent, at constant rate $\lambda$.
\item $E(X) = \mathrm{Var}(X) = \lambda$: equal mean and variance is the Poisson signature.
\item Recurrence: $P(X=r+1) = \dfrac{\lambda}{r+1}P(X=r)$; the mode is the integer part of $\lambda$ (two modes $\lambda-1$ and $\lambda$ if $\lambda$ is an integer).
\item Always \textbf{rescale} $\lambda$ to the interval asked before computing.
\item Tables give $P(X \leq r)$; deduce single values, tails and intervals by subtraction, watching the $r-1$ shifts.
\item $P(X \geq 1) = 1 - e^{-\lambda}$.
\item Independent sums: $\mathrm{Po}(\lambda) + \mathrm{Po}(\mu) = \mathrm{Po}(\lambda+\mu)$; the \emph{difference} is never Poisson.
\end{itemize}
\end{aretenir}

\coursec{Poisson Approximation to the Binomial and Link with Frequency Distributions}

In the previous sections we discovered the Poisson distribution as a model for \emph{rare events occurring at random in time or space}: calls to a call centre, accidents at a junction on the Yaound\'e--Douala road, misprints per page. We also noticed something striking: the Poisson probabilities $P(X=r)=e^{-\lambda}\lambda^r/r!$ are far easier to compute than binomial probabilities, which involve large binomial coefficients $\binom{n}{r}$. This final section exploits that observation. We shall see that when $n$ is large and $p$ is small, the binomial distribution $\mathrm{B}(n,p)$ is \emph{almost indistinguishable} from a Poisson distribution with the same mean $\lambda=np$. We then study an important property of the Poisson distribution — the distribution of the sum of two independent Poisson variables — before closing the chapter by connecting all our theoretical distributions to real data through \emph{frequency distributions}, which is exactly how a statistician checks whether a model fits reality.

\courssub{Why We Need an Approximation}

Imagine a factory producing $n=2000$ bottle tops per day, each defective with probability $p=0.002$. The number $X$ of defective tops follows $\mathrm{B}(2000,\,0.002)$. To find $P(X=5)$ exactly we need
\[
P(X=5)=\binom{2000}{5}(0.002)^5(0.998)^{1995}.
\]
The binomial coefficient $\binom{2000}{5}$ has more than a dozen digits, and $(0.998)^{1995}$ is an extremely small number requiring careful handling. For a GCE candidate with a basic calculator, and even for a computer in extreme cases, this is heavy and error-prone. Yet the \emph{mean} of this distribution is simply $np=4$, a perfectly moderate number. The idea of the Poisson approximation is: \emph{when events are numerous but individually rare, only the mean matters}, and the Poisson law with $\lambda=np$ reproduces the binomial probabilities to excellent accuracy.

\begin{propriete}[Poisson approximation to the binomial]
Let $X\sim\mathrm{B}(n,p)$. If $n$ is \cle{large} and $p$ is \cle{small} so that the mean $\lambda=np$ is \cle{moderate}, then
\[
X\ \text{is approximately}\ \mathrm{Po}(np), \qquad\text{i.e.}\qquad P(X=r)\approx e^{-np}\,\frac{(np)^r}{r!}.
\]
\emph{Rule of thumb used at GCE:} the approximation is acceptable when $n\geq 50$ and $p\leq 0.1$ (some texts use $np\leq 10$ as an additional safeguard). The approximation improves as $n$ grows and $p$ shrinks with $np$ fixed.
\end{propriete}

\begin{attention}[The variance changes slightly]
The approximation preserves the \cle{mean}: both $\mathrm{B}(n,p)$ and $\mathrm{Po}(np)$ have mean $np$. But the binomial variance is $npq$ while the Poisson variance is $np$. Since $q=1-p$, the two variances agree only when $q\approx 1$, i.e.\ when $p$ is small. This is precisely why the condition $p\leq 0.1$ matters: with $p=0.1$, the variance is underestimated by about $10\%$; with $p=0.5$ the approximation would be disastrous.
\end{attention}

\courssub{Heuristic Justification (Enrichment for Strong Candidates)}

\begin{experience}[From the binomial formula to the Poisson formula]
Fix $\lambda>0$ and let $X_n\sim\mathrm{B}(n,\lambda/n)$, so that $np=\lambda$ for every $n$. For a fixed integer $r$,
\begin{align*}
P(X_n=r)&=\binom{n}{r}\left(\frac{\lambda}{n}\right)^{r}\left(1-\frac{\lambda}{n}\right)^{n-r}\\
&=\frac{n(n-1)\cdots(n-r+1)}{r!}\cdot\frac{\lambda^{r}}{n^{r}}\cdot\left(1-\frac{\lambda}{n}\right)^{n}\left(1-\frac{\lambda}{n}\right)^{-r}\\
&=\frac{\lambda^{r}}{r!}\cdot\underbrace{\frac{n(n-1)\cdots(n-r+1)}{n^{r}}}_{\to\,1}\cdot\underbrace{\left(1-\frac{\lambda}{n}\right)^{n}}_{\to\,e^{-\lambda}}\cdot\underbrace{\left(1-\frac{\lambda}{n}\right)^{-r}}_{\to\,1}.
\end{align*}
The first bracket is a product of $r$ factors $\left(1-\tfrac{k}{n}\right)$, each tending to $1$; the classical limit $\left(1-\tfrac{\lambda}{n}\right)^n\to e^{-\lambda}$ gives the exponential; the last bracket tends to $1$. Hence
\[
P(X_n=r)\ \longrightarrow\ e^{-\lambda}\frac{\lambda^{r}}{r!}\quad\text{as } n\to\infty,
\]
which is exactly the Poisson probability. This derivation is \emph{not examinable}, but it explains \emph{why} the rule of thumb works: for large $n$ and small $p$, the binomial formula is already ``almost'' the Poisson formula.
\end{experience}

\courssub{Sum of Independent Poisson Variables (Lesson 132)}

A fundamental property of the Poisson distribution — often used in conjunction with the approximation — is that the sum of independent Poisson variables is again Poisson. This makes the Poisson family \emph{closed under addition}, a fact that simplifies many real-world models (e.g., total calls from two independent switchboards).

\begin{propriete}[Sum of independent Poisson variables]
If $X\sim\mathrm{Po}(\lambda)$ and $Y\sim\mathrm{Po}(\mu)$ are independent, then
\[
X+Y \sim \mathrm{Po}(\lambda+\mu).
\]
\emph{Proof (examinable):} Use the moment generating function (MGF) or the convolution formula.
MGF: $M_X(t)=e^{\lambda(e^t-1)}$, $M_Y(t)=e^{\mu(e^t-1)}$. By independence, $M_{X+Y}(t)=M_X(t)M_Y(t)=e^{(\lambda+\mu)(e^t-1)}$, which is the MGF of $\mathrm{Po}(\lambda+\mu)$.
Convolution: $P(X+Y=r)=\sum_{k=0}^r P(X=k)P(Y=r-k) = \sum_{k=0}^r e^{-\lambda}\frac{\lambda^k}{k!} e^{-\mu}\frac{\mu^{r-k}}{(r-k)!} = e^{-(\lambda+\mu)}\frac{1}{r!}\sum_{k=0}^r \binom{r}{k}\lambda^k\mu^{r-k} = e^{-(\lambda+\mu)}\frac{(\lambda+\mu)^r}{r!}$.
\end{propriete}

\begin{exemplebox}[Combining Poisson processes]
The number of emails received by a teacher follows $\mathrm{Po}(3)$ per hour, and the number of SMS messages follows $\mathrm{Po}(2)$ per hour, independently. The total number of electronic messages received in one hour follows $\mathrm{Po}(3+2)=\mathrm{Po}(5)$. The probability of receiving exactly $4$ messages in an hour is $e^{-5}5^4/4! \approx 0.1755$.
\end{exemplebox}

\courssub{Worked Comparison: Exact Binomial versus Poisson Approximation}

\begin{exoresolu}[Comparing $\mathrm{B}(100,\,0.04)$ with $\mathrm{Po}(4)$]
A machine produces $100$ items per batch; each item is defective independently with probability $0.04$. Let $X$ be the number of defective items in a batch. (a) Compute the exact $P(X=2)$ and $P(X=5)$. (b) Compute the Poisson approximations. (c) Comment.
\tcblower
\textbf{Solution.} Here $n=100\geq 50$, $p=0.04\leq 0.1$ and $np=4$: the rule of thumb is satisfied.

\textbf{(a) Exact binomial values.}
\begin{align*}
P(X=2)&=\binom{100}{2}(0.04)^2(0.96)^{98}=4950\times 0.0016\times (0.96)^{98}\approx 0.1450,\\
P(X=5)&=\binom{100}{5}(0.04)^5(0.96)^{95}\approx 0.1595.
\end{align*}

\textbf{(b) Poisson approximation with $\lambda=np=4$.}
\begin{align*}
P(X=2)&\approx e^{-4}\frac{4^2}{2!}=8e^{-4}\approx 0.1465,\\
P(X=5)&\approx e^{-4}\frac{4^5}{5!}=\frac{1024}{120}e^{-4}\approx 0.1563.
\end{align*}

\textbf{(c) Comment.} The absolute errors are about $0.0015$ and $0.0032$ — roughly $1$--$2\%$ relative error. For practical purposes (and for GCE marking) the approximation is excellent, and it required no binomial coefficients at all.
\end{exoresolu}

The figure below displays the two distributions side by side: the bars are barely distinguishable.

\begin{popfigure}
\begin{tikzpicture}
\begin{axis}[
    ybar, bar width=4pt,
    width=12cm, height=7cm,
    xlabel={$r$}, ylabel={$P(X=r)$},
    xmin=-0.5, xmax=12.5, ymin=0, ymax=0.22,
    xtick={0,1,2,3,4,5,6,7,8,9,10,11,12},
    legend style={at={(0.97,0.95)}, anchor=north east, font=\small},
    axis lines=left, enlargelimits=false]
\addplot+[ybar, fill=popblue, draw=popdark] coordinates {
(0,0.0169) (1,0.0703) (2,0.1450) (3,0.1973) (4,0.1995) (5,0.1595)
(6,0.1052) (7,0.0589) (8,0.0284) (9,0.0120) (10,0.0045) (11,0.0015) (12,0.0005)};
\addplot+[ybar, fill=poporange, draw=popdark, bar shift=4pt] coordinates {
(0,0.0183) (1,0.0733) (2,0.1465) (3,0.1954) (4,0.1954) (5,0.1563)
(6,0.1042) (7,0.0595) (8,0.0298) (9,0.0132) (10,0.0053) (11,0.0019) (12,0.0006)};
\legend{{$\mathrm{B}(100,0.04)$}, {$\mathrm{Po}(4)$}}
\end{axis}
\end{tikzpicture}
\legende{Exact binomial probabilities $\mathrm{B}(100,0.04)$ (blue) and Poisson approximation $\mathrm{Po}(4)$ (orange). The two distributions share the same mean $4$ and are almost identical.}
\end{popfigure}

\begin{methode}[Applying the Poisson approximation in an exam]
\begin{enumerate}
\item \textbf{Check the conditions:} $n\geq 50$, $p\leq 0.1$ (and ideally $np\leq 10$).
\item \textbf{State the approximation explicitly:} ``$X\sim\mathrm{B}(n,p)\approx \mathrm{Po}(\lambda)$ with $\lambda=np=\ldots$''.
\item \textbf{Compute} with the Poisson formula $P(X=r)=e^{-\lambda}\lambda^r/r!$ or the \textbf{cumulative Poisson tables} (for $P(X\leq r)$).
\item \textbf{Conclude} in the context of the question.
\end{enumerate}
\end{methode}

\begin{attention}[Exam tip: name your approximation]
If a question says ``\emph{use a suitable approximation}'', you \textbf{must} write down the approximating distribution and its parameter \emph{before} any calculation, e.g.\ ``Since $n=400$ is large and $p=0.01$ is small, $X\approx\mathrm{Po}(4)$.'' Markers award a specific mark for this statement; a bare numerical answer without it loses credit.
\end{attention}

\courssub{Choosing the Right Model: Binomial, Geometric, Poisson or Approximation}

A recurring GCE skill is recognising \emph{which} distribution a situation describes. Use the following decision criteria.

\begin{center}
\begin{tabularx}{\linewidth}{|l|X|X|}
\hline
\rowcolor{popblueL} \textbf{Situation} & \textbf{Model} & \textbf{Key signal} \\
\hline
Fixed number $n$ of independent trials, count of successes, $p$ constant & $\mathrm{B}(n,p)$ & ``out of $n$ \ldots how many \ldots'' \\
\hline
Trials repeated until the \emph{first} success & Geometric, $P(X=r)=q^{r-1}p$ & ``first time'', ``until one succeeds'' \\
\hline
Events at random in time/space, known average rate $\lambda$ per unit interval & $\mathrm{Po}(\lambda)$ & ``on average $\lambda$ per hour/page/km'' \\
\hline
Binomial with $n\geq 50$, $p\leq 0.1$ & $\mathrm{Po}(np)$ approximation & ``use a suitable approximation'' \\
\hline
\end{tabularx}
\end{center}

\begin{exemplebox}[Quick classification]
\begin{itemize}
\item A student answers $20$ multiple-choice questions at random, each with $4$ options: number correct $\sim\mathrm{B}(20,0.25)$ — \emph{not} approximable by Poisson since $p=0.25>0.1$.
\item A typist makes on average $3$ errors per page: errors per page $\sim\mathrm{Po}(3)$.
\item A player rolls a die until the first six: number of rolls is geometric with $p=\tfrac16$.
\item A lottery sells $5000$ tickets, each winning the jackpot with probability $0.0002$: number of jackpot winners $\sim\mathrm{B}(5000,0.0002)\approx\mathrm{Po}(1)$.
\end{itemize}
\end{exemplebox}

\courssub{Probability Distributions and Frequency Distributions (Lesson 133)}

So far our distributions have been \emph{theoretical}: they predict probabilities. Real life produces \emph{observed frequencies}: a teacher counts how many of her $40$ students scored $0,1,2,\ldots$ points on a quiz; a statistician counts how many league matches ended with $0,1,2,\ldots$ goals. The bridge between the two worlds is the following fundamental idea.

\begin{definition}[Expected (theoretical) frequency]
If a random variable $X$ has $P(X=r)=p_r$ and an experiment is repeated $N$ times independently, the \cle{expected frequency} of the value $r$ is
\[
E_r = N\times P(X=r) = N\,p_r.
\]
The set of pairs $(r, E_r)$ is the \emph{theoretical frequency distribution} associated with the probability model.
\end{definition}

This is simply the statement ``probability = long-run proportion'': if $P(X=2)=0.20$, then out of $N=150$ repetitions we \emph{expect} about $150\times 0.20=30$ occurrences of the value $2$.

\begin{attention}[Expected frequencies need not be integers]
$E_r$ is an \emph{average}, not an actual count. Values like $23.7$ or $8.42$ are perfectly correct. \textbf{Do not round them prematurely}: keep at least one or two decimal places throughout, otherwise accumulated rounding errors distort any later comparison with observed data.
\end{attention}

\courssub{Estimating the Parameter from Data}

In practice the parameter of the model is unknown and must be \emph{estimated from the data}. The natural estimator is the \cle{sample mean}:
\begin{itemize}
\item for a \textbf{Poisson} model, $\lambda$ is the mean, so we take $\hat\lambda=\bar{x}$;
\item for a \textbf{binomial} model with known $n$, the mean is $np$, so we take $\hat p=\bar{x}/n$.
\end{itemize}
We then compute the probabilities $P(X=r)$ with this estimated parameter and multiply by $N$.

\begin{methode}[Fitting a distribution to observed data]
\begin{enumerate}
\item \textbf{Compute the sample mean} $\bar{x}=\dfrac{\sum f_r x_r}{\sum f_r}$ from the observed frequency table.
\item \textbf{Estimate the parameter:} $\hat\lambda=\bar{x}$ (Poisson) or $\hat p=\bar{x}/n$ (binomial).
\item \textbf{Compute the probabilities} $p_r=P(X=r)$ using the fitted model (formula or tables).
\item \textbf{Multiply by the total frequency:} $E_r=N\,p_r$.
\item \textbf{Compare} observed and expected frequencies in a two-row table and comment on the closeness of the fit.
\end{enumerate}
\end{methode}

\begin{exoresolu}[Fitting a Poisson model to goals per match]
Over a season, $N=120$ matches of a Cameroonian league championship gave the following numbers of goals per match:
\begin{center}
\begin{tabular}{|l|cccccc|}
\hline
Goals $r$ & 0 & 1 & 2 & 3 & 4 & 5 or more \\
\hline
Observed frequency $O_r$ & 14 & 31 & 38 & 22 & 11 & 4 \\
\hline
\end{tabular}
\end{center}
Fit a Poisson distribution and compute the expected frequencies.
\tcblower
\textbf{Solution.} \textbf{Step 1 — sample mean.}
\[
\bar{x}=\frac{0(14)+1(31)+2(38)+3(22)+4(11)+5(4)}{120}=\frac{31+76+66+44+20}{120}=\frac{237}{120}=1.975.
\]
(The ``5 or more'' class is counted as $5$; with such a small frequency the effect is negligible.)

\textbf{Step 2 — estimate:} $\hat\lambda=\bar{x}\approx 1.975\approx 2.0$.

\textbf{Step 3 — probabilities from $\mathrm{Po}(2)$:} using $P(X=r)=e^{-2}\,2^r/r!$ with $e^{-2}\approx 0.1353$:
\[
p_0=0.1353,\quad p_1=0.2707,\quad p_2=0.2707,\quad p_3=0.1804,\quad p_4=0.0902,
\]
and $P(X\geq 5)=1-(p_0+p_1+p_2+p_3+p_4)\approx 1-0.9473=0.0527$.

\textbf{Step 4 — expected frequencies $E_r=120\,p_r$:}
\begin{center}
\begin{tabular}{|l|cccccc|}
\hline
Goals $r$ & 0 & 1 & 2 & 3 & 4 & $\geq 5$ \\
\hline
Observed $O_r$ & 14 & 31 & 38 & 22 & 11 & 4 \\
\hline
Expected $E_r$ & 16.24 & 32.48 & 32.48 & 21.65 & 10.83 & 6.32 \\
\hline
\end{tabular}
\end{center}

\textbf{Step 5 — comment.} The expected frequencies track the observed ones closely (e.g.\ $38$ observed versus $32.5$ expected at $r=2$; $22$ versus $21.7$ at $r=3$). A Poisson model with mean about $2$ goals per match describes this championship reasonably well. Note that the expected frequencies were deliberately \emph{not} rounded to integers.
\end{exoresolu}

\begin{popfigure}
\begin{tikzpicture}
\begin{axis}[
    ybar, bar width=6pt,
    width=12cm, height=7cm,
    xlabel={Goals per match $r$}, ylabel={Frequency},
    xmin=-0.5, xmax=5.5, ymin=0, ymax=45,
    xtick={0,1,2,3,4,5},
    xticklabels={0,1,2,3,4,$\geq 5$},
    legend style={at={(0.97,0.95)}, anchor=north east, font=\small},
    axis lines=left]
\addplot+[ybar, fill=popgreen, draw=popdark] coordinates {
(0,14) (1,31) (2,38) (3,22) (4,11) (5,4)};
\addplot+[ybar, fill=popgold, draw=popdark, bar shift=6pt] coordinates {
(0,16.24) (1,32.48) (2,32.48) (3,21.65) (4,10.83) (5,6.32)};
\legend{Observed, Expected $\mathrm{Po}(2)$}
\end{axis}
\end{tikzpicture}
\legende{Observed frequencies of goals per match (green) against expected frequencies from the fitted $\mathrm{Po}(2)$ model (gold). The two profiles are close, supporting the Poisson model.}
\end{popfigure}

\begin{savaistu}[Towards goodness-of-fit tests]
Comparing observed and expected frequencies ``by eye'', as we did above, is the first step of a much deeper theory. In further study you will meet the \emph{chi-squared goodness-of-fit test}, which turns the discrepancies $O_r-E_r$ into a single number $\sum (O_r-E_r)^2/E_r$ and decides objectively whether the model is acceptable. This is beyond the GCE syllabus, but everything you have done here — estimating $\lambda$, computing expected frequencies, keeping decimals — is exactly its preparation.
\end{savaistu}

\begin{exercice}[Binomial fit, Poisson approximation and Sum of Poissons]
\begin{enumerate}
\item In $80$ families with $4$ children, the number of boys was counted, giving mean $2.05$ boys per family. Assuming a binomial model $\mathrm{B}(4,p)$, estimate $p$ and compute the expected frequencies of $0,1,2,3,4$ boys among the $80$ families.
\item A call centre receives on average $1.5$ calls per minute. Calls arrive independently. Use a Poisson model to find the expected number of minutes with no call in one hour ($60$ minutes).
\item The number of defective components in boxes of $200$ is modelled by $\mathrm{B}(200,0.015)$. State a suitable approximation and use it to estimate $P(X\leq 2)$.
\item The number of accidents on Road A follows $\mathrm{Po}(2)$ per month, and on Road B follows $\mathrm{Po}(3)$ per month, independently. Find the probability that there are exactly $4$ accidents in total on both roads in a given month.
\end{enumerate}
\end{exercice}

\begin{corrige}[Exercise 1]
\textbf{1.} $\hat p=\bar{x}/n=2.05/4=0.5125$. With $\mathrm{B}(4,0.5125)$: $p_0=(0.4875)^4\approx 0.0565$, $p_1=4(0.5125)(0.4875)^3\approx 0.2376$, $p_2=6(0.5125)^2(0.4875)^2\approx 0.3745$, $p_3=4(0.5125)^3(0.4875)\approx 0.2625$, $p_4=(0.5125)^4\approx 0.0690$. Expected frequencies ($\times 80$): $4.52,\ 19.01,\ 29.96,\ 21.00,\ 5.52$.

\textbf{2.} Calls per minute $\sim\mathrm{Po}(1.5)$, so $P(X=0)=e^{-1.5}\approx 0.2231$. Expected number of quiet minutes: $60\times 0.2231\approx 13.4$ minutes.

\textbf{3.} $n=200\geq 50$, $p=0.015\leq 0.1$, $np=3$: approximate by $\mathrm{Po}(3)$. Then $P(X\leq 2)\approx e^{-3}\left(1+3+\tfrac{9}{2}\right)=8.5\,e^{-3}\approx 0.4232$ (confirmed by cumulative Poisson tables).

\textbf{4.} Let $X\sim\mathrm{Po}(2)$, $Y\sim\mathrm{Po}(3)$ independent. Total $X+Y\sim\mathrm{Po}(5)$. $P(X+Y=4)=e^{-5}5^4/4! = \frac{625}{24}e^{-5}\approx 0.1755$.
\end{corrige}

\begin{aretenir}[Key points of the section]
\begin{itemize}
\item $\mathrm{B}(n,p)\approx\mathrm{Po}(np)$ when $n\geq 50$ and $p\leq 0.1$; the \textbf{mean} $np$ is preserved, the variance changes from $npq$ to $np$.
\item If $X\sim\mathrm{Po}(\lambda)$ and $Y\sim\mathrm{Po}(\mu)$ are independent, then $X+Y\sim\mathrm{Po}(\lambda+\mu)$.
\item In an exam, \textbf{state the approximating distribution and its parameter} before computing.
\item Choose the model from the situation: fixed $n$ trials $\to$ binomial; waiting for the first success $\to$ geometric; random events at a rate $\to$ Poisson; large $n$, small $p$ $\to$ Poisson approximation.
\item \textbf{Expected frequency} $=$ $N\times P(X=r)$; estimate the parameter from the sample mean ($\hat\lambda=\bar x$ or $\hat p=\bar x/n$).
\item Expected frequencies are \textbf{not rounded}; compare them with observed frequencies to judge the fit of the model.
\end{itemize}
\end{aretenir}

\newpage

\coursec{Integration activity}

\coursec{Special Discrete Probability Distributions}

\courssub{Aims of the Activity}

By the end of this activity, you should be able to:

\begin{itemize}
\item recognise, from a concrete description, which of the \cle{binomial}, \cle{Poisson}, \cle{geometric} or \cle{uniform} models is appropriate, and justify the choice by checking the conditions of each model;
\item estimate the parameter $\lambda$ of a Poisson distribution from observed data and compare \emph{observed frequencies} with \emph{expected frequencies};
\item compute means, variances and probabilities for $B(n,p)$, and use the \cle{Poisson approximation} $Po(np)$ when $n$ is large and $p$ is small;
\item use \cle{cumulative tables} (binomial and Poisson) efficiently;
\item apply the \cle{unit-interval scaling} of the Poisson mean ($\lambda \propto$ length of interval);
\item communicate a complete modelling argument in clear mathematical English, as required at the GCE Advanced Level.
\end{itemize}

\begin{attention}[Working in Groups]
Work in groups of three or four. Every numerical answer must be accompanied by one sentence stating \emph{which model was used and why}. An unjustified number earns no credit in this activity, exactly as in the examination.
\end{attention}

\courssub{Situation}

The Douala City Council, together with a local taxi agency ``Wouri Express'', has commissioned your Further Mathematics class to analyse three sets of data collected during the first term.

\textbf{Data set A — Accidents at Rond-point Deïdo.} A traffic officer recorded the number of road accidents occurring each day at the Rond-point Deïdo junction over $100$ consecutive days. Accidents on different days may be assumed independent, and they occur randomly at an approximately constant average rate. The observed frequencies are:

\begin{center}
\begin{tabular}{|l|cccccccc|}
\hline
\rowcolor{popblueL} Number of accidents per day $x$ & $0$ & $1$ & $2$ & $3$ & $4$ & $5$ & $6$ & $7$ or more \\
\hline Observed frequency (days) & $15$ & $31$ & $26$ & $17$ & $8$ & $2$ & $1$ & $0$ \\
\hline
\end{tabular}
\end{center}

The sample mean is $\bar{x} = 1.82 \approx 1.8$ accidents per day.

\textbf{Data set B — Expired insurance at a checkpoint.} At a police checkpoint on the Douala--Yaoundé road, long-term records show that $2\%$ of vehicles have expired insurance. Each day, exactly $60$ vehicles are checked, independently of one another. Let $X$ be the number of vehicles found with expired insurance in one day, and let $T$ be the number of vehicles checked \emph{up to and including} the first uninsured vehicle.

\textbf{Data set C — Calls at Wouri Express.} The taxi agency receives booking calls randomly, at a constant average rate of $3$ calls per $5$ minutes, calls arriving independently of one another. Let $Y$ be the number of calls received in $5$ minutes, and $W$ the number received in $2$ minutes.

The Council wants a written report answering the tasks below.

\courssub{Tasks}

\begin{enumerate}
\item \textbf{Choosing a model for accidents.} State, with reasons, why a Poisson distribution is a plausible model for the number of accidents per day. Estimate the parameter $\lambda$ from the data.
\item \textbf{Expected frequencies.} Using $Po(1.8)$ and the cumulative Poisson tables (or the recurrence relation $P(X=x+1)=\dfrac{\lambda}{x+1}P(X=x)$), compute $P(X=x)$ for $x = 0, 1, \dots, 6$ and $P(X \geq 7)$, each to 4 decimal places. Hence compute the \emph{expected frequencies} out of $100$ days and display observed and expected frequencies in one table. Comment on the fit.
\item \textbf{Binomial model for uninsured vehicles.} Justify that $X \sim B(60, 0.02)$. Write down $E(X)$ and $\mathrm{Var}(X)$. Compute $P(X = 0)$ and $P(X \geq 2)$ (you may use cumulative binomial tables or direct calculation).
\item \textbf{Poisson approximation.} Explain why $X$ may be approximated by $Po(1.2)$. Recompute $P(X=0)$ and $P(X \geq 2)$ under this approximation and comment on the closeness of the results.
\item \textbf{Geometric model.} Explain why $T$ follows a geometric distribution with $p = 0.02$. Find the probability that the first uninsured vehicle appears \emph{after} the 20th check, i.e. $P(T > 20)$. Write down $E(T)$ and interpret it in context.
\item \textbf{Calls and unit intervals.} State the distributions of $Y$ and $W$, justifying the value of the mean of $W$. Using cumulative Poisson tables, find $P(Y > 5)$ and $P(W = 0)$.
\item \textbf{Synthesis.} Write a short paragraph (5--8 lines) for the City Council summarising your findings: the fit of the accident model, the expected number of uninsured vehicles per day, and the busiest risk the taxi agency should prepare for.
\end{enumerate}

\begin{methode}[Checklist Before Choosing a Model]
\begin{enumerate}
\item Fixed number $n$ of independent trials, two outcomes, constant $p$? $\Rightarrow$ \textbf{Binomial}.
\item Independent events at a constant average rate in time/space? $\Rightarrow$ \textbf{Poisson}.
\item Number of trials until the first success? $\Rightarrow$ \textbf{Geometric}.
\item Binomial with $n$ large ($n \geq 50$) and $p$ small ($p \leq 0.1$)? $\Rightarrow$ approximate by $Po(np)$.
\end{enumerate}
\end{methode}

\begin{popfigure}
\begin{tikzpicture}[scale=0.8, every node/.style={font=\small}]
% Poisson process timeline for accidents
\draw[thick, ->] (0,0) -- (12,0) node[right] {Time (days)};
\foreach \d in {1,2,...,10} {
    \draw (\d, -0.1) -- (\d, 0.1) node[below] {\d};
}
% Random accident events
\foreach \x/\h in {1.2/0.5, 2.8/0.7, 3.1/0.5, 4.5/0.6, 5.9/0.5, 6.2/0.7, 7.8/0.5, 8.3/0.6, 9.1/0.5, 9.7/0.7} {
    \draw[poporange, thick] (\x, 0) -- (\x, \h);
    \fill[poporange] (\x, \h) circle (2pt);
}
\node[poporange] at (6, 1.2) {Accidents occur randomly};
\node[poporange] at (6, 1.0) {at constant average rate};
% Geometric distribution illustration
\begin{scope}[yshift=-3cm]
\draw[thick, ->] (0,0) -- (12,0) node[right] {Vehicle checks};
\foreach \d in {1,2,...,10} {
    \draw (\d, -0.1) -- (\d, 0.1) node[below] {\d};
}
\foreach \x in {1,2,...,9} {
    \draw[popgreen, thick] (\x, 0) -- (\x, 0.4);
    \fill[popgreen] (\x, 0.4) circle (2pt);
}
\draw[popred, thick] (10, 0) -- (10, 0.6);
\fill[popred] (10, 0.6) circle (2pt);
\node[popgreen] at (5, -0.8) {Insured vehicles (failures)};
\node[popred] at (10, -0.8) {First uninsured (success)};
\node[popred] at (10, 1.0) {$T = 10$};
\end{scope}
\end{tikzpicture}
\legende{Left: Poisson process for daily accidents. Right: Geometric setting for first uninsured vehicle.}
\end{popfigure}

\courssub{Model Answers}

\begin{exoresolu}[Task 1: Model for Accidents]
\textbf{Statement:} State, with reasons, why a Poisson distribution is a plausible model for the number of accidents per day. Estimate the parameter $\lambda$ from the data.

\tcblower
\textbf{Detailed Solution:}

The Poisson distribution models the number of events occurring in a fixed interval of time (or space) when the following four conditions hold:
\begin{enumerate}
\item \textbf{Randomness:} Events occur at random, not at fixed intervals.
\item \textbf{Independence:} The occurrence of an event in one interval does not affect the probability of an event in another non-overlapping interval. (Accidents on different days are independent.)
\item \textbf{Constant average rate:} The average number of events per unit interval is constant. (The problem states ``approximately constant average rate''.)
\item \textbf{Singly:} Two events cannot occur at exactly the same instant; in practice, the probability of more than one accident in a very short sub-interval is negligible.
\end{enumerate}
The description of Data Set A explicitly mentions independence and a constant average rate, and accidents are random, discrete events. Therefore, a Poisson distribution is perfectly appropriate.

For a Poisson random variable $X \sim Po(\lambda)$, the expectation is $E(X) = \lambda$. The standard estimator for $\lambda$ is the sample mean $\bar{x}$.

\[
\hat{\lambda} = \bar{x} = \frac{\sum x \cdot f}{\sum f} = \frac{0(15) + 1(31) + 2(26) + 3(17) + 4(8) + 5(2) + 6(1)}{100} = \frac{182}{100} = 1.82.
\]
We round to one decimal place as suggested: $\hat{\lambda} \approx 1.8$. Thus we fit the model $X \sim Po(1.8)$.
\end{exoresolu}

\begin{exoresolu}[Task 2: Expected Frequencies and Goodness of Fit]
\textbf{Statement:} Using $Po(1.8)$ and the cumulative Poisson tables (or the recurrence relation), compute $P(X=x)$ for $x = 0, 1, \dots, 6$ and $P(X \geq 7)$, each to 4 decimal places. Hence compute the expected frequencies out of $100$ days and display observed and expected frequencies in one table. Comment on the fit.

\tcblower
\textbf{Detailed Solution:}

We use the recurrence relation for the Poisson distribution:
\[
P(X=0) = e^{-\lambda} = e^{-1.8} = 0.165298\ldots \approx 0.1653,
\]
\[
P(X=x+1) = \frac{\lambda}{x+1} P(X=x).
\]

\begin{align*}
P(X=1) &= \frac{1.8}{1} \times 0.1653 = 0.29754 \approx 0.2975,\\
P(X=2) &= \frac{1.8}{2} \times 0.2975 = 0.26775 \approx 0.2678,\\
P(X=3) &= \frac{1.8}{3} \times 0.2678 = 0.16068 \approx 0.1607,\\
P(X=4) &= \frac{1.8}{4} \times 0.1607 = 0.072315 \approx 0.0723,\\
P(X=5) &= \frac{1.8}{5} \times 0.0723 = 0.026028 \approx 0.0260,\\
P(X=6) &= \frac{1.8}{6} \times 0.0260 = 0.00780 \approx 0.0078.
\end{align*}

Sum of probabilities for $x=0$ to $6$:
\[
\sum_{x=0}^{6} P(X=x) = 0.1653 + 0.2975 + 0.2678 + 0.1607 + 0.0723 + 0.0260 + 0.0078 = 0.9974.
\]
Hence
\[
P(X \geq 7) = 1 - 0.9974 = 0.0026.
\]

Expected frequencies are obtained by multiplying each probability by $100$ (the total number of days).

\begin{center}
\begin{tabular}{|l|cccccccc|}
\hline
\rowcolor{popblueL} $x$ & $0$ & $1$ & $2$ & $3$ & $4$ & $5$ & $6$ & $\geq 7$ \\
\hline
Observed frequency & $15$ & $31$ & $26$ & $17$ & $8$ & $2$ & $1$ & $0$ \\
\hline
Expected frequency & $16.5$ & $29.8$ & $26.8$ & $16.1$ & $7.2$ & $2.6$ & $0.8$ & $0.3$ \\
\hline
\end{tabular}
\end{center}

\textbf{Comment on the fit:} The expected frequencies are very close to the observed frequencies across all categories. The largest discrepancies are at $x=0$ (observed 15 vs expected 16.5) and $x=1$ (observed 31 vs expected 29.8), but these differences are small relative to the frequencies. The model $Po(1.8)$ provides an excellent fit to the accident data. (A formal $\chi^2$ goodness-of-fit test would confirm this; it is beyond our syllabus, but the visual agreement is already convincing.)
\end{exoresolu}

\begin{exoresolu}[Task 3: Binomial Model for Uninsured Vehicles]
\textbf{Statement:} Justify that $X \sim B(60, 0.02)$. Write down $E(X)$ and $\mathrm{Var}(X)$. Compute $P(X = 0)$ and $P(X \geq 2)$.

\tcblower
\textbf{Detailed Solution:}

\textbf{Justification for Binomial model:}
\begin{itemize}
\item There is a fixed number of trials: $n = 60$ vehicles checked each day.
\item Each trial has two outcomes: ``expired insurance'' (success) or ``valid insurance'' (failure).
\item The probability of success is constant: $p = 0.02$ (2\% of vehicles have expired insurance).
\item The trials are independent: the insurance status of one vehicle does not affect another.
\end{itemize}
Therefore $X$, the number of vehicles with expired insurance in a day, follows a binomial distribution: $X \sim B(60, 0.02)$.

\textbf{Mean and Variance:}
\[
E(X) = np = 60 \times 0.02 = 1.2,
\]
\[
\mathrm{Var}(X) = np(1-p) = 60 \times 0.02 \times 0.98 = 1.176.
\]

\textbf{Probabilities:}
\[
P(X = 0) = (1-p)^n = 0.98^{60} = 0.29755\ldots \approx 0.2976 \quad \text{(4 d.p.)}.
\]
\[
P(X = 1) = \binom{60}{1} p (1-p)^{59} = 60 \times 0.02 \times 0.98^{59} = 1.2 \times 0.30363\ldots = 0.36436\ldots \approx 0.3644.
\]
\[
P(X \geq 2) = 1 - P(X=0) - P(X=1) = 1 - 0.2976 - 0.3644 = 0.3380.
\]
(Using cumulative binomial tables with $n=60, p=0.02$ gives the same result.)
\end{exoresolu}

\begin{exoresolu}[Task 4: Poisson Approximation to the Binomial]
\textbf{Statement:} Explain why $X$ may be approximated by $Po(1.2)$. Recompute $P(X=0)$ and $P(X \geq 2)$ under this approximation and comment on the closeness of the results.

\tcblower
\textbf{Detailed Solution:}

The Poisson approximation to the binomial distribution is valid when:
\begin{itemize}
\item $n$ is large (typically $n \geq 50$),
\item $p$ is small (typically $p \leq 0.1$),
\item $np$ is moderate (typically $np \leq 5$).
\end{itemize}
Here $n = 60 \geq 50$, $p = 0.02 \leq 0.1$, and $np = 1.2 \leq 5$. All conditions are satisfied. The approximating Poisson distribution has parameter $\lambda = np = 1.2$, so $X \approx Po(1.2)$.

\textbf{Approximate probabilities:}
\[
P(X = 0) = e^{-1.2} = 0.30119\ldots \approx 0.3012.
\]
\[
P(X \geq 2) = 1 - P(X=0) - P(X=1) = 1 - e^{-1.2} - 1.2 e^{-1.2} = 1 - e^{-1.2}(1 + 1.2) = 1 - 0.3012 \times 2.2 = 1 - 0.66264 = 0.33736 \approx 0.3374.
\]

\textbf{Comparison:}
\begin{center}
\begin{tabular}{|l|c|c|c|}
\hline
\rowcolor{popblueL} Probability & Exact Binomial & Poisson Approximation & Absolute Error \\
\hline
$P(X=0)$ & $0.2976$ & $0.3012$ & $0.0036$ \\
\hline
$P(X \geq 2)$ & $0.3380$ & $0.3374$ & $0.0006$ \\
\hline
\end{tabular}
\end{center}

The errors are very small (less than $0.004$). The Poisson approximation is excellent here and much quicker to use with tables, especially when $n$ is large and binomial tables for $n=60$ may not be available.
\end{exoresolu}

\begin{exoresolu}[Task 5: Geometric Model for First Uninsured Vehicle]
\textbf{Statement:} Explain why $T$ follows a geometric distribution with $p = 0.02$. Find $P(T > 20)$. Write down $E(T)$ and interpret it in context.

\tcblower
\textbf{Detailed Solution:}

\textbf{Justification for Geometric model:}
The random variable $T$ counts the number of vehicles checked \emph{up to and including} the first uninsured vehicle. This is exactly the setting of the geometric distribution:
\begin{itemize}
\item We perform independent Bernoulli trials (checking vehicles).
\item Each trial has the same success probability $p = 0.02$ (expired insurance).
\item We stop at the first success.
\item $T$ is the total number of trials needed to obtain the first success.
\end{itemize}
Hence $T \sim \mathrm{Geo}(0.02)$, with probability mass function $P(T = t) = (1-p)^{t-1}p = 0.98^{t-1} \times 0.02$ for $t = 1, 2, 3, \ldots$

\textbf{Probability that the first uninsured vehicle appears after the 20th check:}
The event $\{T > 20\}$ means the first 20 vehicles all have valid insurance (i.e., 20 consecutive failures).
\[
P(T > 20) = (1-p)^{20} = 0.98^{20} = 0.667607\ldots \approx 0.6676 \quad \text{(4 d.p.)}.
\]

\textbf{Expectation and interpretation:}
\[
E(T) = \frac{1}{p} = \frac{1}{0.02} = 50.
\]
\textbf{Interpretation:} On average, the checkpoint officers will need to inspect 50 vehicles before finding the first one with expired insurance. This is a long-run average; on any given day the actual number may be much smaller or larger.
\end{exoresolu}

\begin{exoresolu}[Task 6: Calls and Unit Interval Scaling]
\textbf{Statement:} State the distributions of $Y$ and $W$, justifying the value of the mean of $W$. Using cumulative Poisson tables, find $P(Y > 5)$ and $P(W = 0)$.

\tcblower
\textbf{Detailed Solution:}

\textbf{Distribution of $Y$ (calls in 5 minutes):}
Calls arrive randomly, independently, and at a constant average rate. The number of calls in a fixed time interval follows a Poisson distribution. The rate is 3 calls per 5 minutes, so for a 5-minute interval:
\[
Y \sim Po(3).
\]

\textbf{Distribution of $W$ (calls in 2 minutes):}
The Poisson mean is proportional to the length of the time interval. The rate is $3$ calls per $5$ minutes, so the rate per minute is $3/5 = 0.6$ calls/minute. For a 2-minute interval:
\[
\lambda_W = 0.6 \times 2 = 1.2 \quad \text{or directly} \quad \lambda_W = 3 \times \frac{2}{5} = 1.2.
\]
Thus $W \sim Po(1.2)$.

\textbf{Probabilities from cumulative Poisson tables:}
\begin{itemize}
\item For $Y \sim Po(3)$: $P(Y > 5) = 1 - P(Y \leq 5)$. From tables, $P(Y \leq 5) = 0.9161$. Hence $P(Y > 5) = 1 - 0.9161 = 0.0839$.
\item For $W \sim Po(1.2)$: $P(W = 0) = e^{-1.2} = 0.3012$ (directly from formula or tables).
\end{itemize}
\end{exoresolu}

\begin{exoresolu}[Task 7: Synthesis for the City Council]
\textbf{Statement:} Write a short paragraph (5--8 lines) summarising your findings.

\tcblower
\textbf{Detailed Solution:}

``The daily number of accidents at Rond-point Deïdo is well modelled by a Poisson distribution with mean $1.8$: the expected frequencies match the observed ones closely, so the junction averages just under two accidents per day, with about $17$ accident-free days per $100$. At the checkpoint, an average of $1.2$ vehicles per day (out of $60$) have expired insurance, and there is roughly a one-in-three chance ($33.8\%$) of catching at least two offenders in a day; the first offender typically appears around the 50th vehicle checked. Wouri Express should expect busy spells: there is about an $8.4\%$ chance of receiving more than $5$ calls in $5$ minutes, while a completely silent $2$-minute spell occurs with probability about $0.30$.''
\end{exoresolu}

\courssub{Going Further}

\begin{exercice}[Extension Questions]
\begin{enumerate}
\item Using the fitted model $Po(1.8)$, find the probability of \emph{at least one} day with $4$ or more accidents in a week of $7$ independent days. \emph{Hint: first find $p = P(X \geq 4)$, then use a binomial model with $n = 7$.}
\item The taxi agency merges its calls with those of a partner agency receiving calls independently at rate $2$ per $5$ minutes. State the distribution of the total number of calls per $5$ minutes and compute the probability of exactly $4$ calls.
\end{enumerate}
\end{exercice}

\begin{corrige}[Extension Questions]
\textbf{1.} From the Poisson tables (or recurrence) with $\lambda = 1.8$:
\[
P(X \leq 3) = P(X=0)+P(X=1)+P(X=2)+P(X=3) = 0.1653 + 0.2975 + 0.2678 + 0.1607 = 0.8913.
\]
Thus $p = P(X \geq 4) = 1 - 0.8913 = 0.1087$.
Let $N$ be the number of days (out of 7) with $\geq 4$ accidents. Since days are independent, $N \sim B(7, 0.1087)$.
\[
P(N \geq 1) = 1 - P(N=0) = 1 - (1 - 0.1087)^7 = 1 - 0.8913^7 = 1 - 0.4468 = 0.5532.
\]

\textbf{2.} The sum of two independent Poisson random variables is Poisson with mean equal to the sum of the means.
Total calls per 5 minutes $\sim Po(3 + 2) = Po(5)$.
\[
P(\text{Total} = 4) = \frac{e^{-5} \cdot 5^4}{4!} = \frac{0.0067379 \times 625}{24} = \frac{4.2112}{24} = 0.1755 \quad \text{(4 d.p.)}.
\]
\end{corrige}

\begin{aretenir}[Key Points Consolidated by This Activity]
\begin{itemize}
\item A Poisson parameter $\lambda$ is estimated by the sample mean; a model is validated by comparing observed and expected frequencies.
\item Binomial $B(n,p)$: $E(X) = np$, $\mathrm{Var}(X) = np(1-p)$; approximated by $Po(np)$ for large $n$ ($n \geq 50$), small $p$ ($p \leq 0.1$).
\item Geometric $\mathrm{Geo}(p)$: $P(T > k) = (1-p)^k$, $E(T) = \dfrac{1}{p}$.
\item Poisson means scale linearly with the interval length: rate $3$ per $5$ min $\Rightarrow$ mean $1.2$ per $2$ min.
\item The sum of independent Poisson variables is Poisson with the sum of the means.
\end{itemize}
\end{aretenir}

\begin{savaistu}[Did You Know?]
The Poisson distribution is named after Siméon Denis Poisson (1781--1840), a French mathematician who introduced it in 1837. In Cameroon, Poisson processes are used to model:
\begin{itemize}
\item Call arrivals at MTN/Orange call centres (teletraffic engineering).
\item Customer arrivals at banks (e.g., BICEC, Afriland First Bank) for queue management.
\item Defects in manufactured goods (e.g., at SOSUCAM or CICAM factories).
\item Rainfall events in Douala during the rainy season for hydrological planning.
\end{itemize}
The geometric distribution is the discrete analogue of the exponential distribution, which models waiting times between Poisson events.
\end{savaistu}

\newpage

\coursec{Methods and worked examples}

\coursec{Special Discrete Probability Distributions}

\courssub{The Binomial Distribution: Probabilities, Mean and Variance}

\begin{definition}[Binomial Experiment]
A \cle{binomial experiment} satisfies four conditions:
\begin{enumerate}
\item There is a \textbf{fixed} number $n$ of trials.
\item Each trial has exactly two outcomes: \textbf{success} or \textbf{failure}.
\item The probability of success $p$ is \textbf{constant} from trial to trial.
\item The trials are \textbf{independent}.
\end{enumerate}
If $X$ is the number of successes in $n$ trials, we write $X \sim B(n, p)$ and set $q = 1 - p$.
\end{definition}

\begin{propriete}[Probability Mass Function, Mean and Variance of $B(n,p)$]
For $X \sim B(n, p)$ with $q = 1-p$:
\begin{align*}
P(X = r) &= \binom{n}{r} p^{r} q^{\,n-r}, \qquad r = 0, 1, \ldots, n,\\[4pt]
E(X) &= np,\\[4pt]
\operatorname{Var}(X) &= npq,\\[4pt]
\sigma_X &= \sqrt{npq}.
\end{align*}
\textbf{Derivation of $E(X)$ and $\operatorname{Var}(X)$ (examinable):}
Write $X = \sum_{i=1}^{n} X_i$ where $X_i \sim \mathrm{Bernoulli}(p)$ are independent.
Then $E(X_i) = p$, $\operatorname{Var}(X_i) = pq$.
By linearity of expectation and independence:
$E(X) = \sum E(X_i) = np$,
$\operatorname{Var}(X) = \sum \operatorname{Var}(X_i) = npq$.
\end{propriete}

\begin{methode}[Computing binomial probabilities, mean and variance]
\begin{enumerate}
\item \textbf{Check the four conditions} for a binomial model (fixed $n$, two outcomes, constant $p$, independence).
\item \textbf{Identify $n$ and $p$}, write $X \sim B(n, p)$ and $q = 1 - p$.
\item \textbf{Exact probability}: $P(X = r) = \binom{n}{r} p^{r} q^{\,n-r}$.
\item \textbf{Cumulative probabilities}: use complement when shorter:
$P(X \geq 1) = 1 - P(X = 0) = 1 - q^{n}$,
$P(X \geq r) = 1 - P(X \leq r-1)$.
\item \textbf{Quote mean and variance directly}:
$E(X) = np$, $\operatorname{Var}(X) = npq$, $\sigma = \sqrt{npq}$.
\end{enumerate}
\textbf{Reflex:} ``at least one'' $\Rightarrow$ compute $1 - q^{n}$ first.
\end{methode}

\begin{exoresolu}[Quality control in a Douala factory]
A machine in a Douala bottling factory produces caps of which $5\%$ are defective. A quality controller takes a random sample of $20$ caps. Let $X$ be the number of defective caps in the sample.
\begin{enumerate}
\item State the distribution of $X$, justifying your answer.
\item Find the probability that exactly $2$ caps are defective.
\item Find the probability that at least one cap is defective.
\item Calculate the mean and the standard deviation of $X$.
\end{enumerate}
\tcblower
\textbf{Solution.}
\begin{enumerate}
\item There is a fixed number of trials ($n = 20$), each cap is either defective or not, the probability of a defective cap is constant ($p = 0.05$), and the caps are independent (random sample from a large production). Hence
\[ X \sim B(20,\ 0.05), \qquad q = 0.95. \]
\item Using the formula with $r = 2$:
\[ P(X = 2) = \binom{20}{2}(0.05)^{2}(0.95)^{18} = 190 \times 0.0025 \times 0.3972 \approx 0.1887. \]
\item Using the complement:
\[ P(X \geq 1) = 1 - P(X = 0) = 1 - (0.95)^{20} = 1 - 0.3585 = 0.6415. \]
\item Mean and variance:
\[ E(X) = np = 20 \times 0.05 = 1, \qquad \operatorname{Var}(X) = npq = 20 \times 0.05 \times 0.95 = 0.95, \]
\[ \sigma = \sqrt{0.95} \approx 0.975. \]
\end{enumerate}
So on average one defective cap is found per sample of $20$, with standard deviation about $0.97$.
\end{exoresolu}

\begin{attention}[Common pitfalls with the binomial model]
\begin{itemize}
\item Do \textbf{not} use the binomial distribution if sampling is \textbf{without replacement} from a small population (use hypergeometric instead). The binomial is valid when the population is large relative to the sample (rule of thumb: sample size $\le 10\%$ of population).
\item The probability $p$ must be the \textbf{same} for every trial. If $p$ changes, it is not binomial.
\item ``At least $r$'' means $r, r+1, \ldots, n$; ``at most $r$'' means $0, 1, \ldots, r$.
\end{itemize}
\end{attention}

\courssub{Cumulative Binomial Tables and the Mode}

\begin{methode}[Reading cumulative binomial tables]
Cumulative tables give $P(X \leq r)$ for $X \sim B(n,p)$ for selected values of $n$, $p$ and $r$.
\begin{enumerate}
\item Locate the block for your $n$, then the column for your $p$, then the row for $r$: the entry is $P(X \leq r)$.
\item Single value: $P(X = r) = P(X \leq r) - P(X \leq r-1)$ (with $P(X \leq -1) = 0$).
\item Strict inequalities: $P(X < r) = P(X \leq r - 1)$.
\item ``Greater than'': $P(X > r) = 1 - P(X \leq r)$ and $P(X \geq r) = 1 - P(X \leq r-1)$.
\item If $p > 0.5$, use the \cle{symmetry trick}: if $X \sim B(n,p)$ then $Y = n - X \sim B(n, 1-p)$, so
\[ P(X \leq r) = P(Y \geq n - r) = 1 - P(Y \leq n - r - 1). \]
\end{enumerate}
\textbf{Reflex:} always rewrite the required probability in terms of $P(X \leq \cdot)$ \emph{before} opening the table.
\end{methode}

\begin{exoresolu}[Using the tables for $B(10, 0.25)$]
A candidate answers a multiple-choice quiz of $10$ questions by guessing; each question has $4$ options, so the probability of a correct answer is $0.25$. The number of correct answers is $X \sim B(10, 0.25)$. Extract from the cumulative tables:
\begin{center}
\begin{tabular}{|c|c|c|c|}
\hline
\rowcolor{popblueL} $r$ & $P(X \leq r)$, $p=0.25$ & $r$ & $P(X \leq r)$, $p=0.25$ \\
\hline
$0$ & $0.0563$ & $4$ & $0.9219$ \\
\hline
$1$ & $0.2440$ & $5$ & $0.9803$ \\
\hline
$2$ & $0.5256$ & $6$ & $0.9965$ \\
\hline
$3$ & $0.7759$ & $7$ & $0.9996$ \\
\hline
\end{tabular}
\end{center}
Find (a) $P(X = 3)$; (b) $P(X < 2)$; (c) $P(X \geq 4)$; (d) $P(2 \leq X \leq 5)$.
\tcblower
\textbf{Solution.} All probabilities are rewritten using cumulative values.
\begin{enumerate}
\item[(a)] $P(X = 3) = P(X \leq 3) - P(X \leq 2) = 0.7759 - 0.5256 = 0.2503$.
\item[(b)] $P(X < 2) = P(X \leq 1) = 0.2440$.
\item[(c)] $P(X \geq 4) = 1 - P(X \leq 3) = 1 - 0.7759 = 0.2241$.
\item[(d)] $P(2 \leq X \leq 5) = P(X \leq 5) - P(X \leq 1) = 0.9803 - 0.2440 = 0.7363$.
\end{enumerate}
\begin{attention}
A very common error is to write $P(X \geq 4) = 1 - P(X \leq 4)$. Since $X$ is discrete, $P(X \geq 4) = 1 - P(X \leq 3)$: the boundary value must be excluded carefully.
\end{attention}
\end{exoresolu}

\begin{propriete}[Recurrence Formula and Mode of $B(n,p)$]
For $X \sim B(n,p)$ with $q=1-p$:
\[ \frac{P(X = r+1)}{P(X = r)} = \frac{n - r}{r + 1}\cdot\frac{p}{q}, \qquad r = 0, 1, \ldots, n-1. \]
The probabilities increase while $\dfrac{n-r}{r+1}\cdot\dfrac{p}{q} > 1$, i.e. while $r + 1 < (n+1)p$.
Hence the \cle{mode} (most likely value) is
\[ m = \lfloor (n+1)p \rfloor. \]
If $(n+1)p$ is an integer, there are \textbf{two modes}: $m = (n+1)p - 1$ and $m = (n+1)p$.
\end{propriete}

\begin{methode}[Finding the mode via the recurrence relation]
\begin{enumerate}
\item Write down the recurrence formula: $\dfrac{P(X = r+1)}{P(X = r)} = \dfrac{n - r}{r + 1}\cdot\dfrac{p}{q}$.
\item The probabilities increase while $\dfrac{n-r}{r+1}\cdot\dfrac{p}{q} > 1$, i.e. while $r < (n+1)p - 1$.
\item Equivalently, compute $(n+1)p$:
\begin{itemize}
\item If $(n+1)p$ is \textbf{not an integer}, the mode is $\lfloor (n+1)p \rfloor$.
\item If $(n+1)p$ \textbf{is an integer}, there are two modes: $(n+1)p - 1$ and $(n+1)p$.
\end{itemize}
\end{enumerate}
\textbf{Reflex:} compute $(n+1)p$ first; if it is not an integer, take the integer part.
\end{methode}

\begin{exoresolu}[Most likely number of sixes]
A fair die is thrown $12$ times. Let $X$ be the number of sixes obtained.
\begin{enumerate}
\item Show that $\dfrac{P(X = r+1)}{P(X=r)} = \dfrac{12 - r}{5(r+1)}$.
\item Hence find the mode of $X$.
\item Verify your answer by computing $P(X=1)$, $P(X=2)$ and $P(X=3)$.
\end{enumerate}
\tcblower
\textbf{Solution.}
\begin{enumerate}
\item Here $X \sim B(12, \tfrac16)$, so $p = \tfrac16$, $q = \tfrac56$ and $\dfrac{p}{q} = \dfrac{1}{5}$. The recurrence formula gives
\[ \frac{P(X = r+1)}{P(X = r)} = \frac{n - r}{r+1}\cdot\frac{p}{q} = \frac{12 - r}{r+1}\cdot\frac{1}{5} = \frac{12 - r}{5(r+1)}. \]
\item The probabilities increase while $\dfrac{12 - r}{5(r+1)} > 1$, i.e. $12 - r > 5r + 5$, i.e. $7 > 6r$, i.e. $r < \tfrac76$. So the probabilities increase for $r = 0, 1$ (up to $r+1 = 2$) and decrease afterwards. Alternatively, $(n+1)p = 13 \times \tfrac16 = \tfrac{13}{6} \approx 2.17$, so the mode is $m = \lfloor 2.17 \rfloor = 2$.
\item Check:
\begin{align*}
P(X = 1) &= \binom{12}{1}\left(\tfrac16\right)\left(\tfrac56\right)^{11} = 12 \times \tfrac16 \times 0.1346 \approx 0.2692,\\
P(X = 2) &= \binom{12}{2}\left(\tfrac16\right)^{2}\left(\tfrac56\right)^{10} = 66 \times \tfrac{1}{36} \times 0.1615 \approx 0.2961,\\
P(X = 3) &= \binom{12}{3}\left(\tfrac16\right)^{3}\left(\tfrac56\right)^{9} = 220 \times \tfrac{1}{216} \times 0.1938 \approx 0.1974.
\end{align*}
Indeed $P(X=2) > P(X=1)$ and $P(X=2) > P(X=3)$, confirming that the mode is $\boxed{2}$: the most likely outcome is exactly two sixes.
\end{enumerate}
\end{exoresolu}

\begin{aretenir}[Key facts for the Binomial Distribution]
\begin{itemize}
\item Conditions: fixed $n$, two outcomes, constant $p$, independence.
\item $X \sim B(n,p) \Rightarrow P(X=r) = \binom{n}{r}p^r q^{n-r}$.
\item $E(X) = np$, $\operatorname{Var}(X) = npq$.
\item Cumulative tables give $P(X \leq r)$; convert all questions to this form.
\item Mode $m = \lfloor (n+1)p \rfloor$; two modes if $(n+1)p \in \mathbb{N}$.
\item Symmetry: if $p > 0.5$, use $Y = n-X \sim B(n, 1-p)$.
\end{itemize}
\end{aretenir}

\courssub{The Discrete Uniform Distribution}

\begin{definition}[Discrete Uniform Distribution]
A discrete random variable $X$ follows a \cle{uniform distribution} on the set $\{x_1, x_2, \ldots, x_n\}$ if all outcomes are equally likely:
\[ P(X = x_i) = \frac{1}{n}, \qquad i = 1, \ldots, n. \]
The standard case is $X$ taking values $1, 2, \ldots, n$.
\end{definition}

\begin{propriete}[Mean and Variance of the Uniform Distribution on $\{1, 2, \ldots, n\}$]
\[ E(X) = \frac{n+1}{2}, \qquad \operatorname{Var}(X) = \frac{n^{2} - 1}{12}. \]
\textbf{Derivation (examinable):}
\begin{align*}
E(X) &= \frac{1}{n}\sum_{r=1}^{n} r = \frac{1}{n}\cdot\frac{n(n+1)}{2} = \frac{n+1}{2},\\[6pt]
E(X^{2}) &= \frac{1}{n}\sum_{r=1}^{n} r^{2} = \frac{1}{n}\cdot\frac{n(n+1)(2n+1)}{6} = \frac{(n+1)(2n+1)}{6},\\[6pt]
\operatorname{Var}(X) &= E(X^{2}) - [E(X)]^{2} = \frac{(n+1)(2n+1)}{6} - \frac{(n+1)^{2}}{4}\\
&= \frac{n+1}{12}\bigl[2(2n+1) - 3(n+1)\bigr] = \frac{n+1}{12}(n-1) = \frac{n^{2}-1}{12}.
\end{align*}
For a uniform distribution on $\{a, a+1, \ldots, b\}$ with $n = b-a+1$ values:
\[ E(X) = \frac{a+b}{2}, \qquad \operatorname{Var}(X) = \frac{n^{2} - 1}{12} \quad \text{(variance unchanged by translation)}. \]
\end{propriete}

\begin{methode}[Using the uniform distribution on $\{1, 2, \ldots, n\}$]
\begin{enumerate}
\item Recognise the model: $n$ equally likely outcomes, $P(X = r) = \dfrac{1}{n}$ for $r = 1, \ldots, n$.
\item Quote the key results (derivation examinable):
\[ E(X) = \frac{n+1}{2}, \qquad \operatorname{Var}(X) = \frac{n^{2} - 1}{12}. \]
\item For outcomes $a, a+1, \ldots, b$ (with $n = b - a + 1$ values), shift:
$E(X) = \dfrac{a+b}{2}$ and the variance is unchanged, $\operatorname{Var}(X) = \dfrac{n^2 - 1}{12}$.
\end{enumerate}
\textbf{Key identities for the derivation:} $\sum_{r=1}^{n} r = \dfrac{n(n+1)}{2}$ and $\sum_{r=1}^{n} r^{2} = \dfrac{n(n+1)(2n+1)}{6}$.
\end{methode}

\begin{exoresolu}[Deriving the variance and applying it]
\begin{enumerate}
\item Prove that if $X$ is uniformly distributed on $\{1, 2, \ldots, n\}$, then $\operatorname{Var}(X) = \dfrac{n^{2}-1}{12}$.
\item A fair tetrahedral die has faces numbered $1, 2, 3, 4$. Find the mean and variance of the score.
\end{enumerate}
\tcblower
\textbf{Solution.}
\begin{enumerate}
\item First, $E(X) = \dfrac{1}{n}\sum_{r=1}^{n} r = \dfrac{1}{n}\cdot\dfrac{n(n+1)}{2} = \dfrac{n+1}{2}$.
Next,
\[ E(X^{2}) = \frac{1}{n}\sum_{r=1}^{n} r^{2} = \frac{1}{n}\cdot\frac{n(n+1)(2n+1)}{6} = \frac{(n+1)(2n+1)}{6}. \]
Then
\begin{align*}
\operatorname{Var}(X) &= E(X^{2}) - [E(X)]^{2} = \frac{(n+1)(2n+1)}{6} - \frac{(n+1)^{2}}{4}\\
&= \frac{n+1}{12}\Big[2(2n+1) - 3(n+1)\Big] = \frac{n+1}{12}(4n + 2 - 3n - 3)\\
&= \frac{(n+1)(n-1)}{12} = \frac{n^{2} - 1}{12}.
\end{align*}
\item Here $n = 4$: $E(X) = \dfrac{4+1}{2} = 2.5$ and $\operatorname{Var}(X) = \dfrac{16 - 1}{12} = \dfrac{15}{12} = 1.25$.
\end{enumerate}
\end{exoresolu}

\begin{savaistu}[Historical note]
The discrete uniform distribution is the simplest model of ``fairness''. It underlies the classical definition of probability (Laplace, 1812): if all outcomes are equally likely, the probability of an event is the ratio of favourable cases to total cases. In Cameroon, the traditional game of \emph{Nsolo} (or \emph{Awale}) uses seeds and pits; the initial distribution of seeds can be modelled using uniform assumptions when players choose pits at random.
\end{savaistu}

\courssub{The Geometric Distribution}

\begin{definition}[Geometric Distribution]
The \cle{geometric distribution} models the number of independent Bernoulli trials (success probability $p$, failure probability $q=1-p$) needed to obtain the \textbf{first success}. We write $X \sim \mathrm{Geo}(p)$.
\[ P(X = r) = q^{\,r-1} p, \qquad r = 1, 2, 3, \ldots \]
\textbf{Note:} Some textbooks define the geometric distribution as the number of \emph{failures} before the first success (support $0,1,2,\ldots$). The GCE syllabus uses the ``number of trials until first success'' definition (support $1,2,3,\ldots$). Always check which definition is used.
\end{definition}

\begin{propriete}[Mean, Variance and Tail Probabilities of $\mathrm{Geo}(p)$]
For $X \sim \mathrm{Geo}(p)$ with $q=1-p$:
\begin{align*}
E(X) &= \frac{1}{p},\\[4pt]
\operatorname{Var}(X) &= \frac{q}{p^{2}},\\[4pt]
P(X > r) &= q^{r} \quad \text{(first $r$ trials all fail)},\\[4pt]
P(X \leq r) &= 1 - q^{r}.
\end{align*}
The mode is always $1$ (the distribution decreases geometrically).
\textbf{Derivation of $E(X)$ (examinable idea):}
$E(X) = \sum_{r=1}^{\infty} r q^{r-1}p = p \frac{d}{dq}\sum_{r=1}^{\infty} q^{r} = p \frac{d}{dq}\left(\frac{q}{1-q}\right) = \frac{p}{(1-q)^2} = \frac{1}{p}$.
\end{propriete}

\begin{methode}[Modelling ``waiting time until the first success'']
\begin{enumerate}
\item Use the geometric distribution when independent trials with constant success probability $p$ are repeated until the \textbf{first} success; $X$ = number of trials needed.
\item Probabilities: $P(X = r) = q^{\,r-1} p$, $r = 1, 2, 3, \ldots$
\item Tail probabilities (very useful): $P(X > r) = q^{r}$ (the first $r$ trials all fail).
\item Mean and variance:
\[ E(X) = \frac{1}{p}, \qquad \operatorname{Var}(X) = \frac{q}{p^{2}}. \]
\item The mode is always $1$: the distribution decreases geometrically.
\end{enumerate}
\textbf{Reflex:} check whether the question counts trials \emph{up to and including} the first success (geometric) or successes in a \emph{fixed} number of trials (binomial).
\end{methode}

\begin{exoresolu}[First goal scored]
A striker scores with probability $0.2$ on each penalty kick, independently of previous kicks. Let $X$ be the number of kicks taken until he scores his first goal.
\begin{enumerate}
\item Write down the distribution of $X$ and compute $P(X = 4)$.
\item Find the probability that he needs more than $3$ kicks.
\item Find $E(X)$ and $\operatorname{Var}(X)$, and interpret the mean.
\item Find the smallest $n$ such that the probability that he scores within $n$ kicks exceeds $0.9$.
\end{enumerate}
\tcblower
\textbf{Solution.}
\begin{enumerate}
\item $X \sim \mathrm{Geo}(0.2)$ with $p = 0.2$, $q = 0.8$.
\[ P(X = 4) = q^{3}p = (0.8)^{3}(0.2) = 0.512 \times 0.2 = 0.1024. \]
\item $P(X > 3) = q^{3} = (0.8)^{3} = 0.512$.
\item $E(X) = \dfrac{1}{p} = \dfrac{1}{0.2} = 5$ kicks on average; $\operatorname{Var}(X) = \dfrac{q}{p^{2}} = \dfrac{0.8}{0.04} = 20$.
Interpretation: over many sequences of kicks, the average number of kicks needed to score the first goal is $5$.
\item $P(X \leq n) = 1 - P(X > n) = 1 - (0.8)^{n} > 0.9 \iff (0.8)^{n} < 0.1$.
Taking logarithms: $n \ln 0.8 < \ln 0.1$, and since $\ln 0.8 < 0$,
\[ n > \frac{\ln 0.1}{\ln 0.8} = \frac{-2.3026}{-0.2231} \approx 10.32. \]
The smallest integer is $n = 11$: he needs at most $11$ kicks to score with probability greater than $0.9$.
\end{enumerate}
\end{exoresolu}

\begin{attention}[Geometric vs. Binomial]
\begin{itemize}
\item \textbf{Binomial}: fixed number of trials $n$, count successes. $X \sim B(n,p)$.
\item \textbf{Geometric}: no fixed $n$, repeat until first success, count trials. $X \sim \mathrm{Geo}(p)$.
\item If a question says ``find the probability that the first success occurs on the 5th trial'', it is geometric. If it says ``in 10 trials, find the probability of exactly 3 successes'', it is binomial.
\end{itemize}
\end{attention}

\courssub{The Poisson Distribution and Its Tables}

\begin{definition}[Poisson Distribution]
The \cle{Poisson distribution} models the number of events occurring \textbf{randomly}, \textbf{independently}, at a \textbf{constant average rate} in a fixed interval of time or space. We write $X \sim \mathrm{Po}(\lambda)$ where $\lambda > 0$ is the mean number of events in the interval.
\[ P(X = r) = e^{-\lambda}\,\frac{\lambda^{r}}{r!}, \qquad r = 0, 1, 2, \ldots \]
\end{definition}

\begin{propriete}[Mean, Variance and Additive Property of $\mathrm{Po}(\lambda)$]
For $X \sim \mathrm{Po}(\lambda)$:
\[ E(X) = \operatorname{Var}(X) = \lambda. \]
\textbf{Derivation of $E(X)$ (examinable):}
$E(X) = \sum_{r=0}^{\infty} r e^{-\lambda}\frac{\lambda^r}{r!} = e^{-\lambda}\lambda \sum_{r=1}^{\infty} \frac{\lambda^{r-1}}{(r-1)!} = \lambda e^{-\lambda} e^{\lambda} = \lambda$.
\textbf{Additive property:} If $X \sim \mathrm{Po}(\lambda_1)$ and $Y \sim \mathrm{Po}(\lambda_2)$ are \textbf{independent}, then
\[ X + Y \sim \mathrm{Po}(\lambda_1 + \lambda_2). \]
\textbf{Proof (using MGF or convolution):} $P(X+Y = n) = \sum_{k=0}^{n} P(X=k)P(Y=n-k) = e^{-(\lambda_1+\lambda_2)}\frac{(\lambda_1+\lambda_2)^n}{n!}$.
\end{propriete}

\begin{methode}[Modelling random events in a fixed interval]
\begin{enumerate}
\item Use the Poisson distribution for events occurring randomly, independently, at a constant average rate in a fixed interval: $X \sim \mathrm{Po}(\lambda)$,
\[ P(X = r) = e^{-\lambda}\,\frac{\lambda^{r}}{r!}, \qquad E(X) = \operatorname{Var}(X) = \lambda. \]
\item \textbf{Unit interval reflex:} if the rate is given per unit but the question concerns a different interval, rescale $\lambda$ proportionally (rate $\lambda$ per minute $\Rightarrow$ mean $5\lambda$ in $5$ minutes).
\item Cumulative tables give $P(X \leq r)$; manipulate exactly as for binomial tables:
$P(X = r) = P(X \leq r) - P(X \leq r-1)$,
$P(X \geq r) = 1 - P(X \leq r-1)$.
\item If $X \sim \mathrm{Po}(\lambda_1)$ and $Y \sim \mathrm{Po}(\lambda_2)$ are \textbf{independent}, then $X + Y \sim \mathrm{Po}(\lambda_1 + \lambda_2)$.
\end{enumerate}
\end{methode}

\begin{exoresolu}[Calls at a call centre in Yaoundé]
A customer-service centre in Yaoundé receives calls randomly at an average rate of $3$ calls per $10$ minutes.
\begin{enumerate}
\item Find the probability of exactly $2$ calls in a $10$-minute period.
\item Using the cumulative table extract below, find the probability of at least $4$ calls in a $10$-minute period.
\item Find the probability of exactly $5$ calls in a $20$-minute period.
\item Independently, the centre receives emails at an average rate of $1.5$ per $10$ minutes. Find the probability that the total number of contacts (calls plus emails) in $10$ minutes is exactly $3$.
\end{enumerate}
Table extract for $\lambda = 3$: $P(X \leq 2) = 0.4232$, $P(X \leq 3) = 0.6472$.
\tcblower
\textbf{Solution.}
\begin{enumerate}
\item With $\lambda = 3$: $P(X = 2) = e^{-3}\dfrac{3^{2}}{2!} = e^{-3} \times 4.5 \approx 0.04979 \times 4.5 \approx 0.2240$.
\item $P(X \geq 4) = 1 - P(X \leq 3) = 1 - 0.6472 = 0.3528$.
\item New interval, new mean: in $20$ minutes, $\lambda = 2 \times 3 = 6$.
\[ P(X = 5) = e^{-6}\,\frac{6^{5}}{5!} = e^{-6} \times \frac{7776}{120} = e^{-6} \times 64.8 \approx 0.002479 \times 64.8 \approx 0.1606. \]
\item Calls $X \sim \mathrm{Po}(3)$, emails $Y \sim \mathrm{Po}(1.5)$, independent, so $T = X + Y \sim \mathrm{Po}(4.5)$.
\[ P(T = 3) = e^{-4.5}\,\frac{4.5^{3}}{3!} = e^{-4.5} \times \frac{91.125}{6} = e^{-4.5} \times 15.1875 \approx 0.01111 \times 15.1875 \approx 0.1687. \]
\end{enumerate}
\end{exoresolu}

\begin{savaistu}[Siméon Denis Poisson (1781--1840)]
Poisson was a French mathematician who introduced this distribution in 1837 in his work \emph{Recherches sur la probabilité des jugements en matière criminelle et en matière civile}. He derived it as a limit of the binomial distribution when $n \to \infty$, $p \to 0$ with $np = \lambda$ constant. The distribution is now ubiquitous in queueing theory, telecommunications (e.g., modelling call arrivals at MTN or Orange call centres in Cameroon), and reliability engineering.
\end{savaistu}

\courssub{Poisson Approximation to the Binomial}

\begin{propriete}[Poisson Approximation to $B(n,p)$]
If $n$ is \textbf{large} (typically $n \geq 50$), $p$ is \textbf{small} (typically $p \leq 0.1$), and $np = \lambda$ is moderate (say $\lambda \leq 5$), then
\[ B(n, p) \approx \mathrm{Po}(\lambda) \quad \text{with } \lambda = np. \]
\textbf{Justification (examinable idea):} For fixed $\lambda = np$, as $n \to \infty$,
\[ \binom{n}{r} p^{r} q^{\,n-r} = \frac{n(n-1)\cdots(n-r+1)}{r!} \left(\frac{\lambda}{n}\right)^{r} \left(1 - \frac{\lambda}{n}\right)^{n-r} \longrightarrow e^{-\lambda}\frac{\lambda^{r}}{r!}. \]
The approximation works because both distributions then have mean $\lambda$ and variance $\approx \lambda$ (since $npq = \lambda(1-p) \approx \lambda$).
\end{propriete}

\begin{methode}[Approximating $B(n,p)$ by $\mathrm{Po}(np)$]
\begin{enumerate}
\item Check the conditions: $n$ \textbf{large} (typically $n \geq 50$) and $p$ \textbf{small} (typically $p \leq 0.1$), so that $np$ is moderate.
\item Set $\lambda = np$ and replace $X \sim B(n,p)$ by $X \approx \mathrm{Po}(\lambda)$.
\item Compute with the Poisson formula or tables; state the approximation clearly in your answer.
\item Justification (examinable idea): for fixed $\lambda = np$, as $n \to \infty$,
\[ \binom{n}{r} p^{r} q^{\,n-r} \longrightarrow e^{-\lambda}\frac{\lambda^{r}}{r!}. \]
\end{enumerate}
\textbf{Reflex:} the approximation works because both distributions then have mean $\lambda$ and variance $\approx \lambda$ (since $npq = \lambda(1-p) \approx \lambda$).
\end{methode}

\begin{exoresolu}[Rare blood group]
In a certain region, $2\%$ of people have a particular blood group. A mobile clinic screens $150$ people. Let $X$ be the number with this blood group.
\begin{enumerate}
\item Explain why $X$ may be approximated by a Poisson distribution and state its parameter.
\item Use the approximation to find $P(X = 4)$ and $P(X \leq 2)$.
\item Compare $P(X = 4)$ with the exact binomial value.
\end{enumerate}
\tcblower
\textbf{Solution.}
\begin{enumerate}
\item $n = 150$ is large and $p = 0.02$ is small, with $np = 3$ moderate, so the Poisson approximation is appropriate: $X \approx \mathrm{Po}(3)$.
\item 
\[ P(X = 4) \approx e^{-3}\frac{3^{4}}{4!} = e^{-3} \times \frac{81}{24} = e^{-3} \times 3.375 \approx 0.04979 \times 3.375 \approx 0.1680. \]
For the cumulative probability:
\begin{align*}
P(X \leq 2) &\approx e^{-3}\left(1 + 3 + \frac{9}{2}\right) = e^{-3} \times 8.5 \approx 0.04979 \times 8.5 \approx 0.4232.
\end{align*}
\item The exact binomial value is $\binom{150}{4}(0.02)^{4}(0.98)^{146} \approx 0.1693$. The Poisson estimate $0.1680$ differs by only about $0.0013$, a relative error under $1\%$: the approximation is excellent.
\end{enumerate}
\begin{attention}
Never use the Poisson approximation when $p$ is large (e.g. $p = 0.5$): the variance $npq$ then differs greatly from $np$ and the approximation fails. For large $p$, apply the trick to the \emph{failures} (which have small probability $q$) instead.
\end{attention}
\end{exoresolu}

\courssub{Linking Probability Distributions to Frequency Distributions}

\begin{methode}[Fitting a theoretical distribution to observed data]
\begin{enumerate}
\item From the data, estimate the parameter: the sample mean $\bar{x}$ estimates $np$ (binomial) or $\lambda$ (Poisson).
\item Compute the theoretical probabilities $P(X = r)$ for each value $r$.
\item Multiply by the total frequency $N$ to obtain the \cle{expected frequencies} $E_r = N \times P(X = r)$.
\item Compare observed and expected frequencies (e.g. by inspection or a goodness-of-fit idea) and conclude whether the model is suitable.
\end{enumerate}
\end{methode}

\begin{exoresolu}[Fitting a Poisson distribution to accident data]
The number of accidents per week at a junction in Bamenda was recorded over $100$ weeks:
\begin{center}
\begin{tabular}{|l|c|c|c|c|c|}
\hline
\rowcolor{popblueL} Accidents per week & $0$ & $1$ & $2$ & $3$ & $4$ or more \\
\hline
Number of weeks & $25$ & $38$ & $24$ & $10$ & $3$ \\
\hline
\end{tabular}
\end{center}
The mean number of accidents per week is $\bar{x} = 1.3$.
\begin{enumerate}
\item Explain why a Poisson model with $\lambda = 1.3$ might be suitable.
\item Calculate the expected frequencies for $0$, $1$ and $2$ accidents per week (to $1$ d.p.).
\item Comment on the fit.
\end{enumerate}
\tcblower
\textbf{Solution.}
\begin{enumerate}
\item Accidents occur randomly, independently and (presumably) at a roughly constant average rate over the weeks; the Poisson distribution is the natural model, with $\lambda$ estimated by the sample mean $1.3$.
\item With $e^{-1.3} \approx 0.2725$:
\begin{align*}
E_0 &= 100 \times e^{-1.3} = 100 \times 0.2725 = 27.3,\\
E_1 &= 100 \times e^{-1.3} \times 1.3 = 27.3 \times 1.3 = 35.4,\\
E_2 &= 100 \times e^{-1.3} \times \frac{1.3^{2}}{2} = 35.4 \times \frac{1.3}{2} = 23.0.
\end{align*}
\item Comparing: observed $25, 38, 24$ against expected $27.3, 35.4, 23.0$. The agreement is close (differences of only $1$ to $3$ weeks out of $100$), so the Poisson model with $\lambda = 1.3$ fits the data well. This illustrates the general principle: a probability distribution is the theoretical limit of a frequency distribution as the number of observations grows.
\end{enumerate}
\end{exoresolu}

\begin{aretenir}[Choosing the right discrete distribution]
\begin{center}
\begin{tabular}{|l|l|l|}
\hline
\rowcolor{popblueL} \textbf{Situation} & \textbf{Distribution} & \textbf{Key parameter(s)} \\
\hline
Fixed $n$ trials, count successes & Binomial $B(n,p)$ & $n$, $p$ \\
\hline
Wait for first success & Geometric $\mathrm{Geo}(p)$ & $p$ \\
\hline
Events in fixed interval, constant rate & Poisson $\mathrm{Po}(\lambda)$ & $\lambda$ (mean rate $\times$ interval) \\
\hline
Equally likely outcomes $1,\ldots,n$ & Uniform & $n$ \\
\hline
\end{tabular}
\end{center}
\textbf{Approximations:} $B(n,p) \approx \mathrm{Po}(np)$ if $n \ge 50$, $p \le 0.1$.
\textbf{Sum of independent Poissons:} $\mathrm{Po}(\lambda_1) + \mathrm{Po}(\lambda_2) = \mathrm{Po}(\lambda_1+\lambda_2)$.
\end{aretenir}

\begin{exercice}[Mixed practice]
\begin{enumerate}
\item A fair coin is tossed $10$ times. Find the probability of obtaining exactly $6$ heads. State the mean and variance of the number of heads.
\item In a large batch of components, $3\%$ are defective. A sample of $50$ is taken. Use a Poisson approximation to find the probability that the sample contains at most $2$ defectives.
\item The number of emails received per hour by a server follows a Poisson distribution with mean $4$. Find the probability that in a $30$-minute period the server receives exactly $3$ emails.
\item A random variable $X$ is uniformly distributed on the integers $-2, -1, 0, 1, 2$. Find $E(X)$ and $\operatorname{Var}(X)$.
\item A basketball player makes $70\%$ of her free throws. She shoots until she misses. Find the probability that she makes exactly $4$ shots before missing. (Careful: define the success event.)
\end{enumerate}
\end{exercice}

\begin{corrige}[Mixed practice]
\begin{enumerate}
\item $X \sim B(10, 0.5)$. $P(X=6) = \binom{10}{6}0.5^{10} = 210/1024 \approx 0.2051$. $E(X)=5$, $\operatorname{Var}(X)=2.5$.
\item $n=50$, $p=0.03$, $np=1.5$. $X \approx \mathrm{Po}(1.5)$. $P(X \leq 2) = e^{-1.5}(1 + 1.5 + 1.5^2/2) = e^{-1.5} \times 3.625 \approx 0.2231 \times 3.625 \approx 0.8088$.
\item $30$ min $\Rightarrow \lambda = 2$. $P(X=3) = e^{-2} \times 8/6 = 4e^{-2}/3 \approx 0.1804$.
\item Values: $-2,-1,0,1,2$ ($n=5$). Shift of standard uniform on $1..5$ by $-3$. $E(X) = (1+2+3+4+5)/5 - 3 = 3 - 3 = 0$. $\operatorname{Var}(X) = (5^2-1)/12 = 24/12 = 2$.
\item ``Makes exactly 4 shots before missing'' $\Rightarrow$ first miss on 5th trial. Define success = miss ($p=0.3$). $X \sim \mathrm{Geo}(0.3)$. $P(X=5) = 0.7^4 \times 0.3 = 0.2401 \times 0.3 = 0.07203$.
\end{enumerate}
\end{corrige}

\newpage

\coursec{Exercises}

\courssub{I check my knowledge}

\begin{exercice}[MCQ: spotting the right distribution]
For each situation, choose the correct distribution for $X$ from: binomial, discrete uniform, geometric, Poisson. Justify briefly.
\begin{enumerate}
\item $X$ = number of heads in 20 tosses of a fair coin.
\item $X$ = number of the ticket drawn at random from tickets numbered $1$ to $50$.
\item $X$ = number of attempts needed to score a first goal, each attempt succeeding with probability $0.3$.
\item $X$ = number of customers arriving at a bank counter in one hour, arrivals occurring randomly at an average rate of $12$ per hour.
\end{enumerate}
\end{exercice}

\begin{exercice}[True or false?]
State whether each assertion is true or false, justifying your answer.
\begin{enumerate}
\item If $X \sim B(n, p)$, then $\mathrm{Var}(X) = np$.
\item If $X \sim \mathrm{Po}(\lambda)$, then $E(X) = \mathrm{Var}(X)$.
\item If $X$ is uniformly distributed on $\{1, 2, \dots, n\}$, then $E(X) = \dfrac{n+1}{2}$.
\item If $X \sim \mathrm{Po}(2)$ and $Y \sim \mathrm{Po}(3)$ are independent, then $X + Y \sim \mathrm{Po}(6)$.
\end{enumerate}
\end{exercice}

\begin{exercice}[Key definitions and formulas]
\begin{enumerate}
\item State the four conditions (Bernoulli scheme) under which a random variable follows a binomial distribution.
\item Write down the probability mass function of $X \sim B(n, p)$ and of $Y \sim \mathrm{Po}(\lambda)$.
\item State the recurrence relation linking $P(X = k+1)$ and $P(X = k)$ for $X \sim B(n, p)$, and explain how it is used to locate the mode.
\item State the mean and variance of the geometric distribution with parameter $p$ (number of trials up to and including the first success).
\end{enumerate}
\end{exercice}

\begin{exercice}[Direct calculations]
\begin{enumerate}
\item A fair die is thrown $12$ times and $X$ is the number of sixes obtained. Write down the distribution of $X$ and compute $E(X)$ and $\mathrm{Var}(X)$.
\item Compute $P(X = 2)$, giving your answer to 4 decimal places.
\item $Y \sim \mathrm{Po}(4)$. Compute $P(Y = 0)$ and $P(Y = 3)$ to 4 decimal places.
\item $T$ is uniformly distributed on $\{1, 2, \dots, 10\}$. Find $E(T)$ and $\mathrm{Var}(T)$.
\end{enumerate}
\end{exercice}

\courssub{I practise}

\begin{exercice}[Binomial: parameters and mode from the moments]
Given $X \sim B(n, p)$ with $E(X) = 6$ and $\mathrm{Var}(X) = 4.2$:
\begin{enumerate}
\item Find the values of $p$ and $n$.
\item Using the recurrence formula, determine the mode of $X$.
\item Compute the modal probability $P(X = m)$, where $m$ is the mode, to 4 decimal places.
\end{enumerate}
\end{exercice}

\begin{exercice}[Uniform distribution from first principles]
A box contains tickets numbered $1$ to $20$. One ticket is drawn at random and $X$ is the number on the ticket.
\begin{enumerate}
\item Show, from first principles, that $E(X) = 10.5$.
\item Show, from first principles, that $\mathrm{Var}(X) = 33.25$. (You may use $\sum_{k=1}^{n} k^2 = \dfrac{n(n+1)(2n+1)}{6}$.)
\item Find $P(X > 12)$ and $P(X \text{ is even})$.
\end{enumerate}
\end{exercice}

\begin{exercice}[Geometric distribution: the driving test]
The probability that a candidate passes a driving test at each attempt is $0.4$, independently of previous attempts. Let $X$ be the number of attempts needed to pass for the first time.
\begin{enumerate}
\item Find the probability that she passes at the third attempt.
\item Find the probability that she needs more than $4$ attempts.
\item Find the expected number of attempts and the variance of $X$.
\item Find the smallest number $n$ of attempts such that the probability of passing within $n$ attempts exceeds $0.95$.
\end{enumerate}
\end{exercice}

\begin{exercice}[Poisson: calls at an office and sum of independent variables]
Calls arrive at an office at random, at an average rate of $5$ per hour.
\begin{enumerate}
\item Find the probability of exactly $3$ calls in one hour.
\item Find the probability of no call in a $15$-minute interval.
\item Find the probability of at least one call in a $10$-minute interval.
\item Independently, emails arrive at an average rate of $3$ per hour. Find the probability that the total number of calls and emails received in one hour is exactly $4$, and the probability that this total is at least $2$.
\end{enumerate}
\end{exercice}

\courssub{I prepare for the GCE}

\begin{exercice}[GCE-style: binomial tables and a critical value]
A biased coin has probability $0.3$ of showing a head. It is tossed $8$ times and $X$ denotes the number of heads obtained.
\begin{enumerate}
\item State the distribution of $X$ and write down $E(X)$ and $\mathrm{Var}(X)$. \hfill (3 marks)
\item Using cumulative binomial tables, find $P(X \geq 2)$. \hfill (2 marks)
\item Using cumulative binomial tables, find $P(1 \leq X \leq 4)$. \hfill (3 marks)
\item Using tables, or otherwise, find $P(X = 5)$. \hfill (2 marks)
\item Find the smallest integer $r$ such that $P(X \leq r) > 0.9$. \hfill (2 marks)
\item Using the recurrence formula, determine the mode of $X$ and compute the corresponding probability. \hfill (4 marks)
\end{enumerate}
\end{exercice}

\begin{exercice}[GCE-style: Poisson approximation to the binomial]
A factory in Douala produces $400$ items per day, of which $1.5\%$ are defective, independently of one another. Let $X$ be the number of defective items produced in a day.
\begin{enumerate}
\item State the exact distribution of $X$, and compute $E(X)$ and $\mathrm{Var}(X)$. \hfill (3 marks)
\item Explain why $X$ may be approximated by a Poisson distribution, and state the parameter of this approximation. \hfill (2 marks)
\item Using this approximation, find the probability that exactly $6$ defective items are produced in a day. \hfill (2 marks)
\item Using cumulative Poisson tables, find the probability that more than $8$ defective items are produced in a day. \hfill (3 marks)
\item Find the probability that, in a day, the number of defective items lies between $3$ and $9$ inclusive. \hfill (3 marks)
\item On a randomly chosen day, the number of defective items exceeded $10$. Comment on whether this casts doubt on the claimed defect rate of $1.5\%$. \hfill (2 marks)
\end{enumerate}
\end{exercice}

\begin{exercice}[GCE-style: from a frequency distribution to a Poisson model]
The table below gives the observed frequencies of the number of goals scored per match in $50$ football matches of the national championship.
\begin{center}
\begin{tabular}{|l|cccccc|}
\hline
\rowcolor{popblueL} Number of goals $x$ & $0$ & $1$ & $2$ & $3$ & $4$ & $5$ or more \\
\hline Observed frequency & $8$ & $15$ & $13$ & $9$ & $4$ & $1$ \\
\hline
\end{tabular}
\end{center}
\begin{enumerate}
\item Calculate the sample mean number of goals per match. \hfill (3 marks)
\item Explain why a Poisson distribution may be a suitable model, and estimate its parameter $\lambda$. \hfill (2 marks)
\item Using $\lambda$ equal to your estimate (rounded to 1 decimal place), compute $P(X = 0)$, $P(X = 1)$ and $P(X = 2)$, and hence the expected frequencies for $0$, $1$ and $2$ goals per match. \hfill (5 marks)
\item Compare the expected frequencies with the observed frequencies and comment on the goodness of fit of the Poisson model. \hfill (3 marks)
\item Using the model, find the probability that a randomly chosen match produces at least $3$ goals. \hfill (2 marks)
\end{enumerate}
\end{exercice}

\newpage

\coursec{Solutions to the exercises}

\courssub{Solutions to "I check my knowledge"}

\begin{corrige}[Exercise 1 — MCQ: spotting the right distribution]
\begin{enumerate}
\item \textbf{Binomial:} $X \sim B(20, 0.5)$. There is a fixed number $n = 20$ of independent trials, each with two outcomes (head/tail) and constant success probability $p = 0.5$.
\item \textbf{Discrete uniform} on $\{1, 2, \dots, 50\}$: each ticket has the same probability $\frac{1}{50}$ of being drawn, so all outcomes are equally likely.
\item \textbf{Geometric} with parameter $p = 0.3$: we repeat independent identical trials until the \emph{first} success (first goal), and $X$ counts the number of trials needed.
\item \textbf{Poisson} with parameter $\lambda = 12$: events (arrivals) occur randomly, independently, at a constant average rate in a fixed interval of time (one hour).
\end{enumerate}
\end{corrige}

\begin{corrige}[Exercise 2 — True or false?]
\begin{enumerate}
\item \textbf{False.} If $X \sim B(n, p)$ then $\mathrm{Var}(X) = np(1-p) = npq$, not $np$. It is $E(X)$ that equals $np$.
\item \textbf{True.} For $X \sim \mathrm{Po}(\lambda)$, $E(X) = \lambda$ and $\mathrm{Var}(X) = \lambda$; equality of mean and variance is the signature of the Poisson distribution.
\item \textbf{True.} $E(X) = \dfrac{1}{n}\displaystyle\sum_{k=1}^{n} k = \dfrac{1}{n}\cdot\dfrac{n(n+1)}{2} = \dfrac{n+1}{2}$.
\item \textbf{False.} The sum of two independent Poisson variables is Poisson with parameter equal to the \emph{sum} of the parameters: $X + Y \sim \mathrm{Po}(2+3) = \mathrm{Po}(5)$, not $\mathrm{Po}(6)$.
\end{enumerate}
\end{corrige}

\begin{corrige}[Exercise 3 — Key definitions and formulas]
\begin{enumerate}
\item The four conditions of the Bernoulli scheme are:
\begin{itemize}
\item a \textbf{fixed number} $n$ of trials;
\item each trial has exactly \textbf{two outcomes} (success/failure);
\item the probability of success $p$ is \textbf{constant} from trial to trial;
\item the trials are \textbf{independent}.
\end{itemize}
Then $X$ = number of successes follows $B(n, p)$.
\item For $X \sim B(n, p)$: $\;P(X = k) = \dbinom{n}{k} p^k (1-p)^{n-k}$, $k = 0, 1, \dots, n$.\par
For $Y \sim \mathrm{Po}(\lambda)$: $\;P(Y = k) = e^{-\lambda}\dfrac{\lambda^k}{k!}$, $k = 0, 1, 2, \dots$
\item The recurrence relation is
\[ \frac{P(X = k+1)}{P(X = k)} = \frac{n - k}{k + 1}\cdot\frac{p}{1-p}. \]
The probabilities increase while $\dfrac{P(X=k+1)}{P(X=k)} > 1$, i.e. while $k < (n+1)p - 1$, and decrease afterwards. Hence the \cle{mode} is $m = \lfloor (n+1)p \rfloor$ when $(n+1)p$ is not an integer; if $(n+1)p$ is an integer, there are two modes, $(n+1)p - 1$ and $(n+1)p$.
\item For the geometric distribution (number of trials up to and including the first success):
\[ E(X) = \frac{1}{p}, \qquad \mathrm{Var}(X) = \frac{1 - p}{p^2} = \frac{q}{p^2}. \]
\end{enumerate}
\end{corrige}

\begin{corrige}[Exercise 4 — Direct calculations]
\begin{enumerate}
\item Each throw is a Bernoulli trial with $p = P(\text{six}) = \frac{1}{6}$, repeated $n = 12$ times independently, so $X \sim B\!\left(12, \tfrac{1}{6}\right)$.
\[ E(X) = np = 12 \times \tfrac{1}{6} = 2, \qquad \mathrm{Var}(X) = npq = 12 \times \tfrac{1}{6} \times \tfrac{5}{6} = \frac{5}{3} \approx 1.6667. \]
\item $P(X = 2) = \dbinom{12}{2}\left(\dfrac{1}{6}\right)^{2}\left(\dfrac{5}{6}\right)^{10} = 66 \times \dfrac{1}{36} \times 0.16151 \approx \mathbf{0.2961}$.
\item $Y \sim \mathrm{Po}(4)$:
\[ P(Y = 0) = e^{-4} \approx \mathbf{0.0183}; \qquad P(Y = 3) = e^{-4}\,\frac{4^3}{3!} = e^{-4}\times\frac{64}{6} \approx \mathbf{0.1954}. \]
\item $T$ uniform on $\{1, \dots, 10\}$:
\[ E(T) = \frac{n+1}{2} = \frac{11}{2} = 5.5, \qquad \mathrm{Var}(T) = \frac{n^2 - 1}{12} = \frac{100 - 1}{12} = \frac{99}{12} = 8.25. \]
\end{enumerate}
\end{corrige}

\courssub{Solutions to "I practise"}

\begin{corrige}[Exercise 5 — Binomial: parameters and mode from the moments]
\begin{enumerate}
\item We have $np = 6$ and $npq = 4.2$. Dividing: $q = \dfrac{npq}{np} = \dfrac{4.2}{6} = 0.7$, so $p = 1 - q = 0.3$ and $n = \dfrac{6}{0.3} = 20$. Hence $X \sim B(20, 0.3)$.
\item Compute $(n+1)p = 21 \times 0.3 = 6.3$, not an integer. By the recurrence relation, the probabilities increase while $k < (n+1)p - 1 = 5.3$ and decrease from $k = 6$ onwards. Hence the mode is $m = \lfloor 6.3 \rfloor = 6$.
\item $P(X = 6) = \dbinom{20}{6}(0.3)^6(0.7)^{14} = 38760 \times 0.000729 \times 0.0067822 \approx \mathbf{0.1916}$.
\end{enumerate}
\end{corrige}

\begin{corrige}[Exercise 6 — Uniform distribution from first principles]
Each value $k \in \{1, \dots, 20\}$ has probability $\frac{1}{20}$.
\begin{enumerate}
\item Using $\displaystyle\sum_{k=1}^{n} k = \frac{n(n+1)}{2}$ with $n = 20$:
\[ E(X) = \sum_{k=1}^{20} k \cdot \frac{1}{20} = \frac{1}{20}\cdot\frac{20 \times 21}{2} = \frac{210}{20} = 10.5. \qquad \blacksquare \]
\item Using $\displaystyle\sum_{k=1}^{n} k^2 = \frac{n(n+1)(2n+1)}{6}$ with $n = 20$:
\[ E(X^2) = \sum_{k=1}^{20} k^2 \cdot \frac{1}{20} = \frac{1}{20}\times\frac{20 \times 21 \times 41}{6} = \frac{21 \times 41}{6} = \frac{861}{6} = 143.5. \]
Then $\mathrm{Var}(X) = E(X^2) - [E(X)]^2 = 143.5 - 10.5^2 = 143.5 - 110.25 = 33.25$. \qquad $\blacksquare$\par
\emph{Check with the standard formula:} $\mathrm{Var}(X) = \dfrac{n^2 - 1}{12} = \dfrac{400 - 1}{12} = \dfrac{399}{12} = 33.25$. \checkmark
\item $P(X > 12) = P(X \in \{13, \dots, 20\}) = \dfrac{8}{20} = 0.4$.\par
The even values are $\{2, 4, \dots, 20\}$, i.e. $10$ values: $P(X \text{ even}) = \dfrac{10}{20} = 0.5$.
\end{enumerate}
\end{corrige}

\begin{corrige}[Exercise 7 — Geometric distribution: the driving test]
$X \sim$ geometric with $p = 0.4$, $q = 0.6$: $P(X = k) = q^{k-1}p$.
\begin{enumerate}
\item $P(X = 3) = (0.6)^2 \times 0.4 = 0.36 \times 0.4 = \mathbf{0.144}$.
\item "More than $4$ attempts" means the first $4$ attempts all fail:
\[ P(X > 4) = q^4 = (0.6)^4 = \mathbf{0.1296}. \]
\item $E(X) = \dfrac{1}{p} = \dfrac{1}{0.4} = 2.5$ attempts; \quad $\mathrm{Var}(X) = \dfrac{q}{p^2} = \dfrac{0.6}{0.16} = 3.75$.
\item $P(X \leq n) = 1 - q^n = 1 - (0.6)^n > 0.95 \iff (0.6)^n < 0.05$.
Taking logarithms: $n \ln 0.6 < \ln 0.05$, and since $\ln 0.6 < 0$ the inequality reverses on division:
\[ n > \frac{\ln 0.05}{\ln 0.6} \approx \frac{-2.9957}{-0.5108} \approx 5.86. \]
The smallest integer is $\mathbf{n = 6}$. Check: $(0.6)^6 = 0.0467 < 0.05$, so $P(X \le 6) = 0.9533 > 0.95$. \checkmark
\end{enumerate}
\end{corrige}

\begin{corrige}[Exercise 8 — Poisson: calls at an office and sum of independent variables]
Calls: rate $5$ per hour, so in one hour $X \sim \mathrm{Po}(5)$.
\begin{enumerate}
\item $P(X = 3) = e^{-5}\dfrac{5^3}{3!} = e^{-5}\times\dfrac{125}{6} \approx \mathbf{0.1404}$.
\item In $15$ minutes the mean is $\lambda = 5 \times \frac{15}{60} = 1.25$ (the parameter is proportional to the length of the interval — the \emph{unit interval} must be adjusted). With $N \sim \mathrm{Po}(1.25)$:
\[ P(N = 0) = e^{-1.25} \approx \mathbf{0.2865}. \]
\item In $10$ minutes, $\lambda = 5 \times \frac{10}{60} = \frac{5}{6} \approx 0.8333$. With $M \sim \mathrm{Po}(0.8333)$:
\[ P(M \geq 1) = 1 - P(M = 0) = 1 - e^{-0.8333} \approx 1 - 0.4346 = \mathbf{0.5654}. \]
\item Emails: $Y \sim \mathrm{Po}(3)$ per hour, independent of $X \sim \mathrm{Po}(5)$. By the additive property of independent Poisson variables, $X + Y \sim \mathrm{Po}(5 + 3) = \mathrm{Po}(8)$.
\[ P(X + Y = 4) = e^{-8}\frac{8^4}{4!} = e^{-8}\times\frac{4096}{24} \approx \mathbf{0.0573}. \]
\[ P(X + Y \geq 2) = 1 - P(0) - P(1) = 1 - e^{-8}(1 + 8) = 1 - 9e^{-8} \approx 1 - 0.0030 = \mathbf{0.9970}. \]
\end{enumerate}
\end{corrige}

\courssub{Solutions to "I prepare for the GCE"}

\begin{corrige}[Exercise 9 — GCE-style: binomial tables and a critical value]
\begin{enumerate}
\item $X \sim B(8, 0.3)$ \textbf{(B1)}; $E(X) = np = 8 \times 0.3 = 2.4$ \textbf{(M1)}; $\mathrm{Var}(X) = npq = 8 \times 0.3 \times 0.7 = 1.68$ \textbf{(A1)}.
\item From cumulative tables for $B(8, 0.3)$: $P(X \leq 1) = 0.2553$.
\[ P(X \geq 2) = 1 - P(X \leq 1) = 1 - 0.2553 = \mathbf{0.7447} \quad \textbf{(M1 A1)}. \]
\item $P(1 \leq X \leq 4) = P(X \leq 4) - P(X \leq 0) = 0.9420 - 0.0576 = \mathbf{0.8844}$ \textbf{(M1 for the difference, A1 for each value used)}.
\item $P(X = 5) = P(X \leq 5) - P(X \leq 4) = 0.9887 - 0.9420 = \mathbf{0.0467}$ \textbf{(M1 A1)}.\par
\emph{Check by formula:} $\dbinom{8}{5}(0.3)^5(0.7)^3 = 56 \times 0.00243 \times 0.343 = 0.0467$. \checkmark
\item From the tables: $P(X \leq 3) = 0.8059 \leq 0.9$ and $P(X \leq 4) = 0.9420 > 0.9$. Hence the smallest integer is $\mathbf{r = 4}$ \textbf{(M1 A1)}.
\item $(n+1)p = 9 \times 0.3 = 2.7$, not an integer, so the mode is $m = \lfloor 2.7 \rfloor = 2$ \textbf{(M1 A1)}. Equivalently, the ratio $\dfrac{P(X=k+1)}{P(X=k)} = \dfrac{8-k}{k+1}\cdot\dfrac{3}{7}$ is $> 1$ for $k = 0, 1$ and $< 1$ for $k \geq 2$.
\[ P(X = 2) = \binom{8}{2}(0.3)^2(0.7)^6 = 28 \times 0.09 \times 0.117649 \approx \mathbf{0.2965} \quad \textbf{(M1 A1)}. \]
\end{enumerate}
\textbf{Examiner's expectations:} correct reading of cumulative tables (subtracting the right cumulative value), clear statement of the distribution with both parameters, and answers given to 4 decimal places. A common error is writing $P(X \geq 2) = 1 - P(X \leq 2)$: this loses the method mark.
\end{corrige}

\begin{corrige}[Exercise 10 — GCE-style: Poisson approximation to the binomial]
\begin{enumerate}
\item $X \sim B(400, 0.015)$ \textbf{(B1)}; $E(X) = np = 400 \times 0.015 = 6$ \textbf{(A1)}; $\mathrm{Var}(X) = npq = 6 \times 0.985 = 5.91$ \textbf{(A1)}.
\item $n = 400$ is large and $p = 0.015$ is small (with $np = 6$ moderate), so the Poisson approximation is valid \textbf{(M1)}: $X \approx \mathrm{Po}(\lambda)$ with $\lambda = np = 6$ \textbf{(A1)}.
\item $P(X = 6) \approx e^{-6}\dfrac{6^6}{6!} = e^{-6}\times\dfrac{46656}{720} \approx \mathbf{0.1606}$ \textbf{(M1 A1)}.
\item From cumulative Poisson tables with $\lambda = 6$: $P(X \leq 8) = 0.8472$.
\[ P(X > 8) = 1 - P(X \leq 8) = 1 - 0.8472 = \mathbf{0.1528} \quad \textbf{(M1 A1 A1)}. \]
\item $P(3 \leq X \leq 9) = P(X \leq 9) - P(X \leq 2) = 0.9161 - 0.0620 = \mathbf{0.8541}$ \textbf{(M1 A1 A1)}.
\item $P(X > 10) = 1 - P(X \leq 10) = 1 - 0.9574 = 0.0426$ \textbf{(M1)}. Observing more than $10$ defectives has probability about $4.3\%$ under the model: this is a rare event (less than $5\%$), so it raises \emph{some} doubt about the claimed rate of $1.5\%$ — the true rate may be higher — although such an event can still occur by chance; further days should be monitored before concluding \textbf{(A1)}.
\end{enumerate}
\textbf{Examiner's expectations:} the conditions for the approximation must be \emph{stated} ($n$ large, $p$ small, $np$ moderate); in (f) a numerical probability must support the comment — a comment without computation earns no mark.
\end{corrige}

\begin{corrige}[Exercise 11 — GCE-style: from a frequency distribution to a Poisson model]
\begin{enumerate}
\item Total goals: $0(8) + 1(15) + 2(13) + 3(9) + 4(4) + 5(1) = 15 + 26 + 27 + 16 + 5 = 89$ \textbf{(M1)}.
\[ \bar{x} = \frac{89}{50} = 1.78 \text{ goals per match} \quad \textbf{(A1 A1)}. \]
\item Goals are rare, random events occurring independently at an approximately constant average rate per match, so a Poisson model is appropriate \textbf{(B1)}. We estimate $\lambda$ by the sample mean (since $E(X) = \lambda$ for a Poisson distribution): $\hat{\lambda} = 1.78 \approx 1.8$ \textbf{(B1)}.
\item With $\lambda = 1.8$:
\begin{align*}
P(X = 0) &= e^{-1.8} \approx 0.1653 &\Rightarrow\quad &50 \times 0.1653 \approx 8.3,\\
P(X = 1) &= 1.8\,e^{-1.8} \approx 0.2975 &\Rightarrow\quad &50 \times 0.2975 \approx 14.9,\\
P(X = 2) &= e^{-1.8}\frac{1.8^2}{2} \approx 0.2678 &\Rightarrow\quad &50 \times 0.2678 \approx 13.4.
\end{align*}
\textbf{(M1 for the pmf, A1 for each probability, A1 for the expected frequencies.)}
\begin{center}
\begin{tabular}{|l|ccc|}
\hline
\rowcolor{popblueL} Goals $x$ & $0$ & $1$ & $2$ \\
\hline Observed frequency & $8$ & $15$ & $13$ \\
\hline Expected frequency & $8.3$ & $14.9$ & $13.4$ \\
\hline
\end{tabular}
\end{center}
\item The expected frequencies are very close to the observed ones (differences of at most $0.4$) \textbf{(M1 A1)}: the Poisson model with $\lambda = 1.8$ fits the data very well \textbf{(A1)}.
\item $P(X \geq 3) = 1 - [P(0) + P(1) + P(2)] = 1 - (0.1653 + 0.2975 + 0.2678) = 1 - 0.7306 = \mathbf{0.2694}$ \textbf{(M1 A1)}.
\end{enumerate}
\textbf{Examiner's expectations:} the mean must be computed with full working from the frequency table; the estimate $\hat{\lambda} = \bar{x}$ must be \emph{justified} (the mean of a Poisson distribution is $\lambda$); expected frequencies are obtained by multiplying each probability by the total frequency $50$; the comparison must be quantitative, not just "the fit is good".
\end{corrige}

\newpage

\coursec{Summary sheet}

\begin{popfigure}
\begin{center}\tcbox[colback=popink,colframe=popink,coltext=white,arc=9pt,boxsep=4pt,left=6pt,right=6pt]{\bfseries\small Special Discrete Distributions}\end{center}
\vspace{-2pt}
\begin{tcbraster}[raster columns=3,raster equal height=rows,raster column skip=6pt,raster row skip=6pt,enhanced,blankest]
\begin{tcolorbox}[enhanced,colback=popblue,colframe=popblue,coltext=white,boxrule=0.9pt,arc=3pt,left=4pt,right=4pt,top=3pt,bottom=3pt,boxsep=1pt,halign=left]\bfseries\scriptsize Binomial $B(n,p)$ $P(X{=}k){=}\binom{n}{k}p^kq^{n-k}$ $\mu{=}np,\ \sigma^2{=}npq$\end{tcolorbox}
\begin{tcolorbox}[enhanced,colback=poporange,colframe=poporange,coltext=white,boxrule=0.9pt,arc=3pt,left=4pt,right=4pt,top=3pt,bottom=3pt,boxsep=1pt,halign=left]\bfseries\scriptsize Binomial tables $P(X\le r)$ Mode via recurrence\end{tcolorbox}
\begin{tcolorbox}[enhanced,colback=popgreen,colframe=popgreen,coltext=white,boxrule=0.9pt,arc=3pt,left=4pt,right=4pt,top=3pt,bottom=3pt,boxsep=1pt,halign=left]\bfseries\scriptsize Uniform on $\{1,\dots,n\}$ $\mu{=}\frac{n+1}{2}$ $\sigma^2{=}\frac{n^2-1}{12}$\end{tcolorbox}
\begin{tcolorbox}[enhanced,colback=poppurple,colframe=poppurple,coltext=white,boxrule=0.9pt,arc=3pt,left=4pt,right=4pt,top=3pt,bottom=3pt,boxsep=1pt,halign=left]\bfseries\scriptsize Geometric $G(p)$ $P(X{=}k){=}q^{k-1}p$ $\mu{=}\frac1p,\ \sigma^2{=}\frac{q}{p^2}$\end{tcolorbox}
\begin{tcolorbox}[enhanced,colback=popgold,colframe=popgold,coltext=white,boxrule=0.9pt,arc=3pt,left=4pt,right=4pt,top=3pt,bottom=3pt,boxsep=1pt,halign=left]\bfseries\scriptsize Poisson $P(\lambda)$ $P(X{=}k){=}e^{-\lambda}\frac{\lambda^k}{k!}$ $\mu{=}\sigma^2{=}\lambda$\end{tcolorbox}
\begin{tcolorbox}[enhanced,colback=poppink,colframe=poppink,coltext=white,boxrule=0.9pt,arc=3pt,left=4pt,right=4pt,top=3pt,bottom=3pt,boxsep=1pt,halign=left]\bfseries\scriptsize Poisson $\approx$ Binomial $n$ large, $p$ small $\lambda=np$ Sum of Poissons\end{tcolorbox}
\end{tcbraster}

\legende{Mind map of the chapter: the special discrete probability distributions.}
\end{popfigure}

\begin{bilan}
\textbf{Key definitions}
\begin{itemize}
\item \cle{Binomial distribution} $X\sim B(n,p)$: number of successes in $n$ independent identical trials, each with $P(\text{success})=p$, $q=1-p$.
\item \cle{Discrete uniform distribution} on $\{1,2,\dots,n\}$: every outcome equally likely, $P(X=k)=\dfrac{1}{n}$.
\item \cle{Geometric distribution} $X\sim G(p)$: number of trials up to and including the first success; memoryless.
\item \cle{Poisson distribution} $X\sim P(\lambda)$: number of rare events in a fixed interval of time/space, rate $\lambda$ per unit interval.
\item \cle{Mode}: the value of $k$ with the greatest probability.
\end{itemize}

\textbf{Formulas to know by heart}
\begin{itemize}
\item Binomial: $P(X=k)=\dbinom{n}{k}p^{k}q^{\,n-k}$,\quad $E(X)=np$,\quad $\mathrm{Var}(X)=npq$.
\item Recurrence: $\dfrac{P(X=k+1)}{P(X=k)}=\dfrac{n-k}{k+1}\cdot\dfrac{p}{q}$; the mode is the integer part of $(n+1)p$ (if $(n+1)p$ is an integer, both $(n+1)p$ and $(n+1)p-1$ are modes).
\item Uniform: $E(X)=\dfrac{n+1}{2}$,\quad $\mathrm{Var}(X)=\dfrac{n^{2}-1}{12}$.
\item Geometric: $P(X=k)=q^{\,k-1}p$,\quad $P(X>k)=q^{k}$,\quad $E(X)=\dfrac{1}{p}$,\quad $\mathrm{Var}(X)=\dfrac{q}{p^{2}}$.
\item Poisson: $P(X=k)=e^{-\lambda}\dfrac{\lambda^{k}}{k!}$,\quad $E(X)=\mathrm{Var}(X)=\lambda$.
\item Scaling: if $\lambda$ is the rate per unit interval, the rate over an interval of length $t$ is $\lambda t$.
\item Sum of independent Poissons: $X\sim P(\lambda_1)$, $Y\sim P(\lambda_2)$ independent $\Rightarrow X+Y\sim P(\lambda_1+\lambda_2)$.
\item Poisson approximation: $B(n,p)\approx P(np)$ when $n$ is large ($n\ge 50$) and $p$ small ($p\le 0.1$), so that $\lambda=np$ is moderate.
\end{itemize}

\textbf{Using the tables}
\begin{itemize}
\item Cumulative binomial tables give $P(X\le r)$ directly. Then:
$P(X\ge r)=1-P(X\le r-1)$;\quad $P(X=r)=P(X\le r)-P(X\le r-1)$;\quad $P(a\le X\le b)=P(X\le b)-P(X\le a-1)$.
\item Cumulative Poisson tables give $P(X\le r)$ for tabulated $\lambda$; use the same transformations.
\item If $p>0.5$ for the binomial, work with the number of failures $Y\sim B(n,1-p)$.
\end{itemize}

\textbf{Methods}
\begin{itemize}
\item \cle{Identify the model first}: fixed number of trials $\to$ binomial; waiting for first success $\to$ geometric; random rare events in an interval $\to$ Poisson; equally likely scores $\to$ uniform.
\item \cle{Mode of a binomial}: compute $\dfrac{P(X=k+1)}{P(X=k)}\ge 1 \iff k\le (n+1)p-1$; probabilities increase up to the mode then decrease.
\item \cle{Frequency link}: expected frequency $=$ (total frequency) $\times P(X=k)$; compare with observed frequencies to judge goodness of fit.
\end{itemize}

\textbf{Common mistakes}
\begin{itemize}
\item Confusing $P(X\le r)$ with $P(X<r)$: $P(X<r)=P(X\le r-1)$.
\item Forgetting $q=1-p$, or writing $q^{n-k}$ with the wrong exponent.
\item Using the Poisson approximation when $p$ is not small.
\item For the geometric law, forgetting that $X$ counts the successful trial too: $P(X=k)=q^{k-1}p$, not $q^{k}p$.
\item Forgetting to adjust $\lambda$ when the interval changes (e.g.\ per hour $\to$ per 15 minutes: divide by 4).
\item Adding Poisson variables that are not independent.
\end{itemize}

\textbf{GCE tips}
\begin{itemize}
\item Quote the distribution and parameters first: ``$X\sim B(10,0.3)$'' earns the first mark.
\item Show the formula before substituting; examiners reward method marks.
\item Give final probabilities to 4 decimal places (as in the tables).
\item For ``mean and variance'' questions, state $E(X)$ and $\mathrm{Var}(X)$, not the standard deviation, unless asked.
\item Always check: probabilities must lie in $[0,1]$ and $\sum P(X=k)=1$.
\end{itemize}
\end{bilan}

\findecours

\end{document}
